% Options for packages loaded elsewhere
\PassOptionsToPackage{unicode}{hyperref}
\PassOptionsToPackage{hyphens}{url}
\PassOptionsToPackage{dvipsnames,svgnames,x11names}{xcolor}
%
\documentclass[
  13pt]{extreport}
\usepackage{amsmath,amssymb}
\usepackage{setspace}
\usepackage{iftex}
\ifPDFTeX
  \usepackage[T1]{fontenc}
  \usepackage[utf8]{inputenc}
  \usepackage{textcomp} % provide euro and other symbols
\else % if luatex or xetex
  \usepackage{unicode-math} % this also loads fontspec
  \defaultfontfeatures{Scale=MatchLowercase}
  \defaultfontfeatures[\rmfamily]{Ligatures=TeX,Scale=1}
\fi
\usepackage{lmodern}
\ifPDFTeX\else
  % xetex/luatex font selection
  \setmainfont[]{Times New Roman}
\fi
% Use upquote if available, for straight quotes in verbatim environments
\IfFileExists{upquote.sty}{\usepackage{upquote}}{}
\IfFileExists{microtype.sty}{% use microtype if available
  \usepackage[]{microtype}
  \UseMicrotypeSet[protrusion]{basicmath} % disable protrusion for tt fonts
}{}
\usepackage{xcolor}
\usepackage[left=3cm,right=2cm,top=2cm,bottom=2cm]{geometry}
\usepackage{longtable,booktabs,array}
\usepackage{calc} % for calculating minipage widths
% Correct order of tables after \paragraph or \subparagraph
\usepackage{etoolbox}
\makeatletter
\patchcmd\longtable{\par}{\if@noskipsec\mbox{}\fi\par}{}{}
\makeatother
% Allow footnotes in longtable head/foot
\IfFileExists{footnotehyper.sty}{\usepackage{footnotehyper}}{\usepackage{footnote}}
\makesavenoteenv{longtable}
\setlength{\emergencystretch}{3em} % prevent overfull lines
\providecommand{\tightlist}{%
  \setlength{\itemsep}{0pt}\setlength{\parskip}{0pt}}
\setcounter{secnumdepth}{-\maxdimen} % remove section numbering
\ifLuaTeX
\usepackage[bidi=basic]{babel}
\else
\usepackage[bidi=default]{babel}
\fi
\babelprovide[main,import]{vietnamese}
\ifPDFTeX
\else
\babelfont{rm}[]{Times New Roman}
\fi
% get rid of language-specific shorthands (see #6817):
\let\LanguageShortHands\languageshorthands
\def\languageshorthands#1{}
\ifLuaTeX
  \usepackage{selnolig}  % disable illegal ligatures
\fi
\IfFileExists{bookmark.sty}{\usepackage{bookmark}}{\usepackage{hyperref}}
\IfFileExists{xurl.sty}{\usepackage{xurl}}{} % add URL line breaks if available
\urlstyle{same}
\hypersetup{
  pdflang={vi},
  colorlinks=true,
  linkcolor={black},
  filecolor={Maroon},
  citecolor={black},
  urlcolor={blue},
  pdfcreator={LaTeX via pandoc}}

\author{}
\date{}

\begin{document}

{
\hypersetup{linkcolor=}
\setcounter{tocdepth}{2}
\tableofcontents
}
\setstretch{1.2}
BỘ GIÁO DỤC VÀ ĐÀO TẠO TRƯỜNG ĐẠI HỌC CẦN THƠ TRƯỜNG KINH TẾ

─────────────────────────────────────────

ĐỖ THÙY HƯƠNG

CHUYÊN ĐỀ TIẾN SĨ SỐ 1

THỰC TRẠNG VỀ HIỆU QUẢ HOẠT ĐỘNG KINH DOANH CỦA CÁC DOANH NGHIỆP Ở CHÂU
Á

Ngành: Quản trị kinh doanh Mã ngành: 9340101 Mã nghiên cứu sinh:
P1323001

NGƯỜI HƯỚNG DẪN KHOA HỌC TS. NGUYỄN MINH CẢNH

Cần Thơ, năm 2026

\begin{center}\rule{0.5\linewidth}{0.5pt}\end{center}

\hypertarget{lux1eddi-cam-ux111oan}{%
\subsection{LỜI CAM ĐOAN}\label{lux1eddi-cam-ux111oan}}

Tôi xin cam đoan chuyên đề tiến sĩ số 1 với tiêu đề ``Thực trạng về hiệu
quả hoạt động kinh doanh của các doanh nghiệp ở Châu Á'' là công trình
nghiên cứu của riêng tôi dưới sự hướng dẫn khoa học của TS. Nguyễn Minh
Cảnh, Trường Đại học Cần Thơ. Các kết quả phân tích, số liệu, bảng biểu,
hình vẽ trình bày trong chuyên đề đều trung thực, có nguồn gốc rõ ràng.
Tài liệu tham khảo được trích dẫn theo chuẩn APA 7th. Những trích dẫn từ
các nghiên cứu khác đều được nêu rõ trong danh mục tài liệu tham khảo.
Tôi xin chịu trách nhiệm hoàn toàn về nội dung khoa học của chuyên đề
này.

\emph{Cần Thơ, ngày \ldots{} tháng \ldots{} năm 2026}

\emph{Nghiên cứu sinh}

\textbf{Đỗ Thùy Hương}

\begin{center}\rule{0.5\linewidth}{0.5pt}\end{center}

\hypertarget{tuxf3m-tux1eaft}{%
\subsection{TÓM TẮT}\label{tuxf3m-tux1eaft}}

Chuyên đề này thực hiện phân tích mô tả -- chẩn đoán thực trạng hiệu quả
hoạt động kinh doanh của doanh nghiệp ở \textbf{các nền kinh tế châu Á
(phạm vi chính)} trong giai đoạn \textbf{2006--2026}, mở rộng so sánh
với các quốc đảo nhỏ đang phát triển (SIDS) Thái Bình Dương như trường
hợp biên. Phạm vi không gian gồm \textbf{49 nền kinh tế}, phân thành 6
nhóm theo khung Biến thiên chế độ bối cảnh thể chế (ICRV): Advanced đổi
mới sáng tạo dẫn dắt, Advanced tài nguyên dẫn dắt, trung bình cao,
chuyển đổi thu nhập trung bình--thấp, đang nổi, và SIDS Thái Bình Dương
(trường hợp biên mở rộng). Phần phân tích mô tả (Bảng 2.3.3--2.3.8) được
tính trực tiếp từ dữ liệu vi mô World Bank Enterprise Surveys (WBES)
trên \textbf{pool mô tả gồm 86.701 doanh nghiệp · 102 cặp nền kinh tế ×
năm}, kế thừa và mở rộng từ pool 17 nền châu Á mới nổi
(\textasciitilde40.633 doanh nghiệp) đã được tác giả công bố trước đó
(Đỗ \& Phan, 2026 --- VEFR). Khung phân loại thống nhất với bộ dữ liệu
ước lượng của luận án (khung phân tích 91.982 doanh nghiệp / 49 nền;
pool phân loại 96.415 doanh nghiệp / 52 nền).

\begin{quote}
\emph{Ghi chú dữ liệu: Chuyên đề sử dụng hai khung dữ liệu bổ trợ, thống
nhất về danh sách 49 nền kinh tế và nhãn nhóm ICRV. (i) \textbf{Pool mô
tả} phục vụ toàn bộ thống kê §2.3.3--2.3.8: 86.701 doanh nghiệp · 102
cặp nền kinh tế × năm (2006--2026), tính trực tiếp từ tệp vi mô WBES ---
mỗi cặp nền kinh tế × năm lấy đúng một cross-section chuẩn (loại các
điều tra phi chính thức, siêu nhỏ và panel mở rộng), chỉ nhận các đợt
khảo sát từ 2006 trở đi khi WBES áp dụng bộ câu hỏi so sánh được toàn
cầu. Pool này chứa đầy đủ các biến mô tả (R\&D, ISO, đổi mới, tham
nhũng, giới) vốn không thuộc tập biến ước lượng. (ii) \textbf{Khung ước
lượng} của luận án và Chuyên đề 2: 91.982 doanh nghiệp / 49 nền (pool
phân loại 96.415 / 52 nền), được hợp nhất và hài hòa hóa theo quy trình
tại Phụ lục A của luận án --- mọi kết quả kinh tế lượng được dẫn chiếu
theo khung này. Trường hợp biên SIDS: Nhóm VI gồm 8 nền (7 Pacific +
Timor-Leste); mẫu chính của nghiên cứu SIDS là 7 nền Pacific.}
\end{quote}

Phương pháp chủ đạo là thống kê mô tả đa chiều, kết hợp đồ thị, bảng so
sánh xuyên quốc gia -- xuyên ngành và phân tích hai biến. Hiệu quả hoạt
động kinh doanh được tiếp cận theo bốn chiều: năng suất lao động, lợi
nhuận và biên lợi nhuận, tăng trưởng, đổi mới sáng tạo và cấu trúc doanh
nghiệp.

Kết quả cho thấy hiệu quả doanh nghiệp châu Á có \textbf{phân tán lớn và
không hội tụ}. Bảy tiểu cảnh điển hình tồn tại đồng thời (Singapore --
Vùng Vịnh Saudi/Qatar/Kuwait -- Việt Nam -- Trung Quốc -- Emerging Asia
-- Mongolia -- SIDS Thái Bình Dương) với điểm mạnh hiệu quả khác nhau.
Chuyên đề phát hiện \textbf{phân nhóm con Advanced} (innovation-driven
so với resource-driven) và xác nhận pattern \textbf{chi phí buộc phải
quốc tế hóa} ở 7 SIDS Thái Bình Dương --- đóng góp định hướng cho luận
án và Chuyên đề tiến sĩ số 2.

\textbf{Từ khóa}: hiệu quả doanh nghiệp; châu Á; WBES; thực trạng; năng
suất lao động; doanh nghiệp vừa và nhỏ; phân nhóm con Advanced; SIDS
Thái Bình Dương.

\begin{center}\rule{0.5\linewidth}{0.5pt}\end{center}

\hypertarget{abstract}{%
\subsection{ABSTRACT}\label{abstract}}

This specialty essay conducts a descriptive-diagnostic analysis of the
firm performance landscape across \textbf{Asian economies (primary
scope)} during \textbf{2006--2026}, with comparative extension to
Pacific SIDS as a boundary case. The spatial scope covers \textbf{49
economies}, classified into six sub-regimes of ICRV (Institutional
Context Regime Variation): Advanced innovation-driven, Advanced
resource-driven, Upper-middle, Lower-middle transition, Emerging, and
Pacific SIDS (boundary case extension). The descriptive analysis (Tables
2.3.3--2.3.8) is computed directly from World Bank Enterprise Surveys
(WBES) micro-data over a descriptive pool of \textbf{86,701 firms · 102
economy-year pairs} (one standard cross-section per economy-year, waves
from 2006 onward), building on and extending the 17-economy pool of
emerging Asia (\textasciitilde40,633 firms) previously published by the
author (Do \& Phan, 2026 --- VEFR). The classification is identical to
the dissertation's estimation dataset (analytic frame 91,982 firms / 49
economies; classification pool 96,415 firms / 52 economies).

Adopting multidimensional descriptive statistics, charts,
cross-country--cross-industry comparison tables and bivariate analyses,
the study examines firm performance along four axes: labor productivity,
profitability and margins, growth, and innovation and structural
composition. Findings indicate substantial and persistent dispersion in
firm performance across Asian economies with no convergence pattern; six
salient archetypes (Singapore, Saudi/Qatar/Kuwait, Vietnam, China,
broader Emerging Asia, Mongolia --- primary scope) plus Pacific SIDS as
boundary case coexist with distinct performance strengths; bivariate
correlations vary in sign and magnitude across regimes, suggesting
nonlinear and moderation models. The essay identifies a sub-grouping of
the Advanced regime (innovation- vs resource-driven) within primary
Asian scope and confirms the forced internationalization penalty pattern
in 7 Pacific SIDS boundary case.

\textbf{Keywords}: firm performance; Asia; WBES; landscape; labor
productivity; SMEs; Advanced sub-grouping; Pacific SIDS boundary case.

\begin{center}\rule{0.5\linewidth}{0.5pt}\end{center}

\hypertarget{mux1ee5c-lux1ee5c}{%
\subsection{MỤC LỤC}\label{mux1ee5c-lux1ee5c}}

\begin{itemize}
\tightlist
\item
  Lời cam đoan
\item
  Tóm tắt / Abstract
\item
  Danh mục bảng
\item
  Danh mục hình
\item
  Danh mục từ viết tắt
\item
  \textbf{2.1 Giới thiệu}

  \begin{itemize}
  \tightlist
  \item
    2.1.1 Đặt vấn đề và tính cấp thiết
  \item
    2.1.2 Mục tiêu chuyên đề
  \item
    2.1.3 Nội dung và phạm vi nghiên cứu
  \item
    2.1.4 Giới hạn nghiên cứu
  \item
    2.1.5 Ý nghĩa khoa học và thực tiễn
  \end{itemize}
\item
  \textbf{2.2 Phương pháp nghiên cứu}
\item
  \textbf{2.3 Nội dung nghiên cứu}

  \begin{itemize}
  \tightlist
  \item
    2.3.1 Cơ sở lý luận về hiệu quả hoạt động kinh doanh
  \item
    2.3.2 Khung phân loại thể chế ICRV
  \item
    2.3.3 Thực trạng năng suất lao động và quốc tế hóa
  \item
    2.3.4 Đổi mới sáng tạo và năng lực số hóa
  \item
    2.3.5 Cấu trúc doanh nghiệp
  \item
    2.3.6 Biến thiên theo thời gian (2006--2026)
  \item
    2.3.7 Bảy tiểu cảnh điển hình
  \item
    2.3.8 Yếu tố giải thích sơ bộ
  \item
    2.3.9 Thảo luận tổng hợp
  \end{itemize}
\item
  \textbf{2.4 Kết luận}

  \begin{itemize}
  \tightlist
  \item
    2.4.1 Kết luận chính
  \item
    2.4.2 Hàm ý và đề xuất
  \end{itemize}
\item
  Tài liệu tham khảo
\item
  Phụ lục A --- Phạm vi nhóm dữ liệu WBES
\end{itemize}

\begin{center}\rule{0.5\linewidth}{0.5pt}\end{center}

\hypertarget{danh-mux1ee5c-bux1ea3ng}{%
\subsection{DANH MỤC BẢNG}\label{danh-mux1ee5c-bux1ea3ng}}

\begin{longtable}[]{@{}
  >{\raggedright\arraybackslash}p{(\columnwidth - 4\tabcolsep) * \real{0.3333}}
  >{\raggedright\arraybackslash}p{(\columnwidth - 4\tabcolsep) * \real{0.3333}}
  >{\raggedright\arraybackslash}p{(\columnwidth - 4\tabcolsep) * \real{0.3333}}@{}}
\toprule\noalign{}
\begin{minipage}[b]{\linewidth}\raggedright
Bảng
\end{minipage} & \begin{minipage}[b]{\linewidth}\raggedright
Tiêu đề
\end{minipage} & \begin{minipage}[b]{\linewidth}\raggedright
Trang
\end{minipage} \\
\midrule\noalign{}
\endhead
\bottomrule\noalign{}
\endlastfoot
Bảng 2.3.1.1 & Ma trận đối chiếu 3 trạng thái U-curve theo 3 điều kiện
cấu trúc & 2.3 \\
Bảng 2.3.2.1 & 49 nền kinh tế phân theo 6 nhóm ICRV & 2.3 \\
Bảng 2.3.2.2 & Đặc điểm các nhóm ICRV (tiêu chí + đại diện + pattern dự
kiến) & 2.3 \\
Bảng 2.3.3.1 & Phân tán năng suất lao động trong từng đợt khảo sát theo
nhóm ICRV & 2.3 \\
Bảng 2.3.3.2 & FSTS, doanh nghiệp xuất khẩu và CAGR việc làm theo phân
nhóm con & 2.3 \\
Bảng 2.3.4.1 & Đổi mới sáng tạo và áp dụng số (\%) & 2.3 \\
Bảng 2.3.4.2 & Khung 5 lĩnh vực chỉ số WBES & 2.3 \\
Bảng 2.3.5.1 & Cấu trúc doanh nghiệp theo phân nhóm con (\%) & 2.3 \\
Bảng 2.3.5.2 & Bốn rào cản hàng đầu Asia-Pacific WBES & 2.3 \\
Bảng 2.3.6.1 & Thay đổi theo thời gian --- Δ điểm phần trăm 2018--2026
so với 2006--2012 & 2.3 \\
Bảng 2.3.6.2 & Phân biệt 3 chỉ số tham nhũng WBES --- Bribery vs Graft &
2.3 \\
Bảng 2.3.6.3 & Khung phân loại 9 ngành ISIC Rev.~4 & 2.3 \\
Bảng 2.3.6.4 & Dị biệt nội bộ khối châu Á mới nổi (7 nền Nhóm IV--V) &
2.3 \\
Bảng 2.3.6.5 & 12 nền kinh tế trong đợt khảo sát 2025 & 2.3 \\
Bảng 2.3.7.1 & Ma trận so sánh đa chiều bảy tiểu cảnh điển hình & 2.3 \\
Bảng 2.3.8.1 & Hệ số tương quan Pearson theo 6 nhóm ICRV (pool mô tả 49
nền) & 2.3 \\
Bảng 2.3.8.2 & Pattern giới tính trong lãnh đạo cấp cao và sở hữu ---
phân tầng theo khu vực & 2.3 \\
\end{longtable}

\begin{center}\rule{0.5\linewidth}{0.5pt}\end{center}

\hypertarget{danh-mux1ee5c-huxecnh}{%
\subsection{DANH MỤC HÌNH}\label{danh-mux1ee5c-huxecnh}}

\begin{longtable}[]{@{}
  >{\raggedright\arraybackslash}p{(\columnwidth - 4\tabcolsep) * \real{0.3333}}
  >{\raggedright\arraybackslash}p{(\columnwidth - 4\tabcolsep) * \real{0.3333}}
  >{\raggedright\arraybackslash}p{(\columnwidth - 4\tabcolsep) * \real{0.3333}}@{}}
\toprule\noalign{}
\begin{minipage}[b]{\linewidth}\raggedright
Hình
\end{minipage} & \begin{minipage}[b]{\linewidth}\raggedright
Tiêu đề
\end{minipage} & \begin{minipage}[b]{\linewidth}\raggedright
Trang
\end{minipage} \\
\midrule\noalign{}
\endhead
\bottomrule\noalign{}
\endlastfoot
Hình 2.3.3.1 & Phân bố pool mô tả 86.701 doanh nghiệp (2006--2026) theo
3 thế hệ schema WBES & 2.3 \\
Hình 2.3.3.2 & Phân bố pool mô tả theo 6 nhóm ICRV (n=86.701 doanh
nghiệp) & 2.3 \\
Hình 2.3.3.3 & Phân phối năng suất chuẩn hóa trong đợt theo 6 nhóm ICRV
(kernel density) & 2.3 \\
Hình 2.3.4.1 & Slope chart 2006--2012 vs 2018--2026 (Δ điểm phần trăm
theo nhóm ICRV) & 2.3 \\
Hình 2.3.5.1 & Phân phối quy mô doanh nghiệp theo nhóm ICRV (n=86.701) &
2.3 \\
Hình 2.3.6.1 & Spider chart 5 chiều × 2 mốc thời gian (Advanced,
Emerging, SIDS) & 2.3 \\
Hình 2.3.7.1 & Mông Cổ: tiến triển 2009--2019 (3 đợt khảo sát WBES) &
2.3 \\
\end{longtable}

\begin{center}\rule{0.5\linewidth}{0.5pt}\end{center}

\hypertarget{danh-mux1ee5c-tux1eeb-viux1ebft-tux1eaft}{%
\subsection{DANH MỤC TỪ VIẾT
TẮT}\label{danh-mux1ee5c-tux1eeb-viux1ebft-tux1eaft}}

\begin{longtable}[]{@{}
  >{\raggedright\arraybackslash}p{(\columnwidth - 4\tabcolsep) * \real{0.3333}}
  >{\raggedright\arraybackslash}p{(\columnwidth - 4\tabcolsep) * \real{0.3333}}
  >{\raggedright\arraybackslash}p{(\columnwidth - 4\tabcolsep) * \real{0.3333}}@{}}
\toprule\noalign{}
\begin{minipage}[b]{\linewidth}\raggedright
Từ viết tắt
\end{minipage} & \begin{minipage}[b]{\linewidth}\raggedright
Tiếng Anh
\end{minipage} & \begin{minipage}[b]{\linewidth}\raggedright
Tiếng Việt
\end{minipage} \\
\midrule\noalign{}
\endhead
\bottomrule\noalign{}
\endlastfoot
ADB & Asian Development Bank & Ngân hàng Phát triển Châu Á \\
AIPI & AI Productivity Impact & Tác động năng suất của AI \\
ASEAN & Association of Southeast Asian Nations & Hiệp hội các quốc gia
Đông Nam Á \\
BEE & Business Environment and Enterprise & --- \\
BREADY & WBES 2018+ schema & Thế hệ bảng hỏi WBES mới từ 2018 \\
CAGR & Compound Annual Growth Rate & Tốc độ tăng trưởng hàng năm kép \\
CDCM & Context-Contingent Digital-and-Capability Model & Mô hình Áp dụng
số và Năng lực phụ thuộc bối cảnh \\
CĐ1 & Chuyên đề tiến sĩ số 1 & Specialty Essay 1 \\
CĐ2 & Chuyên đề tiến sĩ số 2 & Specialty Essay 2 \\
CIEM & Central Institute for Economic Management & Viện Nghiên cứu Quản
lý Kinh tế Trung ương \\
CRM & Customer Relationship Management & Quản lý quan hệ khách hàng \\
DAI & Digital Adoption Index & Chỉ số áp dụng số \\
DAI (Tier-1) & Website binary & Sự hiện diện số cơ bản \\
DAI (Tier-2) & Website + e-payment composite & Chuyển đổi số cơ bản \\
DPL & Digital Paradox Lifecycle & Vòng đời nghịch lý số \\
EAP & East Asia and Pacific & Đông Á và Thái Bình Dương \\
ECA & Europe and Central Asia & Châu Âu và Trung Á \\
ERP & Enterprise Resource Planning & Hoạch định nguồn lực doanh
nghiệp \\
FDA & Foundational Digital Adoption & Áp dụng số nền tảng \\
FDI & Foreign Direct Investment & Đầu tư trực tiếp nước ngoài \\
FSTS & Foreign Trade Sales to Total Sales & Tỷ lệ doanh thu xuất khẩu
trên tổng doanh thu \\
GCC & Gulf Cooperation Council & Hội đồng Hợp tác vùng Vịnh \\
GDP & Gross Domestic Product & Tổng sản phẩm quốc nội \\
GNI & Gross National Income & Tổng thu nhập quốc gia \\
GVC & Global Value Chain & Chuỗi giá trị toàn cầu \\
HC1 & HC1 Robust Standard Errors & Sai số chuẩn vững HC1 \\
ICRV & Institutional Context Regime Variation & Biến thiên chế độ bối
cảnh thể chế \\
IFC & International Finance Corporation & Tổ chức Tài chính Quốc tế \\
IMF & International Monetary Fund & Quỹ Tiền tệ Quốc tế \\
IPR & Intellectual Property Rights & Quyền sở hữu trí tuệ \\
ISIC & International Standard Industrial Classification & Phân loại
ngành kinh tế quốc tế \\
ISO & International Organization for Standardization & Tổ chức Tiêu
chuẩn Quốc tế \\
IV & Instrumental Variable & Biến công cụ \\
LP & Labor Productivity & Năng suất lao động \\
MENA & Middle East and North Africa & Trung Đông và Bắc Phi \\
MIRAB & Migration, Remittances, Aid, Bureaucracy & Mô hình kinh tế SIDS
Pacific \\
MNE & Multinational Enterprise & Doanh nghiệp đa quốc gia \\
ODA & Official Development Assistance & Hỗ trợ phát triển chính thức \\
OLI & Ownership-Location-Internalization & Mô hình chiết trung OLI \\
PICS3 & WBES 2009--2012 schema & Thế hệ bảng hỏi WBES 2009--2012 \\
PICs & Pacific Island Countries & Các quốc đảo Thái Bình Dương \\
PLMS & Pacific Labour Mobility Scheme & Chương trình lao động di chuyển
Thái Bình Dương \\
PPP & Purchasing Power Parity & Ngang giá sức mua \\
R\&D & Research and Development & Nghiên cứu và Phát triển \\
RBV & Resource-Based View & Quan điểm dựa trên nguồn lực \\
SIDS & Small Island Developing States & Quốc đảo nhỏ đang phát triển \\
SME & Small and Medium Enterprise & Doanh nghiệp vừa và nhỏ \\
TCI & Technological Capability Index & Chỉ số năng lực công nghệ \\
TMT & Top Management Team & Nhóm quản trị cấp cao \\
UNCTAD & United Nations Conference on Trade and Development & --- \\
WBES & World Bank Enterprise Surveys & Điều tra doanh nghiệp của Ngân
hàng Thế giới \\
WGI & World Governance Indicators & Chỉ số quản trị quốc gia \\
\end{longtable}

\begin{center}\rule{0.5\linewidth}{0.5pt}\end{center}

\hypertarget{giux1edbi-thiux1ec7u}{%
\subsection{2.1 GIỚI THIỆU}\label{giux1edbi-thiux1ec7u}}

\hypertarget{ux111ux1eb7t-vux1ea5n-ux111ux1ec1-vuxe0-tuxednh-cux1ea5p-thiux1ebft}{%
\subsubsection{2.1.1 Đặt vấn đề và tính cấp
thiết}\label{ux111ux1eb7t-vux1ea5n-ux111ux1ec1-vuxe0-tuxednh-cux1ea5p-thiux1ebft}}

Khu vực châu Á đã trở thành động lực tăng trưởng của kinh tế thế giới
trong gần hai thập niên qua. Theo Asian Development Bank (ADB, 2024),
châu Á đóng góp khoảng 60\% tăng trưởng GDP toàn cầu giai đoạn
2010--2023, với tổng quy mô kinh tế chiếm xấp xỉ 40\% GDP toàn cầu
(World Bank, 2024). Bên trong khu vực, hiệu quả hoạt động của doanh
nghiệp --- đơn vị tế bào của nền kinh tế --- là yếu tố quyết định mức độ
chuyển hoá tăng trưởng quốc gia thành thịnh vượng và năng lực cạnh tranh
dài hạn (Cusolito \& Maloney, 2018). Vì vậy, đánh giá thực trạng hiệu
quả hoạt động kinh doanh của doanh nghiệp ở châu Á có ý nghĩa lý luận
lẫn chính sách.

Bức tranh hiệu quả doanh nghiệp châu Á 2006--2026 được định hình bởi
\textbf{bảy lớp bối cảnh đan xen}. Bốn lớp đầu gồm: hậu khủng hoảng tài
chính toàn cầu 2008--2009; tái cấu trúc chuỗi giá trị toàn cầu sau chiến
tranh thương mại Mỹ--Trung 2018 (UNCTAD, 2023); đại dịch COVID-19
2020--2022 (World Bank, 2023); và làn sóng chuyển đổi số tăng tốc từ
2018 (Verhoef et al., 2021; Banalieva \& Dhanaraj, 2019). Lớp thứ năm là
\textbf{giai đoạn AI bùng nổ và củng cố hậu COVID 2023--2025} (Stallkamp
\& Schotter, 2021). \textbf{Lớp thứ sáu là xung đột Trung Đông 2026} ---
IMF (2026, tháng 4) điều chỉnh giảm tăng trưởng toàn cầu xuống 3,1\% năm
2026; ADB (2026, tháng 4) dự báo Pacific 3,4\%, Đông Nam Á 4,7\%, Việt
Nam 7,0\%. Riêng Việt Nam, ADB (2026, tháng 4) ghi nhận GDP lịch sử 2025
đạt 8,0\% và cập nhật dự báo 7,2\% (2026) → 7,0\% (2027). \textbf{Lớp
thứ bảy là bất định chính sách thuế quan Mỹ 2026} --- Mỹ chiếm xấp xỉ
30\% tổng kim ngạch xuất khẩu Việt Nam; thuế toàn cầu 10\% tạm thời tạo
áp lực kép: hiệu ứng front-loading ngắn hạn và trì hoãn đầu tư dài hạn.
\textbf{Lưu ý thời điểm}: Lớp bối cảnh thứ sáu và thứ bảy (các sự kiện
2026) phần lớn nằm \emph{ngoài} cửa sổ dữ liệu của chuyên đề (các đợt
khảo sát WBES từ 2006 đến đầu 2026); chúng được nêu để định khung diễn
giải vĩ mô và định hướng chính sách, chứ không phản ánh trong các thống
kê mô tả.

Trong bối cảnh đó, hệ thống thông tin về hiệu quả doanh nghiệp châu Á
còn phân mảnh. World Bank Enterprise Surveys (WBES) là cơ sở dữ liệu vi
mô tin cậy nhất hiện có (World Bank Enterprise Surveys, 2025). Đỗ và
Phan (2026 --- VEFR) đã mở rộng phạm vi WBES đến 17 nền kinh tế châu Á
mới nổi với \textasciitilde40.633 doanh nghiệp. Tuy nhiên, vẫn còn ít
tổng hợp xuyên \textbf{49 nền kinh tế} trên một giai đoạn dài. Chuyên đề
này lấp khoảng trống thực tiễn đó với \textbf{pool mô tả gồm 86.701
doanh nghiệp · 102 cặp nền kinh tế × năm (2006--2026)} --- gấp hơn 2 lần
phạm vi của Đỗ \& Phan (2026 --- VEFR) --- và khung phân loại thống nhất
với bộ dữ liệu ước lượng của luận án (91.982 doanh nghiệp / 49 nền; pool
phân loại 96.415 DN / 52 nền).

\textbf{Ẩn dụ minh họa --- Fiji và Singapore}: Trường hợp Fiji 2025 với
tỷ lệ website 74,8\% vượt Singapore 66,1\% là minh chứng rõ nhất cho
ranh giới giữa ``lý thuyết Uppsala vận hành ở điểm tối ưu'' và ``lý
thuyết Uppsala bị biến dạng ở điều kiện biên''. Đối với Fiji và 6 SIDS
Pacific khác, số hóa không phải công cụ tối ưu hóa chuỗi cung ứng mà là
phương tiện sinh tồn --- dẫn vào các luận điểm về điều kiện biên của mô
hình U-curve và tiểu cảnh điển hình SIDS Pacific trong phần 2.3.

\hypertarget{mux1ee5c-tiuxeau-chuyuxean-ux111ux1ec1}{%
\subsubsection{2.1.2 Mục tiêu chuyên
đề}\label{mux1ee5c-tiuxeau-chuyuxean-ux111ux1ec1}}

\textbf{Mục tiêu chung}: Hệ thống hóa và phân tích mô tả thực trạng hiệu
quả hoạt động kinh doanh của các doanh nghiệp ở \textbf{49 nền kinh tế
châu Á và Thái Bình Dương} (phạm vi chính châu Á + SIDS Thái Bình Dương
mở rộng trường hợp biên) trong giai đoạn \textbf{2006--2026}, trên pool
mô tả 86.701 doanh nghiệp · 102 cặp nền kinh tế × năm --- xem ghi chú dữ
liệu ở Tóm tắt.

\textbf{Mục tiêu cụ thể}:

\begin{enumerate}
\def\labelenumi{\arabic{enumi}.}
\tightlist
\item
  Hệ thống hóa khái niệm và đa chiều đo lường hiệu quả hoạt động kinh
  doanh trong văn liệu hiện hành;
\item
  Đề xuất khung phân loại 6 nhóm ICRV cho 49 nền kinh tế;
\item
  Phân tích thực trạng đa chiều hiệu quả doanh nghiệp ở 41 nền kinh tế
  châu Á đại lục và Đông Á (phạm vi chính) và so sánh với 8 nền SIDS
  (trường hợp biên);
\item
  Nhận diện đặc điểm và yếu tố giải thích sơ bộ cho dị biệt hiệu quả
  giữa các phân nhóm;
\item
  Đề xuất khoảng trống thực tiễn để định hướng nghiên cứu tiếp theo.
\end{enumerate}

\hypertarget{nux1ed9i-dung-vuxe0-phux1ea1m-vi-nghiuxean-cux1ee9u}{%
\subsubsection{2.1.3 Nội dung và phạm vi nghiên
cứu}\label{nux1ed9i-dung-vuxe0-phux1ea1m-vi-nghiuxean-cux1ee9u}}

\textbf{Đối tượng nghiên cứu}: Doanh nghiệp đang hoạt động ở các nền
kinh tế trong pool đáp ứng tiêu chí mẫu của WBES --- chủ yếu là doanh
nghiệp khu vực ngoài quốc doanh, có ít nhất 5 lao động, thuộc các ngành
phi nông nghiệp.

\textbf{Phạm vi không gian}:

\textbf{A. Phạm vi chính --- 41 nền kinh tế châu Á}: - \textbf{(i)
Advanced đổi mới sáng tạo dẫn dắt (5 nước)}: Singapore, Hong Kong SAR,
Hàn Quốc, Đài Loan, Israel. - \textbf{(ii) Advanced tài nguyên dẫn dắt
(6 nước)}: Saudi Arabia, Qatar, Kuwait, Bahrain, Brunei, Oman. -
\textbf{(iii) Trung bình cao --- Upper-middle (6 nước)}: Trung Quốc,
Malaysia, Thái Lan, Kazakhstan, Armenia, Georgia. - \textbf{(iv) Chuyển
đổi TB--thấp --- Lower\_mid\_transition (7 nước)}: Việt Nam, Ấn Độ,
Indonesia, Philippines, Mông Cổ, Bangladesh, Pakistan. - \textbf{(v)
Đang nổi --- Emerging (17 nước)}: Sri Lanka, Jordan, Lào, Campuchia,
Myanmar, Nepal, Bhutan, Maldives, Uzbekistan, Tajikistan, Kyrgyzstan,
Turkmenistan, Afghanistan, Azerbaijan, Iraq, Lebanon, Yemen.

\begin{quote}
\textbf{Ghi chú nhất quán nhãn nhóm:} Nhãn nhóm và danh sách thành viên
trong §2.1.3 và Bảng 2.3.2.1 thống nhất với biến phân loại
\texttt{icrv\_label} của bộ dữ liệu ước lượng (P7) và Chuyên đề 2: Nhóm
IV = \texttt{Lower\_mid\_transition} (7 nền đông dân), Nhóm V =
\texttt{Emerging} (17 nền), Nhóm VI = \texttt{SIDS\_small} (8 nền). Pool
mô tả phủ \textbf{đủ 49/49 nền}: Nhóm I 5/5, II 6/6 (đủ khối GCC), III
6/6, IV 7/7, V 17/17 (gồm Lebanon, Yemen), VI 8/8 (7 Pacific +
Timor-Leste).
\end{quote}

\textbf{B. Trường hợp biên mở rộng --- Nhóm VI SIDS (8 nền, n=2.295,
2,6\% pool)}: Fiji, Papua New Guinea, Solomon Islands, Tonga, Vanuatu,
Samoa, Kiribati (2025) và Timor-Leste. Mẫu chính của nghiên cứu chuyên
sâu SIDS là 7 nền Thái Bình Dương (n=1.781).

\textbf{Phạm vi thời gian}: \textbf{2006--2026} (19 năm khảo sát, 102
cặp nền kinh tế × năm). Ba thế hệ schema WBES trong pool mô tả:
PICS3/tiền-Standardized (2006--2012, 18 đợt, n=8.048), Standardized
(2013--2017, 27 đợt, n=23.940), BREADY/BEE/EAP Core (2018--2026, 57 đợt,
n=54.713 --- chiếm 63,1\% pool).

\hypertarget{giux1edbi-hux1ea1n-nghiuxean-cux1ee9u}{%
\subsubsection{2.1.4 Giới hạn nghiên
cứu}\label{giux1edbi-hux1ea1n-nghiuxean-cux1ee9u}}

Bảy giới hạn cấu trúc được thừa nhận minh bạch:

\textbf{(1) Cross-sectional không cho phép suy luận nhân quả mạnh}: WBES
là repeated cross-sections, không phải panel cân bằng. Kết luận nhân quả
dựa trên ước lượng biến công cụ từ nghiên cứu Việt Nam và không thể khái
quát trực tiếp sang 49 nền kinh tế trong khung phân tích.

\textbf{(2) Chỉ số áp dụng số bị giới hạn ở Tier-1 và Tier-2}: Bộ câu
hỏi PICS3/Standardized chỉ có website nhị phân (Tier-1). BREADY có
Tier-2 nhưng không đo Tier-3/4. Chuyên đề chỉ phát biểu về phần bù
Tier-1/2, không phải toàn bộ phổ áp dụng số.

\textbf{(3) Chỉ số TCI tổng hợp với trọng số đồng đều}: Trung bình đơn
giản 4 thành phần (ISO, R\&D, đổi mới sản phẩm, công nghệ nước ngoài)
chưa được xác nhận bằng phân tích nhân tố --- trọng số tối ưu có thể
khác nhau giữa các nhóm ICRV.

\textbf{(4) Thiếu một số biến lãnh đạo trong các đợt khảo sát cũ}: Biến
kinh nghiệm quản lý và giới tính lãnh đạo không có đủ trong tất cả các
đợt khảo sát WBES. Phân tích liên quan phải thực hiện trên mẫu con, tạo
ra khả năng lệch chọn lọc.

\textbf{(5) Không tách biệt tác động COVID-19 khỏi tác động áp dụng số}:
Bộ câu hỏi BREADY 2018--2025 bao gồm giai đoạn gián đoạn COVID
(2020--2022) và tăng tốc AI (2023+). Kết quả từ BREADY 2025 nên được
giải thích là ``hậu đại dịch kết hợp tăng tốc số''.

\textbf{(6) Phạm vi SIDS chưa đầy đủ}: 9/\textasciitilde52 SIDS theo
UNCTAD (dữ liệu WBES hiện có). Tuvalu, Nauru, Marshall Islands,
Micronesia chưa có WBES. Kết luận về chi phí quốc tế hóa buộc phải có
thể chưa đại diện toàn bộ SIDS Pacific.

\textbf{(7) Không có dữ liệu cường độ số theo ngành}: WBES không phân
biệt chỉ số áp dụng số theo ngành. Doanh nghiệp sản xuất và dịch vụ có
hệ số DAI Tier-1 giống nhau nhưng cường độ số thực tế khác nhau. Hiệu
ứng cố định ngành kiểm soát được một phần nhưng không giải quyết triệt
để vấn đề này.

\hypertarget{uxfd-nghux129a-khoa-hux1ecdc-vuxe0-thux1ef1c-tiux1ec5n}{%
\subsubsection{2.1.5 Ý nghĩa khoa học và thực
tiễn}\label{uxfd-nghux129a-khoa-hux1ecdc-vuxe0-thux1ef1c-tiux1ec5n}}

\textbf{Bốn đóng góp phạm vi chính (40 nước châu Á)}:

\begin{enumerate}
\def\labelenumi{(\arabic{enumi})}
\item
  \textbf{Bức tranh thực trạng đa chiều} cho hiệu quả doanh nghiệp châu
  Á 2006--2026 --- phạm vi rộng nhất trong văn liệu kinh doanh quốc tế
  hiện hành (Wu et al., 2022 dừng ở 2020; Arte \& Larimo, 2022 dừng ở
  2019).
\item
  \textbf{Phân nhóm con Advanced} (innovation-driven so với
  resource-driven) --- phát hiện lý thuyết mới chưa có trong văn liệu
  hiện hành. Hai nhóm cùng mức thu nhập cao nhưng cấu trúc năng lực đối
  lập: R\&D 20,5\% so với 7,0\%; doanh nghiệp xuất khẩu 28,4\% so với
  11,7\%; nữ quản lý cấp cao 27,5\% so với 4,6\% --- phản ánh hai cơ chế
  tăng trưởng hoàn toàn khác nhau.
\item
  \textbf{Nhảy vọt số} ở các nhóm đang nổi và chuyển đổi 2018--2026 ---
  mở rộng pattern từ 17 nền lên 49 nền. Phân tách Tier-1 (website nhị
  phân, tiệm cận bão hòa ở Advanced) và Tier-2 (tổng hợp đa thành phần,
  đang phát triển ở các nhóm IV--VI).
\item
  \textbf{Bằng chứng cho khung lý thuyết phi tuyến và điều tiết đa tầng}
  --- gradient cường độ rõ rệt của các hệ số tương quan xuyên 6 nhóm thể
  chế (đặc biệt: tương quan FSTS--năng suất giảm dần từ +0,163 ở Nhóm I
  xuống mất ý nghĩa ở SIDS) là cơ sở thực tiễn cho các giả thuyết của
  Chuyên đề 2.
\end{enumerate}

\textbf{Một đóng góp trường hợp biên mở rộng (SIDS Pacific)}:

\begin{enumerate}
\def\labelenumi{(\arabic{enumi})}
\setcounter{enumi}{4}
\tightlist
\item
  \textbf{Chi phí buộc phải quốc tế hóa} ở 7 SIDS Thái Bình Dương
  (Briguglio, 1995; Bertram, 2006) --- bằng chứng cấp doanh nghiệp đầu
  tiên trong văn liệu kinh doanh quốc tế. Kiribati 2025 là trường hợp
  cực đoan nhất với FSTS 1,0\%, FDI 0,7\%, website 18,7\%.
\end{enumerate}

\begin{center}\rule{0.5\linewidth}{0.5pt}\end{center}

\hypertarget{phux1b0ux1a1ng-phuxe1p-nghiuxean-cux1ee9u}{%
\subsection{2.2 PHƯƠNG PHÁP NGHIÊN
CỨU}\label{phux1b0ux1a1ng-phuxe1p-nghiuxean-cux1ee9u}}

Chuyên đề tiếp cận theo lối \textbf{mô tả -- chẩn đoán}, không thiết lập
mô hình hồi quy đa biến để kiểm định nhân quả. Phương pháp kết hợp ba kỹ
thuật: (a) thống kê mô tả đa biến (trung vị, trung bình, độ phân tán,
kurtosis) theo phân nhóm thể chế và theo ngành; (b) trực quan hóa dữ
liệu (heatmap, kernel density, spider chart, scatter plot --- 7 hình);
(c) phân tích hai biến.

\textbf{Nguồn dữ liệu}: World Bank Enterprise Surveys (WBES) --- cơ sở
dữ liệu vi mô doanh nghiệp cấp quốc gia lớn nhất toàn cầu, được thiết kế
theo phương pháp lấy mẫu xác suất phân tầng đảm bảo tính đại diện (World
Bank Enterprise Surveys, 2025). Các bảng thống kê mô tả §2.3.3--2.3.8
dưới đây được tính trên \textbf{pool mô tả gồm 86.701 doanh nghiệp · 49
nền kinh tế · 102 cặp nền kinh tế × năm (2006--2026)} --- mỗi cặp nền
kinh tế × năm lấy đúng một cross-section chuẩn, chỉ nhận đợt khảo sát từ
2006 trở đi (xem ghi chú dữ liệu ở Tóm tắt và Phụ lục A). Khung phân
loại thống nhất với bộ dữ liệu ước lượng của luận án: 49 nền kinh tế
(91.982 doanh nghiệp; pool phân loại 96.415 DN / 52 nền).

\textbf{Các biến đo lường chính}:

\begin{itemize}
\item
  \textbf{Biến năng suất lao động (LP)}:
  \texttt{LP\ =\ ln(doanh\ thu\ bán\ hàng\ hàng\ năm\ /\ số\ lao\ động\ toàn\ thời\ gian)}
  (d2/l1), tính theo nội tệ của từng nền kinh tế. Vì đơn vị tiền tệ khác
  nhau giữa các nước, chuyên đề tuân thủ một \textbf{nguyên tắc so sánh
  xuyên tiền tệ} xuyên suốt: (a) các thống kê \emph{phân tán} được tính
  \textbf{trong từng cặp nền kinh tế × năm} (bất biến với đơn vị tiền tệ
  vì thừa số quy đổi triệt tiêu trong thang log); (b) so sánh \emph{mức}
  giữa các nền kinh tế dùng chỉ số không thứ nguyên (ROS --- tỷ suất
  sinh lời trên doanh thu); (c) trong khung ước lượng của luận án, LP
  được chuẩn hóa z-score trong từng cặp nền kinh tế × năm và mô hình
  luôn có hiệu ứng cố định quốc gia (Phụ lục A của luận án). Giá trị LP
  được winsorize ở mức 1/99 phân vị trong cụm quốc gia × năm để kiểm
  soát ngoại lai.
\item
  \textbf{Biến FSTS (Foreign Sales to Total Sales)}:
  \texttt{FSTS\ =\ (d3b\ +\ d3c)\ /\ 100}, trong đó d3b là tỷ lệ doanh
  thu \textbf{xuất khẩu gián tiếp} (bán qua trung gian xuất khẩu) và d3c
  là tỷ lệ doanh thu \textbf{xuất khẩu trực tiếp} --- theo nhãn biến
  chính thức của bộ công cụ WBES. Quan sát được gán FSTS = 0 nếu doanh
  nghiệp báo cáo toàn bộ doanh thu nội địa (d3a = 100); gán khuyết nếu
  không trả lời. Tổng xuất khẩu trực tiếp + gián tiếp là định nghĩa tham
  gia xuất khẩu chuẩn của WBES (World Bank Enterprise Surveys, 2025) và
  được dùng thống nhất cho thống kê mô tả của chuyên đề cũng như khung
  ước lượng đa quốc gia của luận án. \textbf{Ghi chú thao tác hóa}: hai
  nghiên cứu thành phần đơn quốc gia (P3 Việt Nam, P5 Trung Quốc) đo
  cường độ quốc tế hóa bằng riêng xuất khẩu trực tiếp (d3c), vì xuất
  khẩu gián tiếp ủy thác chức năng quốc tế hóa cho trung gian thương mại
  --- doanh nghiệp không trực tiếp tích lũy kinh nghiệm thị trường nước
  ngoài, nên cơ chế học hỏi và yêu cầu nguồn lực khác về bản chất
  (Hessels \& Terjesen, 2010; Peng \& Ilinitch, 1998); trong khi tỷ
  trọng doanh thu xuất khẩu (FSTS) là thước đo cường độ quốc tế hóa
  chuẩn của văn liệu I--P (Lu \& Beamish, 2001; Sullivan, 1994). Hai
  cách thao tác hóa trả lời hai câu hỏi bổ trợ --- mức độ \emph{tham
  gia} thị trường quốc tế (tổng) và mức độ \emph{cam kết trực tiếp} của
  doanh nghiệp (trực tiếp) --- và được ghi rõ trong phần dữ liệu của
  từng nghiên cứu.
\item
  \textbf{Biến TCI (Technological Capability Index)}: Trung bình chuẩn
  hóa z của tối thiểu 3 trong 4 thành phần: ISO (b8), R\&D (h8), đổi mới
  sản phẩm (h1), công nghệ nước ngoài (e6). Yêu cầu ≥3 non-missing để
  tính tổng hợp.
\item
  \textbf{Biến DAI (Digital Adoption Index)}: Spec 1 (toàn bộ
  2006--2026): website nhị phân (c22b/e8) --- \textbf{Tier-1 Sự hiện
  diện số cơ bản}. Spec 2 (BREADY 2018+ trở lên): tổng hợp chuẩn hóa z
  website + cường độ thanh toán điện tử (k33/k38) --- \textbf{Tier-2
  Chuyển đổi số cơ bản}.
\end{itemize}

\textbf{Hài hòa hóa xuyên ba thế hệ schema}: World Bank Enterprise
Surveys trải qua ba thế hệ schema trong giai đoạn nghiên cứu. Thế hệ 1
--- PICS3/MENA-WBES (2006--2012, 18 đợt, n=8.048): bảng hỏi PICS3 cho
khu vực Đông Nam Á và Trung Á. Thế hệ 2 --- Standardized (2013--2017, 27
đợt, n=23.940): áp dụng đồng nhất xuyên toàn cầu từ 2013. Thế hệ 3 ---
BREADY/BEE/EAP Core (2018--2026, 57 đợt, n=54.713): tích hợp thanh toán
điện tử, điện toán đám mây, ngân hàng di động vào các module bổ sung,
gây \textbf{đứt gãy schema} quan sát được: Ấn Độ FSTS sụt 7,7\%→2,7\%
trong đợt 2025.

\textbf{Nguyên tắc minh bạch dữ liệu} (Aguinis et al., 2019): chuyên đề
trình bày rõ (i) định nghĩa tổng thể (49 nền kinh tế châu Á và Thái Bình
Dương, phân tách 41 phạm vi chính và 8 trường hợp biên); (ii) quy trình
lấy mẫu (lấy mẫu chính thức WBES); (iii) lý giải quy mô mẫu (pool mô tả
86.701 DN / 102 cặp; khung ước lượng 91.982 DN; pool phân loại 96.415
DN); (iv) xử lý giá trị thiếu (FSTS = d3b + d3c, winsorize 1/99); (v)
xây dựng biến; (vi) độ tin cậy và giá trị (TCI đa thành phần; DAI đơn
thành phần Spec 1, đa thành phần Spec 2); (vii) phương pháp ước lượng
(mô tả trong CĐ1; OLS và biến công cụ trong CĐ2); (viii) khoa học mở ---
toàn bộ bảng và hình tái lập được từ gói tái lập kèm theo luận án.

\textbf{Hệ thống phòng thủ ba tầng cho đứt gãy schema BREADY 2025}: Đợt
2025 (12 nền kinh tế, n=14.712 doanh nghiệp, 17,0\% pool) sử dụng schema
mới BREADY gây biến động bất thường: Ấn Độ FSTS giảm từ 7,7\% xuống
2,7\%; R\&D nhiều nước tăng đột biến do hiệu ứng bảng hỏi chứ không phải
thực tế. Ba đề xuất phương pháp: (3a) biến giả
\texttt{Post\_BREADY\_2024} hấp thụ hiệu ứng schema tĩnh; (3b) mô hình
neo (anchor model) --- chạy hồi quy với dữ liệu đến 2024, khóa hệ số,
kiểm định ổn định với dữ liệu đầy đủ bằng Chow test; (3c) panel hậu đại
dịch độc lập --- tách 2025 thành tập xác thực ngoài mẫu cho kỷ nguyên
hậu COVID và AI.

\begin{center}\rule{0.5\linewidth}{0.5pt}\end{center}

\hypertarget{nux1ed9i-dung-nghiuxean-cux1ee9u}{%
\subsection{2.3 NỘI DUNG NGHIÊN
CỨU}\label{nux1ed9i-dung-nghiuxean-cux1ee9u}}

\hypertarget{cux1a1-sux1edf-luxfd-luux1eadn-vux1ec1-hiux1ec7u-quux1ea3-houx1ea1t-ux111ux1ed9ng-kinh-doanh}{%
\subsubsection{2.3.1 Cơ sở lý luận về hiệu quả hoạt động kinh
doanh}\label{cux1a1-sux1edf-luxfd-luux1eadn-vux1ec1-hiux1ec7u-quux1ea3-houx1ea1t-ux111ux1ed9ng-kinh-doanh}}

\hypertarget{khuxe1i-niux1ec7m-hiux1ec7u-quux1ea3-houx1ea1t-ux111ux1ed9ng-kinh-doanh}{%
\paragraph{2.3.1.1 Khái niệm hiệu quả hoạt động kinh
doanh}\label{khuxe1i-niux1ec7m-hiux1ec7u-quux1ea3-houx1ea1t-ux111ux1ed9ng-kinh-doanh}}

Hiệu quả hoạt động kinh doanh (firm performance) là khái niệm \textbf{đa
chiều} phản ánh mức độ đạt được các mục tiêu chiến lược của doanh nghiệp
theo thời gian và trong bối cảnh thể chế cụ thể. Venkatraman và
Ramanujam (1986) phân biệt hai tầng: (a) \textbf{hiệu quả tài chính} ---
doanh thu, lợi nhuận, tỷ suất sinh lợi; và (b) \textbf{hiệu quả hoạt
động rộng hơn} --- thị phần, đổi mới, năng suất, tăng trưởng việc làm.
Lý thuyết dựa trên nguồn lực (Barney, 1991; Wernerfelt, 1984) xem hiệu
quả doanh nghiệp là kết quả của việc tích lũy và khai thác \textbf{nguồn
lực chiến lược} (có giá trị, hiếm, không thể sao chép, không thể thay
thế --- VRIN), trong khi quan điểm thể chế (North, 1990; Scott, 1995;
Peng, 2001) nhấn mạnh rằng hiệu quả này bị định hình mạnh bởi
\textbf{môi trường thể chế} nơi doanh nghiệp hoạt động.

Vượt ra ngoài khung hai tầng của Venkatraman và Ramanujam (1986), văn
liệu phương pháp luận về đo lường hiệu quả tổ chức (Combs et al., 2005;
Richard et al., 2009) cho phép tổng hợp thành \textbf{ba nhóm thước đo
chính}: (i) \textbf{thước đo dựa trên kế toán} (ROA, ROE, ROS) phản ánh
hiệu quả tài chính lịch sử; (ii) \textbf{thước đo dựa trên thị trường
tài chính} (Tobin's Q, tổng lợi nhuận cổ đông) phản ánh kỳ vọng tương
lai do nhà đầu tư định giá; và (iii) \textbf{thước đo dựa trên năng suất
và tăng trưởng} (năng suất lao động, tăng trưởng doanh thu và việc làm)
phản ánh năng lực vận hành thực chất. Richard và cộng sự (2009) cảnh báo
rằng việc dựa vào một thước đo đơn lẻ làm suy giảm hiệu lực cấu trúc
(construct validity) và khuyến nghị đo lường đa nguồn, đa chiều. Trong
bối cảnh dữ liệu WBES --- vốn không cung cấp giá cổ phiếu nên loại trừ
nhóm (ii), và có dữ liệu lợi nhuận thưa thớt làm hạn chế nhóm (i) ---
\textbf{năng suất lao động} (thuộc nhóm iii) là lựa chọn chủ đạo có cơ
sở lý thuyết và kỹ thuật vững nhất giữa các phương án (ROA/ROS/Tobin's
Q): nó so sánh được xuyên quốc gia, ít bị méo mó bởi cơ cấu vốn và chế
độ thuế đặc thù từng quốc gia, và phản ánh trực tiếp năng lực chuyển hóa
đầu vào thành đầu ra. Để bù đắp giới hạn của một thước đo đơn lẻ theo
khuyến nghị của Combs và cộng sự (2005), chuyên đề bổ sung ba chiều hiệu
quả hỗ trợ (lợi nhuận, tăng trưởng, đổi mới) được trình bày ở mục
2.3.1.2.

Trong bối cảnh nghiên cứu xuyên quốc gia, khái niệm hiệu quả cần được
\textbf{tiêu chuẩn hóa} để so sánh được. Theo Cusolito và Maloney
(2018), \textbf{năng suất lao động} --- đo bằng giá trị gia tăng trên
lao động hoặc doanh thu trên lao động --- là thước đo phổ biến nhất
trong các nghiên cứu WBES vì: (a) tính khả dụng cao xuyên 3 thế hệ
schema; (b) ít bị biến dạng bởi cơ cấu vốn hay cơ chế thuế quốc gia; (c)
phản ánh trực tiếp năng lực chuyển đổi đầu vào thành đầu ra. Trong
chuyên đề này, năng suất lao động được tính bằng:

\textbf{LP = ln(Doanh thu bán hàng hàng năm PPP / Số lao động toàn thời
gian)}

Giá trị LP được winsorize ở mức 1/99 phân vị trong cụm quốc gia × năm để
kiểm soát ngoại lai và được điều chỉnh theo sức mua tương đương nhằm đảm
bảo tính so sánh xuyên quốc gia (World Bank Enterprise Surveys, 2025).

\hypertarget{bux1ed1n-chiux1ec1u-ux111o-lux1b0ux1eddng-hiux1ec7u-quux1ea3}{%
\paragraph{2.3.1.2 Bốn chiều đo lường hiệu
quả}\label{bux1ed1n-chiux1ec1u-ux111o-lux1b0ux1eddng-hiux1ec7u-quux1ea3}}

Hiệu quả doanh nghiệp trong chuyên đề này được đo lường theo \textbf{bốn
chiều} tương ứng với bốn nhóm biến WBES:

\textbf{Chiều 1 --- Năng suất lao động}. Biến đo lường chính: LP =
ln(doanh thu hàng năm PPP / số lao động toàn thời gian). Đây là
\textbf{thước đo chính} xuyên suốt phân tích, được lựa chọn vì: (a) có
thể tính cho hầu hết các nền kinh tế trong pool; (b) không bị nhiễu bởi
cơ cấu thuế và chính sách chuyển giá quốc gia; (c) có giá trị tham chiếu
rõ ràng trong văn liệu quốc tế (Hsieh \& Klenow, 2009, 2014).

\textbf{Chiều 2 --- Lợi nhuận và biên lợi nhuận}. Biến đo lường: (a) tỷ
suất lợi nhuận hoạt động từ câu hỏi WBES về doanh thu sau chi phí lao
động; (b) chỉ số gián tiếp qua khả năng tái đầu tư. Hạn chế: nhiều doanh
nghiệp trong WBES không tiết lộ lợi nhuận cụ thể --- chiều này mang tính
hỗ trợ thay vì chủ đạo.

\textbf{Chiều 3 --- Tăng trưởng}. Biến đo lường: (a) tốc độ tăng trưởng
việc làm hàng năm kép giữa 3 năm trước và thời điểm khảo sát; (b) tỷ lệ
doanh nghiệp xuất khẩu như proxy cho mở rộng thị trường. Tăng trưởng
việc làm được ưu tiên vì: (a) ít nhạy cảm với thay đổi giá cả so với
doanh thu; (b) phản ánh khả năng tạo việc làm bền vững --- quan trọng
cho mục tiêu chính sách.

\textbf{Chiều 4 --- Đổi mới sáng tạo và cấu trúc doanh nghiệp}. Biến đo
lường: (a) R\&D (tỷ lệ doanh nghiệp có chi tiêu R\&D dương); (b) đổi mới
sản phẩm mới; (c) đổi mới quy trình mới; (d) chứng nhận ISO; (e) website
(Tier-1 Sự hiện diện số cơ bản); (f) FDI ≥10\% (cơ cấu sở hữu). Chiều
thứ tư phản ánh \textbf{năng lực tương lai} --- khả năng duy trì hiệu
quả trong dài hạn qua đổi mới và nâng cấp công nghệ.

\hypertarget{nux103m-meta-analysis-lux1edbn-vux1ec1-quan-hux1ec7-quux1ed1c-tux1ebf-huxf3a-hiux1ec7u-quux1ea3-19802024}{%
\paragraph{2.3.1.3 Năm meta-analysis lớn về quan hệ quốc tế hóa -- hiệu
quả
(1980--2024)}\label{nux103m-meta-analysis-lux1edbn-vux1ec1-quan-hux1ec7-quux1ed1c-tux1ebf-huxf3a-hiux1ec7u-quux1ea3-19802024}}

Văn liệu về quan hệ quốc tế hóa → hiệu quả là một trong những dòng
nghiên cứu tích lũy nhiều nhất trong lý thuyết kinh doanh quốc tế. Năm
meta-analysis lớn cung cấp nền tảng thực nghiệm:

\textbf{(1) Bausch \& Krist (2007)}: Tổng hợp \textbf{68 nghiên cứu}
giai đoạn 1980--2005. Hệ số tương quan trung bình r=0,045 --- tích cực
nhưng nhỏ. Phát hiện chính: đo lường quốc tế hóa và hiệu quả rất không
đồng nhất giữa các nghiên cứu, hạn chế khả năng so sánh. Khoảng trống:
ít nghiên cứu châu Á mới nổi.

\textbf{(2) Kirca et al.~(2012)}: Tổng hợp \textbf{154 mẫu} với phân
tích điều tiết meta-analytic. Phát hiện: quan hệ quốc tế hóa → hiệu quả
dương và có ý nghĩa thống kê; nhưng bị điều tiết bởi loại doanh nghiệp,
phương thức đo lường, bối cảnh quốc gia. Quan trọng: các doanh nghiệp
châu Á và các nền kinh tế mới nổi biểu hiện pattern khác biệt so với
doanh nghiệp OECD.

\textbf{(3) Marano et al.~(2016)}: Tập trung vào thể chế hóa nghiên cứu
quốc tế hóa → hiệu quả. Phát hiện: chất lượng thể chế nước chủ nhà là
biến điều tiết quan trọng nhất --- giải thích phần lớn phương sai không
giải thích được trong các meta-analysis trước. Đặt nền tảng cho khung
ICRV trong chuyên đề này.

\textbf{(4) Schwens et al.~(2018)}: Tổng hợp các yếu tố điều tiết mức độ
tác động của quốc tế hóa. Phát hiện: năng lực hấp thụ (Cohen \&
Levinthal, 1990), chất lượng lãnh đạo, và phát triển thể chế là các biến
điều tiết quan trọng nhất.

\textbf{(5) Wu, Xu và Yang (2022)}: Cập nhật meta-analysis đến 2020, tập
trung vào \textbf{doanh nghiệp đa quốc gia mới nổi châu Á}. Phát hiện:
quan hệ quốc tế hóa → hiệu quả ở doanh nghiệp mới nổi châu Á có hình
dạng U ngược rõ nét hơn so với doanh nghiệp OECD; điểm uốn thường nằm
trong phạm vi 30--50\% FSTS.

\textbf{Khoảng trống từ 5 meta-analyses}: Không một meta-analysis nào
trong số trên có dữ liệu WBES xuyên 49 nền kinh tế với phân loại ICRV 6
nhóm. Không có nghiên cứu nào phân biệt tường minh Tier-1 và Tier-2 áp
dụng số như biến điều tiết riêng biệt. Khoảng trống này là nền tảng cho
phần thảo luận về các khoảng trống nghiên cứu trong mục 2.3.9.

\hypertarget{ba-khung-luxfd-thuyux1ebft-nux1ec1n}{%
\paragraph{2.3.1.4 Ba khung lý thuyết
nền}\label{ba-khung-luxfd-thuyux1ebft-nux1ec1n}}

Ba khung lý thuyết cung cấp nền tảng giải thích cho các pattern thực
trạng được quan sát:

\textbf{Khung 1 --- Lý thuyết Uppsala} (Johanson \& Vahlne, 1977, 2009):
Doanh nghiệp quốc tế hóa theo \textbf{lộ trình tiệm tiến}, bắt đầu từ
các thị trường gần về tâm lý rồi mở rộng dần. Phiên bản 2009 cập nhật:
mạng lưới quan hệ kinh doanh thay thế khoảng cách tâm lý như cơ chế dẫn
dắt. Hàm ý cho chuyên đề: (a) pattern FSTS thấp ở SIDS Pacific (6,3\%)
và FSTS cao liên kết FDI ở Việt Nam (23,2\%) phù hợp với đường cong học
hỏi Uppsala; (b) các công cụ số theo Banalieva \& Dhanaraj (2019) có thể
rút ngắn lộ trình tiệm tiến ở giai đoạn đầu --- nhưng cần Tier-2 để phát
huy đầy đủ.

\textbf{Khung 2 --- Quan điểm dựa trên nguồn lực và năng lực động}
(Barney, 1991; Teece et al., 1997; Wernerfelt, 1984): Doanh nghiệp đạt
lợi thế cạnh tranh bền vững thông qua nguồn lực VRIN và khả năng tái cấu
hình năng lực theo biến động môi trường. Hàm ý cho chuyên đề: (a) TCI
(năng lực công nghệ) là proxy cho nguồn lực khó sao chép --- giải thích
tại sao Advanced innovation-driven (ISO 29,9\%, R\&D 16,7\%) có năng
suất vượt trội; (b) năng lực hấp thụ (Cohen \& Levinthal, 1990) giải
thích tại sao tác động lan tỏa FDI tích cực hơn ở Ấn Độ với R\&D 19,8\%
so với Việt Nam 6,1\%.

\textbf{Khung 3 --- Lý thuyết thể chế} (North, 1990; Scott, 1995; Peng,
2001, 2003): Thể chế --- luật pháp, chuẩn mực, quy tắc phi chính thức
--- định hình chi phí giao dịch và cơ hội chiến lược cho doanh nghiệp.
Peng (2001, 2003) xây dựng \textbf{khung thể chế dựa trên giao dịch} cho
các nền kinh tế mới nổi --- lý giải tại sao quan hệ quốc tế hóa → hiệu
quả không đồng nhất xuyên chế độ thể chế. Hàm ý: (a) gradient cường độ
của các hệ số tương quan FDI/FSTS/TCI/DAI xuyên 6 nhóm ICRV --- đặc biệt
việc tương quan FSTS--năng suất giảm dần và mất ý nghĩa ở SIDS (Bảng
2.3.8.1) --- là bằng chứng mô tả cho điều tiết thể chế; (b) Xu (2024)
làm sâu thêm bằng phân biệt \textbf{de jure} (luật trên giấy) so với
\textbf{de facto} (thực thi thực tế) --- giải thích tại sao Việt Nam và
Trung Quốc có gap lớn giữa hai tầng này; (c) khung ICRV 6 nhóm là sự vận
hành hóa Lý thuyết thể chế theo hướng này.

\hypertarget{ux111iux1ec1u-kiux1ec7n-biuxean-cho-u-curve-lu-beamish-2004}{%
\paragraph{2.3.1.5 Điều kiện biên cho U-curve Lu \& Beamish
(2004)}\label{ux111iux1ec1u-kiux1ec7n-biuxean-cho-u-curve-lu-beamish-2004}}

\textbf{Ba điều kiện cấu trúc bắt buộc cho mô hình U-curve vận hành}:

\textbf{(1) Điều kiện thị trường nội địa khả thi}: Doanh nghiệp có thể
bán hàng nội địa ở quy mô đủ để tích lũy năng lực kinh doanh trước khi
quốc tế hóa. Điều kiện này thỏa mãn ở Singapore (5,9 triệu dân, GDP bình
quân đầu người \textasciitilde80.000 USD), Trung Quốc (1,4 tỷ dân), Việt
Nam (98 triệu dân) --- \textbf{không} thỏa mãn ở Kiribati (130.000 dân,
GDP bình quân đầu người \textasciitilde1.700 USD).

\textbf{(2) Điều kiện chi phí thương mại kiểm soát được}: Logistics, vận
tải, quy định xuất nhập khẩu nằm trong khoảng kinh tế. Điều kiện này
thỏa mãn ở châu Á đại lục --- \textbf{không} thỏa mãn đầy đủ ở SIDS
Pacific xa hub thương mại (Fiji--Auckland 2.100 km, Kiribati--Suva 2.700
km).

\textbf{(3) Điều kiện thể chế vận hành}: Khung pháp lý, quyền sở hữu trí
tuệ, thực thi hợp đồng, hạ tầng ngân hàng ở mức tối thiểu. Điều kiện này
thỏa mãn ở Advanced và Trung bình cao --- \textbf{không} thỏa mãn đầy đủ
ở SIDS cô lập (Kiribati ISO 1,4\%, FDI 0,7\%) hay các nền đang qua xung
đột thuộc Nhóm V; nhóm Chuyển đổi TB-thấp có gap de jure--de facto lớn
(Xu, 2024).

\textbf{Bảng 2.3.1.1}. \emph{Ma trận đối chiếu 3 trạng thái U-curve theo
3 điều kiện cấu trúc.}

\begin{longtable}[]{@{}
  >{\raggedright\arraybackslash}p{(\columnwidth - 10\tabcolsep) * \real{0.1250}}
  >{\raggedright\arraybackslash}p{(\columnwidth - 10\tabcolsep) * \real{0.1250}}
  >{\centering\arraybackslash}p{(\columnwidth - 10\tabcolsep) * \real{0.2083}}
  >{\centering\arraybackslash}p{(\columnwidth - 10\tabcolsep) * \real{0.2083}}
  >{\centering\arraybackslash}p{(\columnwidth - 10\tabcolsep) * \real{0.2083}}
  >{\raggedright\arraybackslash}p{(\columnwidth - 10\tabcolsep) * \real{0.1250}}@{}}
\toprule\noalign{}
\begin{minipage}[b]{\linewidth}\raggedright
Trạng thái
\end{minipage} & \begin{minipage}[b]{\linewidth}\raggedright
Đại diện
\end{minipage} & \begin{minipage}[b]{\linewidth}\centering
ĐK 1 (thị trường)
\end{minipage} & \begin{minipage}[b]{\linewidth}\centering
ĐK 2 (chi phí TM)
\end{minipage} & \begin{minipage}[b]{\linewidth}\centering
ĐK 3 (thể chế)
\end{minipage} & \begin{minipage}[b]{\linewidth}\raggedright
Hành vi U-curve
\end{minipage} \\
\midrule\noalign{}
\endhead
\bottomrule\noalign{}
\endlastfoot
\textbf{A. Điểm tối ưu của U-curve} & Singapore, Hong Kong, Hàn Quốc,
Đài Loan & ✓ & ✓ & ✓ & U-curve vận hành chuẩn --- FSTS có mức tối ưu cho
hiệu quả tối đa \\
\textbf{B. Bẫy chuyển hướng nguồn lực} & Việt Nam, Indonesia, Trung
Quốc, Mông Cổ & ✓ & ✓ & ⚠ \emph{(gap de jure--de facto)} & U-curve vận
hành nhưng có bẫy chuyển hướng nguồn lực --- cần phân biệt quy tắc hình
thức và năng lực thực thi \\
\textbf{C. Sụp đổ đường cong} & SIDS Pacific (Fiji, Kiribati, Tonga) & ✗
& ✗ & ⚠--✗ & U-curve không vận hành --- quốc tế hóa là buộc phải; Fiji
website 74,8\% \textgreater{} Singapore 66\% là minh chứng \\
\end{longtable}

\hypertarget{khung-4-tux1ea7ng-chuyux1ec3n-ux111ux1ed5i-sux1ed1-vuxe0-muxf4-huxecnh-cdcm}{%
\paragraph{2.3.1.6 Khung 4 Tầng chuyển đổi số và Mô hình
CDCM}\label{khung-4-tux1ea7ng-chuyux1ec3n-ux111ux1ed5i-sux1ed1-vuxe0-muxf4-huxecnh-cdcm}}

\textbf{Bốn tầng bậc số hóa theo Verhoef et al.~(2021)}:

\begin{longtable}[]{@{}
  >{\raggedright\arraybackslash}p{(\columnwidth - 6\tabcolsep) * \real{0.2500}}
  >{\raggedright\arraybackslash}p{(\columnwidth - 6\tabcolsep) * \real{0.2500}}
  >{\raggedright\arraybackslash}p{(\columnwidth - 6\tabcolsep) * \real{0.2500}}
  >{\raggedright\arraybackslash}p{(\columnwidth - 6\tabcolsep) * \real{0.2500}}@{}}
\toprule\noalign{}
\begin{minipage}[b]{\linewidth}\raggedright
Tầng
\end{minipage} & \begin{minipage}[b]{\linewidth}\raggedright
Tên (Tiếng Việt)
\end{minipage} & \begin{minipage}[b]{\linewidth}\raggedright
Đo lường WBES
\end{minipage} & \begin{minipage}[b]{\linewidth}\raggedright
Quan sát được?
\end{minipage} \\
\midrule\noalign{}
\endhead
\bottomrule\noalign{}
\endlastfoot
\textbf{Tier 1} & Sự hiện diện số cơ bản & website nhị phân, email & Có
(toàn pool 2006--2026) \\
\textbf{Tier 2} & Giao tiếp số và thương mại cơ bản & thanh toán điện tử
(k33, k38) & Có một phần (BREADY 2018+) \\
\textbf{Tier 3} & Tích hợp quy trình số & ERP, CRM & Không quan sát
được \\
\textbf{Tier 4} & Năng lực số động & AI, platform & Không quan sát
được \\
\end{longtable}

\textbf{Tuyên bố ranh giới đo lường}: Dữ liệu WBES của chuyên đề (pool
mô tả 86.701 DN; khung ước lượng 91.982 DN / 49 nền) quan sát Tier 1--2,
không Tier 3--4. Mọi lập luận về ``năng lực số'' phải hạn chế ở phạm vi
\textbf{Áp dụng số nền tảng (Foundational Digital Adoption --- FDA)}.

\textbf{Mô hình Năng lực số phụ thuộc bối cảnh (CDCM)} --- hợp nhất bằng
chứng Singapore và Việt Nam:

\textbf{(a) Hình dạng đường cong quốc tế hóa -- hiệu quả}: Việt Nam → U
ngược, điểm uốn ≈ 39--46\% FSTS; Singapore → tuyến tính dương + bậc hai
nhẹ, điểm uốn ≈ 88,6\% FSTS (điểm uốn này nằm sát/ngoài miền dữ liệu với
khoảng tin cậy bootstrap rất rộng, nên được diễn giải là quan hệ
\emph{tăng gần đơn điệu} trong miền quan sát chứ không phải một ngưỡng
bất lợi đã định vị chắc chắn --- nhất quán với diễn giải tại luận án
Chương 4 §4.3.1).

\textbf{(b) Cơ chế DAI}: Việt Nam (Tier 1 = website only) → \textbf{phụ
thuộc giai đoạn}: tương tác DAI×FSTS âm ở FSTS cao --- Tier 1 trở thành
điểm nghẽn. Singapore (Tier 1+2) → \textbf{mở rộng có điều kiện}: tương
tác DAI×FSTS² dương mạnh (β=3,119, p=0,005) --- Tier 1+2 là đòn bẩy.

\textbf{(c) Vai trò TCI}: Việt Nam --- \emph{lợi thế khan hiếm} (TCI z
biến công cụ β=1,639, p\textless0,001, ước lượng 2SLS vững). Singapore
--- \emph{yếu tố vệ sinh} nâng sàn năng suất nhưng không điều tiết đường
cong quốc tế hóa -- hiệu quả.

\begin{longtable}[]{@{}
  >{\raggedright\arraybackslash}p{(\columnwidth - 0\tabcolsep) * \real{0.0556}}@{}}
\toprule\noalign{}
\endhead
\bottomrule\noalign{}
\endlastfoot
\#\#\# 2.3.2 Khung phân loại thể chế ICRV \\
\#\#\#\# 2.3.2.1 Lý do phân loại \\
Châu Á không phải là một khối đồng nhất về thể chế. Các meta-analysis
(Marano et al., 2016; Wu et al., 2022) đã ghi nhận rằng bối cảnh thể chế
là biến điều tiết quan trọng nhất trong quan hệ quốc tế hóa → hiệu quả
--- nhưng không có nghiên cứu nào thiết lập hệ thống phân loại thể chế
toàn diện cho 49 nền kinh tế châu Á (khung phân tích). Khung phân loại
ICRV (Institutional Context Regime Variation --- Biến thiên chế độ bối
cảnh thể chế) được đề xuất trong chuyên đề này nhằm lấp khoảng trống đó,
dựa trên bốn nguyên tắc: \\
(a) \textbf{Tính phân biệt}: mỗi nhóm phải có mô hình hiệu quả doanh
nghiệp khác biệt thực sự --- không chỉ khác nhau về mức GDP bình quân
đầu người; \\
(b) \textbf{Tính đại diện lý thuyết}: phân loại phải phản ánh được các
cơ chế thể chế khác nhau --- thể chế đổi mới sáng tạo, thể chế tài
nguyên, thể chế chuyển đổi, thể chế MIRAB --- theo Varieties of
Capitalism (Hall \& Soskice, 2001) và Lý thuyết thể chế (North,
1990); \\
(c) \textbf{Tính khả dụng dữ liệu}: mỗi nhóm phải có đủ quan sát WBES để
phân tích có ý nghĩa thống kê (tối thiểu 500 doanh nghiệp); \\
(d) \textbf{Tính minh bạch và tái tạo}: tiêu chí phân loại phải rõ ràng
và có thể áp dụng nhất quán. \\
\#\#\#\# 2.3.2.2 Tiêu chí phân loại sáu nhóm \\
\textbf{Nhóm I --- Advanced đổi mới sáng tạo dẫn dắt}: WGI Rule of Law
\textgreater+0,80; GNI bình quân đầu người \textgreater30.000 USD
(phương pháp Atlas WB); GDP tăng trưởng dựa trên đổi mới công nghệ, dịch
vụ tài chính và hàm lượng tri thức cao. Bằng chứng WBES: R\&D
\textgreater10\%, ISO \textgreater20\%, độ lệch chuẩn log năng suất
trong đợt thấp nhất sáu nhóm (\textasciitilde1,04 --- phân tán hẹp, ít
phân bổ sai). \\
\textbf{Nhóm II --- Advanced tài nguyên dẫn dắt}: WGI Rule of Law trung
bình (+0,00 → +0,50); GNI bình quân đầu người \textgreater20.000 USD;
GDP tăng trưởng dựa trên xuất khẩu tài nguyên thiên nhiên (dầu mỏ, khí
đốt, khoáng sản). Bằng chứng WBES: FDI cao hướng khai thác tài nguyên;
phân tán năng suất trong đợt thấp ở các nền lớn nhất khối (Saudi Arabia
0,49; Qatar 0,33 --- đặc trưng nhà nước tô san bằng phân phối). \\
\textbf{Nhóm III --- Trung bình cao (Upper-middle)}: WGI Rule of Law
0,00 → +0,80; GNI bình quân đầu người 5.000--20.000 USD; nền kinh tế
chuyển đổi từ sản xuất sang dịch vụ, với R\&D đang tăng (đặc biệt Trung
Quốc). Pattern WBES: FSTS sản xuất cao, doanh nghiệp xuất khẩu tăng mạnh
giữa hai kỳ khảo sát (+12,5 điểm phần trăm). \\
\textbf{Nhóm IV --- Chuyển đổi thu nhập trung bình--thấp (Lower-middle
transition)}: WGI Rule of Law -0,50 → 0,00; GNI bình quân đầu người
1.000--5.000 USD; tốc độ tăng trưởng cao nhưng gap de jure--de facto lớn
(Xu, 2024); sản xuất FDI-driven dẫn dắt xuất khẩu ở khối ASEAN. Pattern
WBES: FSTS phân cực cao/thấp trong nội bộ nhóm; FDI thấp nhất (3,1\%) do
doanh nghiệp nội địa chi phối mẫu. \\
\textbf{Nhóm V --- Đang nổi (Emerging)}: WGI Rule of Law \textless{}
-0,50 hoặc voids thể chế lớn; GNI bình quân đầu người phần lớn
\textless3.000 USD; gồm nhiều nền đang qua giai đoạn xung đột hoặc
chuyển đổi sau xung đột. Pattern WBES: phân tán năng suất trong đợt cao
nhất (sd \textasciitilde1,39; P90/P10 \textasciitilde30 lần), nhưng đổi
mới sản phẩm khá cao (đổi mới bằng nguồn lực hạn chế). \\
\textbf{Nhóm VI --- SIDS Thái Bình Dương (Trường hợp biên mở rộng)}: Dân
số \textless1 triệu; cô lập địa lý \textgreater2.000 km từ hub thương
mại; phụ thuộc ODA + remittances \textgreater{} 20\% GDP; mô hình MIRAB
(Bertram, 2006). Pattern WBES: n=2.295 (2,6\% pool); FDI cao nhất
(20,0\% --- do MNE du lịch/viễn thông); đổi mới sản phẩm cao nhất
(34,2\% --- đổi mới bằng nguồn lực hạn chế); website khá cao ngoại trừ
Kiribati (18,7\%). \\
\#\#\#\# 2.3.2.3 Danh sách 49 nền kinh tế theo nhóm ICRV (căn theo nhãn
dữ liệu) \\
\textbf{Bảng 2.3.2.1}. \emph{Phân loại 49 nền kinh tế vào 6 nhóm ICRV
theo nhãn dữ liệu \texttt{icrv\_label}, dùng thống nhất với Chuyên đề 2
(Bảng 2.8).} \\
\textbar{} Nhóm ICRV (icrv\_label) \textbar{} Tên nhóm \textbar{} Quốc
gia \textbar{} N doanh nghiệp \textbar{}
\textbar---\textbar---\textbar---\textbar:--:\textbar{} \textbar{}
\textbf{I} --- \texttt{Advanced\_innovation} \textbar{} Advanced --- đổi
mới sáng tạo \textbar{} Singapore, Hong Kong SAR, Hàn Quốc, Đài Loan,
Israel \textbar{} 4.222 \textbar{} \textbar{} \textbf{II} ---
\texttt{Advanced\_resource} \textbar{} Advanced --- tài nguyên
\textbar{} Saudi Arabia, Qatar, Kuwait, Bahrain, Brunei, Oman \textbar{}
2.269 \textbar{} \textbar{} \textbf{III} --- \texttt{Upper\_mid}
\textbar{} Trung bình cao \textbar{} Trung Quốc, Malaysia, Thái Lan,
Kazakhstan, Armenia, Georgia \textbar{} 13.993 \textbar{} \textbar{}
\textbf{IV} --- \texttt{Lower\_mid\_transition} \textbar{} Chuyển đổi
thu nhập trung bình--thấp \textbar{} Việt Nam, Ấn Độ, Indonesia,
Philippines, Mông Cổ, Bangladesh, Pakistan \textbar{} 50.926 \textbar{}
\textbar{} \textbf{V} --- \texttt{Emerging} \textbar{} Đang nổi
\textbar{} Sri Lanka, Jordan, Lào, Campuchia, Myanmar, Nepal, Bhutan,
Maldives, Uzbekistan, Tajikistan, Kyrgyzstan, Turkmenistan, Afghanistan,
Azerbaijan, Iraq, Lebanon, Yemen (17 nước) \textbar{} 18.569 \textbar{}
\textbar{} \textbf{VI} --- \texttt{SIDS\_small} \textbar{} SIDS (trường
hợp biên) \textbar{} Fiji, Kiribati, Papua New Guinea, Samoa, Solomon
Islands, Tonga, Vanuatu, Timor-Leste \textbar{} 1.885 \textbar{}
\textbar{} \textbf{Tổng} \textbar{} \textbar{} \textbf{49 nền kinh tế}
\textbar{} \textbf{91.864} (≈91.982 mẫu phân tích) \textbar{} \\
\emph{Ghi chú: Cột N theo khung phân tích của bộ dữ liệu ước lượng
(91.864 quan sát phân nhóm; tổng mẫu phân tích ≈91.982). Nhãn nhóm theo
biến phân loại \texttt{icrv\_label}, dùng thống nhất trong luận án và
hai chuyên đề. Pool phân loại đầy đủ = 96.415 DN/52 nền (gồm Comoros,
Turkey, Cyprus ngoài châu Á --- loại khỏi khung 49 nền). SIDS: mẫu chính
của nghiên cứu chuyên sâu SIDS là 7 nền Pacific; kiểm định bền vững mở
rộng 9 nền. Nhật Bản được WBES khảo sát lần đầu năm 2025 (sau thời điểm
chốt khung phân tích) nên không thuộc khung 49 nền; đây là hướng mở rộng
tự nhiên cho nghiên cứu tiếp theo.} \\
\#\#\#\# 2.3.2.4 Đặc điểm các nhóm ICRV \\
\textbf{Bảng 2.3.2.2}. \emph{Đặc điểm điển hình của 6 nhóm ICRV --- WGI,
GNI, cơ chế tăng trưởng, pattern WBES dự kiến.} \\
\textbar{} Nhóm \textbar{} WGI Rule of Law \textbar{} GNI/capita
\textbar{} Cơ chế tăng trưởng \textbar{} FSTS dự kiến \textbar{} sd log
LP (trong đợt) \textbar{} Cơ chế quốc tế hóa → hiệu quả \textbar{}
\textbar---\textbar---\textbar---\textbar---\textbar---\textbar---\textbar---\textbar{}
\textbar{} I --- Innovation \textbar{} \textgreater+0,80 \textbar{}
\textgreater\$30k \textbar{} R\&D + dịch vụ tri thức \textbar{} Trung
bình (6--15\%) \textbar{} \textasciitilde1,04 \textbar{} U ngược nhẹ /
tuyến tính dương; điểm uốn \textasciitilde88,6\% (SG) \textbar{}
\textbar{} II --- Resource \textbar{} +0,00→+0,50 \textbar{}
\textgreater\$20k \textbar{} Xuất khẩu dầu/khí \textbar{} Thấp (1--4\%)
\textbar{} \textasciitilde1,14 \textbar{} Phi tuyến suy yếu; nhà nước tô
phân phối lại \textbar{} \textbar{} III --- Upper-mid \textbar{}
0,00→+0,80 \textbar{} \$5k--\$20k \textbar{} Sản xuất + chuyển đổi
\textbar{} Trung bình (8--12\%) \textbar{} \textasciitilde1,37
\textbar{} U ngược; điểm uốn \textasciitilde47\% \textbar{} \textbar{}
IV --- Lower-mid transition \textbar{} -0,50→0,00 \textbar{} \$1k--\$5k
\textbar{} FDI-driven sản xuất \textbar{} Biến động (7--23\%) \textbar{}
\textasciitilde1,31 \textbar{} U ngược; điểm uốn 39--46\% (VNM) --- bẫy
chuyển hướng nguồn lực \textbar{} \textbar{} V --- Emerging \textbar{}
\textless-0,50 \textbar{} \textless\$1k--\$3k \textbar{} Tự cấp + phi
chính thức \textbar{} Thấp--trung bình \textbar{} \textasciitilde1,39
\textbar{} Yếu; voids thể chế → đổi mới bằng nguồn lực hạn chế
\textbar{} \textbar{} VI --- SIDS \textbar{} Đặc biệt \textbar{}
\$1k--\$8k \textbar{} MIRAB + du lịch \textbar{} Rất thấp (1--12\%)
\textbar{} \textasciitilde1,32 \textbar{} Chi phí buộc phải quốc tế hóa
--- U-curve không vận hành \textbar{} \\
\emph{Nguồn: Tổng hợp từ World Bank Enterprise Surveys (2025); WGI từ
Kaufmann et al.~(2011); GNI từ World Bank (2025, tháng 7). Cột sd log
LP: trung vị độ lệch chuẩn ln(năng suất lao động) tính trong từng đợt
khảo sát (Bảng 2.3.3.1). Pattern quốc tế hóa → hiệu quả từ các nghiên
cứu thành phần của luận án (Đỗ \& Phan, 2026 --- VEFR, JABS, JFAR,
MIR).} \\
\end{longtable}

\hypertarget{thux1ef1c-trux1ea1ng-nux103ng-suux1ea5t-lao-ux111ux1ed9ng-vuxe0-quux1ed1c-tux1ebf-huxf3a}{%
\subsubsection{2.3.3 Thực trạng năng suất lao động và quốc tế
hóa}\label{thux1ef1c-trux1ea1ng-nux103ng-suux1ea5t-lao-ux111ux1ed9ng-vuxe0-quux1ed1c-tux1ebf-huxf3a}}

\hypertarget{nguux1ed3n-dux1eef-liux1ec7u-vuxe0-cux1ea5u-truxfac-pool}{%
\paragraph{2.3.3.1 Nguồn dữ liệu và cấu trúc
pool}\label{nguux1ed3n-dux1eef-liux1ec7u-vuxe0-cux1ea5u-truxfac-pool}}

\textbf{Phạm vi tổng hợp dữ liệu}. Pool mô tả gồm \textbf{86.701 doanh
nghiệp}, 49 nền kinh tế, \textbf{102 cặp nền kinh tế × năm}, giai đoạn
2006--2026 (bao gồm Kiribati 2025 --- đợt khảo sát WBES đầu tiên của nền
kinh tế này). Tổng hợp này kế thừa và mở rộng từ pool 17 nền châu Á mới
nổi (\textasciitilde40.633 doanh nghiệp) của Đỗ \& Phan (2026 --- VEFR),
gấp hơn 2 lần. Phân bố theo nhóm ICRV: Nhóm IV Chuyển đổi TB-thấp 45.003
(51,9\%), Nhóm V Đang nổi 18.957 (21,9\%), Nhóm III Trung bình cao
13.993 (16,1\%), Nhóm I Advanced đổi mới 4.222 (4,9\%), Nhóm VI SIDS
2.295 (2,6\%), Nhóm II Advanced tài nguyên 2.231 (2,6\%). Có \textbf{12
nền kinh tế khảo sát năm 2025} với 14.712 doanh nghiệp (17,0\% pool);
đợt khảo sát Nepal 2025 dùng bảng hỏi BREADY Micro (doanh nghiệp siêu
nhỏ) --- không phải cross-section chuẩn nên không thuộc pool, chỉ được
nhắc đến như tham chiếu.

\emph{Hình 2.3.3.1. Phân bố pool mô tả 86.701 doanh nghiệp (2006--2026)
theo 3 thế hệ schema WBES: PICS3 (2006--2012, 18 đợt, n=8.048),
Standardized (2013--2017, 27 đợt, n=23.940), BREADY (2018--2026, 57 đợt,
n=54.713). Các đợt 2024--2026 đóng góp 25.691 doanh nghiệp (29,6\% pool,
30 đợt) --- phản ánh chu kỳ khảo sát WBES dày lên rõ rệt sau đại dịch.}

\emph{Hình 2.3.3.2. Bar chart phân bố theo 6 nhóm ICRV: IV Chuyển đổi
TB-thấp (45.003, 51,9\%), V Đang nổi (18.957, 21,9\%), III Trung bình
cao (13.993, 16,1\%), I Advanced đổi mới (4.222, 4,9\%), VI SIDS (2.295,
2,6\%), II Advanced tài nguyên (2.231, 2,6\%).}

\textbf{Trường hợp biên}: Nhóm VI đủ \textbf{8/8 nền SIDS} (FJI, PNG,
SLB, TON, VUT, WSM, KIR, TLS). Quy trình hài hòa hóa: FSTS =
\texttt{d3b\ +\ d3c}; winsorize log năng suất ở mức 1/99 phân vị trong
cụm quốc gia × năm.

\hypertarget{phuxe2n-tuxe1n-nux103ng-suux1ea5t-lao-ux111ux1ed9ng}{%
\paragraph{2.3.3.2 Phân tán năng suất lao
động}\label{phuxe2n-tuxe1n-nux103ng-suux1ea5t-lao-ux111ux1ed9ng}}

\textbf{Nguyên tắc đo lường.} Doanh thu WBES được báo cáo bằng
\textbf{nội tệ} của từng nền kinh tế, nên không thể so sánh trực tiếp
\emph{mức} log năng suất giữa các nước. Mọi thống kê phân tán ở Bảng
2.3.3.1 vì vậy được tính \textbf{trong từng cặp nền kinh tế × năm} ---
trên thang log, thừa số quy đổi tiền tệ là hằng số cộng và triệt tiêu
hoàn toàn trong độ lệch chuẩn và các tỷ số phân vị --- rồi lấy trung vị
giữa các đợt trong mỗi nhóm ICRV. Đây là thước đo phân tán hợp lệ và bất
biến với đơn vị tiền tệ.

\textbf{Bảng 2.3.3.1}. \emph{Phân tán năng suất lao động trong từng đợt
khảo sát, theo nhóm ICRV (trung vị giữa các đợt; pool mô tả 102 cặp nền
kinh tế × năm).}

\begin{longtable}[]{@{}lrrrrr@{}}
\toprule\noalign{}
Nhóm ICRV & Số đợt & n\_firms & sd ln(LP) & P90/P10 & P75/P25 \\
\midrule\noalign{}
\endhead
\bottomrule\noalign{}
\endlastfoot
I --- Advanced đổi mới & 6 & 3.620 & \textbf{1,04} & \textbf{12,2} &
3,5 \\
II --- Advanced tài nguyên & 6 & 2.082 & 1,14 & 20,8 & 4,8 \\
III --- Trung bình cao & 17 & 12.481 & 1,37 & 25,0 & 5,4 \\
IV --- Chuyển đổi TB-thấp & 17 & 42.798 & 1,31 & 27,8 & 5,6 \\
V --- Đang nổi & 40 & 16.595 & \textbf{1,39} & \textbf{30,3} & 5,8 \\
VI --- SIDS & 16 & 1.916 & 1,32 & 26,3 & 5,4 \\
\end{longtable}

\emph{Ghi chú: sd, P90/P10 và P75/P25 tính trong từng đợt khảo sát (bất
biến tiền tệ), sau đó lấy trung vị giữa các đợt của nhóm; n\_firms là số
quan sát có năng suất hợp lệ (doanh thu và lao động dương). Trong nội bộ
Nhóm II, các nền lớn nhất có phân tán đặc biệt hẹp (Saudi Arabia 0,49;
Qatar 0,33) --- đặc trưng nhà nước tô; trung vị nhóm bị kéo lên bởi các
nền nhỏ (Kuwait 1,16; Brunei 1,12). Kiribati 2025: sd 1,48 --- rộng nhất
trong các đợt SIDS.}

\emph{Hình 2.3.3.3. Đồ thị mật độ kernel của năng suất chuẩn hóa trong
đợt cho 6 nhóm ICRV --- Nhóm I hẹp nhất (P90/P10 ≈ 12 lần), Nhóm V rộng
nhất (≈ 30 lần).}

\textbf{Ba phát hiện từ phân tán trong đợt}: (1) \textbf{Gradient thể
chế của phân tán là thực và đơn điệu theo hướng yếu dần của thể chế}:
doanh nghiệp ở phân vị 90 năng suất gấp doanh nghiệp ở phân vị 10 khoảng
\textbf{12 lần} tại Nhóm I, nhưng \textbf{25--30 lần} tại Nhóm III--V
--- nhất quán với giả thuyết phân bổ sai nguồn lực (Hsieh \& Klenow,
2009, 2014), với lưu ý sd log năng suất lao động là proxy gián tiếp cho
phân tán TFPR nên đây là bằng chứng mô tả phù hợp, không phải phân rã
trực tiếp. (2) \textbf{So sánh mức giữa các nền kinh tế} dùng chỉ số
không thứ nguyên \textbf{ROS} (tỷ suất sinh lời trên doanh thu, tính
trên khung ước lượng của luận án): ROS trung vị giảm từ \textbf{0,50
(Nhóm I) → 0,45 (III) → 0,35 (IV)}; Nhóm V--VI có outlier mạnh nên chỉ
đọc trung vị (V≈0,41; VI≈0,47) --- thể chế yếu đi kèm biên lợi nhuận
mỏng hơn, nhất quán với Lý thuyết thể chế (chi phí giao dịch cao bào mòn
biên). (3) Suy diễn về \textbf{điều tiết thể chế} trong quan hệ quốc tế
hóa → hiệu quả được neo vào bằng chứng suy diễn ở meta-analysis của luận
án (\(Q_M=17{,}35\)) và gradient điểm uốn xuyên nhóm ICRV của nghiên cứu
đa quốc gia (P7), chứ không chỉ vào thống kê mô tả.

\hypertarget{thux1ef1c-trux1ea1ng-quux1ed1c-tux1ebf-huxf3a-vuxe0-tux103ng-trux1b0ux1edfng-viux1ec7c-luxe0m}{%
\paragraph{2.3.3.3 Thực trạng quốc tế hóa và tăng trưởng việc
làm}\label{thux1ef1c-trux1ea1ng-quux1ed1c-tux1ebf-huxf3a-vuxe0-tux103ng-trux1b0ux1edfng-viux1ec7c-luxe0m}}

\textbf{Bảng 2.3.3.2}. \emph{FSTS, doanh nghiệp xuất khẩu và CAGR việc
làm theo phân nhóm con.}

\begin{longtable}[]{@{}
  >{\raggedright\arraybackslash}p{(\columnwidth - 6\tabcolsep) * \real{0.2500}}
  >{\raggedright\arraybackslash}p{(\columnwidth - 6\tabcolsep) * \real{0.2500}}
  >{\raggedright\arraybackslash}p{(\columnwidth - 6\tabcolsep) * \real{0.2500}}
  >{\raggedright\arraybackslash}p{(\columnwidth - 6\tabcolsep) * \real{0.2500}}@{}}
\toprule\noalign{}
\begin{minipage}[b]{\linewidth}\raggedright
Nhóm ICRV (icrv\_label)
\end{minipage} & \begin{minipage}[b]{\linewidth}\raggedright
FSTS (\%)
\end{minipage} & \begin{minipage}[b]{\linewidth}\raggedright
Doanh nghiệp xuất khẩu (\%)
\end{minipage} & \begin{minipage}[b]{\linewidth}\raggedright
CAGR việc làm (\%)
\end{minipage} \\
\midrule\noalign{}
\endhead
\bottomrule\noalign{}
\endlastfoot
I --- Adv. đổi mới & 13,4 & 28,4 & 3,0 \\
II --- Adv. tài nguyên¹ & 3,9 & 11,7 & 3,7 \\
III --- Trung bình cao & 10,2 & 21,3 & 3,8 \\
IV --- Chuyển đổi TB-thấp & 8,1 & 14,3 & 2,8 \\
V --- Đang nổi & 10,6 & 18,8 & 4,5 \\
VI --- SIDS² & 6,7 & 14,7 & 5,3 \\
\end{longtable}

\emph{Nguồn: tính từ dữ liệu vi mô WBES trên pool mô tả 49 nền (FSTS =
d3b + d3c; CAGR việc làm từ l1/l2 trên 3 năm tài khóa). ¹Nhóm II: đủ 6/6
nền GCC (Bahrain, Brunei, Kuwait, Oman, Qatar, Saudi Arabia). ²Nhóm VI:
đủ 8/8 nền (7 Pacific + Timor-Leste; n=2.295).}

Trung vị FSTS bằng 0\% --- phân phối phân cực mạnh với chỉ khoảng
12--28\% doanh nghiệp tham gia xuất khẩu tùy theo nhóm (cao nhất ở Nhóm
I, thấp nhất ở Nhóm II). SIDS có tốc độ tăng trưởng việc làm hàng năm
kép cao nhất (5,3\%), phản ánh quy mô nhỏ và tăng trưởng từ nền thấp.
Nhóm I và V có FSTS trung bình cao nhất (13,4\% và 10,6\%) nhưng cơ chế
xuất khẩu khác nhau: Nhóm I chủ yếu dịch vụ và hàng hóa hàm lượng tri
thức cao; các nhóm III--V chủ yếu sản xuất theo chuỗi giá trị toàn cầu.

\begin{center}\rule{0.5\linewidth}{0.5pt}\end{center}

\hypertarget{ux111ux1ed5i-mux1edbi-suxe1ng-tux1ea1o-vuxe0-nux103ng-lux1ef1c-sux1ed1-huxf3a}{%
\subsubsection{2.3.4 Đổi mới sáng tạo và năng lực số
hóa}\label{ux111ux1ed5i-mux1edbi-suxe1ng-tux1ea1o-vuxe0-nux103ng-lux1ef1c-sux1ed1-huxf3a}}

\hypertarget{khung-5-lux129nh-vux1ef1c-chux1ec9-sux1ed1-wbes}{%
\paragraph{2.3.4.1 Khung 5 lĩnh vực chỉ số
WBES}\label{khung-5-lux129nh-vux1ef1c-chux1ec9-sux1ed1-wbes}}

WBES tổ chức các chỉ số doanh nghiệp thành \textbf{5 lĩnh vực chính
thức}: (1) Quy định pháp lý và thuế; (2) Tài chính và tín dụng; (3) Cơ
sở hạ tầng và khí hậu; (4) Thương mại và cạnh tranh; (5) Tham nhũng và
phi chính thức.

\textbf{Bảng 2.3.4.2}. \emph{Khung 5 lĩnh vực chỉ số WBES --- định vị
phạm vi đo lường.}

\begin{longtable}[]{@{}
  >{\raggedright\arraybackslash}p{(\columnwidth - 8\tabcolsep) * \real{0.2000}}
  >{\raggedright\arraybackslash}p{(\columnwidth - 8\tabcolsep) * \real{0.2000}}
  >{\raggedright\arraybackslash}p{(\columnwidth - 8\tabcolsep) * \real{0.2000}}
  >{\raggedright\arraybackslash}p{(\columnwidth - 8\tabcolsep) * \real{0.2000}}
  >{\raggedright\arraybackslash}p{(\columnwidth - 8\tabcolsep) * \real{0.2000}}@{}}
\toprule\noalign{}
\begin{minipage}[b]{\linewidth}\raggedright
\#
\end{minipage} & \begin{minipage}[b]{\linewidth}\raggedright
Lĩnh vực
\end{minipage} & \begin{minipage}[b]{\linewidth}\raggedright
Ví dụ chỉ số WBES
\end{minipage} & \begin{minipage}[b]{\linewidth}\raggedright
CĐ1 đo trực tiếp
\end{minipage} & \begin{minipage}[b]{\linewidth}\raggedright
Hàm ý CĐ2
\end{minipage} \\
\midrule\noalign{}
\endhead
\bottomrule\noalign{}
\endlastfoot
(1) & Quy định pháp lý + Thuế & Thuế thời gian (giờ/năm); Cấp phép &
(gợi ý CĐ2) & \texttt{time\_tax\_hours},
\texttt{tax\_inspection\_freq} \\
(2) & Tài chính + Tín dụng & Hạn chế tín dụng; Số hóa tài chính & (gợi ý
--- §2.3.5.2 rào cản \#1) & \texttt{credit\_constraint\_dummy} \\
(3) & Cơ sở hạ tầng + Khí hậu & Quản lý năng lượng; Mất điện & §2.3.5.2
rào cản \#4 & \texttt{power\_outage\_intensity} từ c30 \\
(4) & Thương mại + Cạnh tranh & Thời gian thông quan; FSTS & FSTS đã đo
§2.3.3 & \texttt{customs\_days} + \texttt{informal\_competition} \\
(5) & Tham nhũng + Phi chính thức & Hối lộ (Bribery); Tham ô (Graft) &
§2.3.6.2 & \texttt{bribery\_incidence\_pct} \\
\end{longtable}

\hypertarget{thux1ef1c-trux1ea1ng-ux111ux1ed5i-mux1edbi-suxe1ng-tux1ea1o-vuxe0-uxe1p-dux1ee5ng-sux1ed1}{%
\paragraph{2.3.4.2 Thực trạng đổi mới sáng tạo và áp dụng
số}\label{thux1ef1c-trux1ea1ng-ux111ux1ed5i-mux1edbi-suxe1ng-tux1ea1o-vuxe0-uxe1p-dux1ee5ng-sux1ed1}}

\textbf{Bảng 2.3.4.1}. \emph{Đổi mới sáng tạo và áp dụng số (\%).}

\begin{longtable}[]{@{}
  >{\raggedright\arraybackslash}p{(\columnwidth - 10\tabcolsep) * \real{0.1667}}
  >{\raggedright\arraybackslash}p{(\columnwidth - 10\tabcolsep) * \real{0.1667}}
  >{\raggedright\arraybackslash}p{(\columnwidth - 10\tabcolsep) * \real{0.1667}}
  >{\raggedright\arraybackslash}p{(\columnwidth - 10\tabcolsep) * \real{0.1667}}
  >{\raggedright\arraybackslash}p{(\columnwidth - 10\tabcolsep) * \real{0.1667}}
  >{\raggedright\arraybackslash}p{(\columnwidth - 10\tabcolsep) * \real{0.1667}}@{}}
\toprule\noalign{}
\begin{minipage}[b]{\linewidth}\raggedright
Nhóm ICRV (icrv\_label)
\end{minipage} & \begin{minipage}[b]{\linewidth}\raggedright
Sản phẩm mới
\end{minipage} & \begin{minipage}[b]{\linewidth}\raggedright
Quy trình mới
\end{minipage} & \begin{minipage}[b]{\linewidth}\raggedright
R\&D
\end{minipage} & \begin{minipage}[b]{\linewidth}\raggedright
ISO
\end{minipage} & \begin{minipage}[b]{\linewidth}\raggedright
Website
\end{minipage} \\
\midrule\noalign{}
\endhead
\bottomrule\noalign{}
\endlastfoot
I --- Adv. đổi mới & 24,3 & 14,7 & \textbf{20,5} & \textbf{35,0} &
\textbf{61,7} \\
II --- Adv. tài nguyên¹ & 17,0 & 8,6 & 7,0 & 15,2 & 50,2 \\
III --- Trung bình cao & 26,7 & 18,9 & 21,5 & 25,4 & 53,9 \\
IV --- Chuyển đổi TB-thấp & 17,1 & 14,7 & 15,5 & 25,3 & 48,3 \\
V --- Đang nổi & 24,4 & 16,3 & 17,5 & 12,7 & 40,9 \\
VI --- SIDS² & \textbf{34,2} & \textbf{24,7} & 10,2 & 12,2 & 41,5 \\
\end{longtable}

\emph{Nguồn: tính từ dữ liệu vi mô WBES trên pool mô tả 49 nền (biến
h1/h5/h8/b8/c22b; mã 1=có, 2=không; mã âm = thiếu). ¹Nhóm II: đủ 6/6 nền
GCC --- đặc trưng dầu mỏ: R\&D thấp (7,0\%), nữ quản lý cấp cao rất thấp
(4,6\%), website 50,2\%. ²Nhóm VI: đủ 8/8 nền --- đổi mới sản phẩm cao
nhất mọi nhóm (34,2\%) và website 41,5\%: pattern ``thích nghi và nhảy
vọt số'' của SIDS.}

Với dữ liệu đầy đủ 8/8 nền Nhóm VI, SIDS thể hiện pattern \textbf{thích
nghi và nhảy vọt số} rõ nét: tỷ lệ đổi mới sản phẩm \textbf{cao nhất mọi
nhóm (34,2\%)} phản ánh đổi mới bằng nguồn lực hạn chế --- sáng tạo
không chính thức bù đắp thiếu hụt nguồn lực; website 41,5\% ngang nhóm
Đang nổi dù GNI thấp hơn đáng kể. Gradient năng lực thể chế vẫn rõ ở các
biến nền tảng: R\&D giảm từ 20,5\% (Nhóm I) xuống 10,2\% (Nhóm VI); ISO
35,0\%→12,2\%; website 61,7\%→41,5\%. Phân tách rõ TCI so với DAI --- kế
thừa Đỗ \& Phan (2026 --- VEFR).

\hypertarget{tuxe1i-ux111ux1ecbnh-huxecnh-chux1ec9-sux1ed1-uxe1p-dux1ee5ng-sux1ed1-tier-1-sux1ef1-hiux1ec7n-diux1ec7n-sux1ed1-cux1a1-bux1ea3n}{%
\paragraph{2.3.4.3 Tái định hình Chỉ số áp dụng số --- Tier-1 Sự hiện
diện số cơ
bản}\label{tuxe1i-ux111ux1ecbnh-huxecnh-chux1ec9-sux1ed1-uxe1p-dux1ee5ng-sux1ed1-tier-1-sux1ef1-hiux1ec7n-diux1ec7n-sux1ed1-cux1a1-bux1ea3n}}

Spec 1 (toàn bộ 2006--2026) dùng biến nhị phân duy nhất (website
có/không). \textbf{Hệ quả phân tích đáng chú ý}: phần bù tương quan DAI
ở Nhóm I là thấp nhất trong sáu nhóm (+0,098 --- Bảng 2.3.8.1), trong
khi tỷ lệ có website ở nhóm này đã đạt 61,7\%. Khi một chỉ báo nhị phân
tiệm cận bão hòa, nó mất dần khả năng phân biệt năng suất: phương sai
của biến co lại và phần doanh nghiệp ``không website'' còn lại ngày càng
mang tính chọn lọc đặc thù (quy mô siêu nhỏ, ngành không cần hiện diện
web) --- nên hệ số đo được phản ánh \textbf{giới hạn của thước đo
Tier-1} nhiều hơn là giới hạn của số hóa.

\textbf{Kiểm định loại trừ ICT (đề xuất cho Chuyên đề 2)}: so sánh hệ số
DAI của phân nhóm Advanced trước và sau khi loại các doanh nghiệp ICT
(ISIC J 58--63 --- nơi website là điều kiện tồn tại, không phải lựa chọn
đầu tư) sẽ tách được phần ``bão hòa thước đo'' khỏi phần bù số thực.
Theo logic đó, biến Spec 1 được đặt tên chính xác là \textbf{``Sự hiện
diện số cơ bản (Tier-1)''} --- phạm vi khiêm tốn và phản ánh đúng thực
tế đo lường.

\begin{center}\rule{0.5\linewidth}{0.5pt}\end{center}

\hypertarget{cux1ea5u-truxfac-doanh-nghiux1ec7p}{%
\subsubsection{2.3.5 Cấu trúc doanh
nghiệp}\label{cux1ea5u-truxfac-doanh-nghiux1ec7p}}

\hypertarget{phuxe2n-bux1ed1-quy-muxf4-vuxe0-sux1edf-hux1eefu}{%
\paragraph{2.3.5.1 Phân bố quy mô và sở
hữu}\label{phuxe2n-bux1ed1-quy-muxf4-vuxe0-sux1edf-hux1eefu}}

\textbf{Bảng 2.3.5.1}. \emph{Cấu trúc doanh nghiệp theo phân nhóm con
(\%).}

\begin{longtable}[]{@{}llll@{}}
\toprule\noalign{}
Nhóm ICRV (icrv\_label) & SME & Doanh nghiệp xuất khẩu & FDI ≥10\% \\
\midrule\noalign{}
\endhead
\bottomrule\noalign{}
\endlastfoot
I --- Adv. đổi mới & 79,5 & 28,4 & 10,6 \\
II --- Adv. tài nguyên¹ & 77,4 & 11,7 & 13,6 \\
III --- Trung bình cao & 78,7 & 21,3 & 8,9 \\
IV --- Chuyển đổi TB-thấp & 74,0 & 14,3 & 3,1 \\
V --- Đang nổi & 86,0 & 18,8 & 7,1 \\
VI --- SIDS² & \textbf{91,8} & 14,7 & \textbf{20,0} \\
\end{longtable}

\emph{Nguồn: tính từ dữ liệu vi mô WBES trên pool mô tả 49 nền (SME =
\textless100 lao động thường xuyên, biến l1; FDI = sở hữu nước ngoài
≥10\%, biến b2b). ¹Nhóm II: đủ 6/6 nền GCC. ²Nhóm VI: đủ 8/8 nền --- SME
cao nhất (91,8\%) và FDI cao (20,0\%, do MNE du lịch/viễn thông).}

\emph{Hình 2.3.5.1. Stacked bar chart: Nhóm VI (SIDS, đủ 8/8 nền) có tỷ
lệ SME cao nhất (91,8\%), Nhóm V kế tiếp (85,7\%), Nhóm IV thấp nhất
(74,0\%).}

Tỷ lệ FDI có dạng chữ U với cực tiểu ở Nhóm IV --- Chuyển đổi TB-thấp
(3,1\%). SIDS có FDI cao (20,0\%) do du lịch và viễn thông được dẫn dắt
bởi doanh nghiệp đa quốc gia. Hai mô hình FDI hoàn toàn khác nhau: (i)
MNE hub tại Advanced (Singapore FDI 31,5\% --- doanh nghiệp đa quốc gia
chọn Singapore làm trung tâm khu vực); (ii) du lịch/viễn thông tại SIDS
(FDI nhắm vào khu vực không cạnh tranh với lao động địa phương).

\hypertarget{bux1ed1n-ruxe0o-cux1ea3n-huxe0ng-ux111ux1ea7u-cux1ee7a-khu-vux1ef1c-tux1b0-nhuxe2n-uxe1-thuxe1i}{%
\paragraph{2.3.5.2 Bốn rào cản hàng đầu của khu vực tư nhân
Á-Thái}\label{bux1ed1n-ruxe0o-cux1ea3n-huxe0ng-ux111ux1ea7u-cux1ee7a-khu-vux1ef1c-tux1b0-nhuxe2n-uxe1-thuxe1i}}

\textbf{Bảng 2.3.5.2}. \emph{Bốn rào cản hàng đầu Asia-Pacific WBES ---
phân tầng theo phân nhóm con thể chế.}

\begin{longtable}[]{@{}
  >{\raggedright\arraybackslash}p{(\columnwidth - 12\tabcolsep) * \real{0.1429}}
  >{\raggedright\arraybackslash}p{(\columnwidth - 12\tabcolsep) * \real{0.1429}}
  >{\raggedright\arraybackslash}p{(\columnwidth - 12\tabcolsep) * \real{0.1429}}
  >{\raggedright\arraybackslash}p{(\columnwidth - 12\tabcolsep) * \real{0.1429}}
  >{\raggedright\arraybackslash}p{(\columnwidth - 12\tabcolsep) * \real{0.1429}}
  >{\raggedright\arraybackslash}p{(\columnwidth - 12\tabcolsep) * \real{0.1429}}
  >{\raggedright\arraybackslash}p{(\columnwidth - 12\tabcolsep) * \real{0.1429}}@{}}
\toprule\noalign{}
\begin{minipage}[b]{\linewidth}\raggedright
\#
\end{minipage} & \begin{minipage}[b]{\linewidth}\raggedright
Rào cản
\end{minipage} & \begin{minipage}[b]{\linewidth}\raggedright
Tên chuẩn WB
\end{minipage} & \begin{minipage}[b]{\linewidth}\raggedright
Cường độ Nhóm V--VI
\end{minipage} & \begin{minipage}[b]{\linewidth}\raggedright
Cường độ Nhóm III--IV
\end{minipage} & \begin{minipage}[b]{\linewidth}\raggedright
Cường độ Nhóm I--II
\end{minipage} & \begin{minipage}[b]{\linewidth}\raggedright
Lý thuyết liên kết
\end{minipage} \\
\midrule\noalign{}
\endhead
\bottomrule\noalign{}
\endlastfoot
(1) & \textbf{Tiếp cận tài chính} & Credit constraint & Cao & Cao & Thấp
& Khanna \& Palepu (2010) --- voids thể chế \\
(2) & \textbf{Lực lượng lao động thiếu kỹ năng} & Skilled labor shortage
& Cao & Trung bình & Thấp & Cohen \& Levinthal (1990) --- năng lực hấp
thụ \\
(3) & \textbf{Cạnh tranh từ doanh nghiệp phi chính thức} & Informal
sector competition & Cao & Cao & Thấp & La Porta \& Shleifer (2008,
2014) kinh tế ngầm \\
(4) & \textbf{Nguồn cung điện không đáng tin cậy} & Power outages & Cao
(đặc biệt SIDS) & Trung bình & Thấp & IFC PSD Blueprint trụ cột nguồn
lực vận hành \\
\end{longtable}

Bốn rào cản này tạo ra bốn biến kiểm soát quan trọng cho Chuyên đề 2:
\texttt{credit\_constraint\_dummy}, \texttt{labor\_skill\_gap},
\texttt{informal\_competition}, \texttt{power\_outage\_intensity}. Sự
hiện diện đồng thời của nhiều voids thể chế tại Nhóm V giải thích tại
sao quan hệ quốc tế hóa → hiệu quả ở nhóm này yếu và không đơn điệu.

\begin{center}\rule{0.5\linewidth}{0.5pt}\end{center}

\hypertarget{biux1ebfn-thiuxean-theo-thux1eddi-gian-20062026}{%
\subsubsection{2.3.6 Biến thiên theo thời gian
(2006--2026)}\label{biux1ebfn-thiuxean-theo-thux1eddi-gian-20062026}}

\hypertarget{xu-hux1b0ux1edbng-thay-ux111ux1ed5i-chuxednh}{%
\paragraph{2.3.6.1 Xu hướng thay đổi
chính}\label{xu-hux1b0ux1edbng-thay-ux111ux1ed5i-chuxednh}}

\textbf{Bảng 2.3.6.1}. \emph{Δ điểm phần trăm khi so sánh các đợt
2018--2026 với các đợt 2006--2012 (gộp doanh nghiệp trong từng kỳ).}

\begin{longtable}[]{@{}lrrrrr@{}}
\toprule\noalign{}
Nhóm ICRV & Δ Website & Δ DN xuất khẩu & Δ FDI ≥10\% & Δ R\&D & Δ ISO \\
\midrule\noalign{}
\endhead
\bottomrule\noalign{}
\endlastfoot
III --- Trung bình cao & +11,2 & +12,5 & +1,6 & ---¹ & +5,2 \\
IV --- Chuyển đổi TB-thấp & \textbf{+20,6} & \textbf{−11,3} & −6,7 &
---¹ & +3,4 \\
V --- Đang nổi & \textbf{+21,6} & +1,5 & −2,5 & −31,6² & 0,0 \\
VI --- SIDS & \textbf{+30,8} & −2,6 & −5,3 & ---¹ & −10,2 \\
\end{longtable}

\emph{¹ Δ R\&D chỉ tính được cho Nhóm V: biến chi R\&D (h8) không có
trong các đợt 2006--2012 của Nhóm III, IV và VI. Nhóm I và II không có
đợt khảo sát nào trước 2013 trong pool nên không so sánh dọc được. ² Mức
giảm R\&D ở Nhóm V chủ yếu là hiệu ứng schema (BREADY đo R\&D khác
PICS3), không phản ánh suy giảm đầu tư thực.}

\textbf{Nhảy vọt số} là xu hướng nổi bật nhất: website tăng +21 đến +31
điểm phần trăm ở Nhóm IV, V và VI --- phù hợp luận điểm nhảy vọt số của
Banalieva \& Dhanaraj (2019). Ngược chiều, tỷ lệ doanh nghiệp xuất khẩu
của Nhóm IV giảm 11,3 điểm phần trăm --- chủ yếu do đứt gãy schema tại
Ấn Độ 2025 (FSTS đo lường thấp hơn hệ thống) cộng hưởng với cấu trúc FDI
lắp ráp; trong khi Nhóm III tăng 12,5 điểm phần trăm theo chuỗi giá trị
toàn cầu. FDI giảm nhẹ phổ biến ở các nhóm IV--VI.

\emph{Hình 2.3.4.1. Slope chart minh họa thay đổi 5 chỉ số giữa 2 kỳ
khảo sát cho 4 nhóm ICRV có dữ liệu hai kỳ (III--VI). Pattern nổi bật:
website tăng mạnh ở IV/V/VI; FDI giảm nhẹ phổ biến; R\&D Nhóm V giảm
mạnh do hiệu ứng schema.}

\emph{Hình 2.3.6.1. Radar/spider chart cho 3 nhóm (III, V, VI) trên 5
chiều (FSTS, Đổi mới sản phẩm, R\&D, Website, ISO) với 2 kỳ khảo sát
2006--2012 vs 2018--2026.}

\hypertarget{phuxe2n-biux1ec7t-3-chux1ec9-sux1ed1-tham-nhux169ng-wbes}{%
\paragraph{2.3.6.2 Phân biệt 3 chỉ số tham nhũng
WBES}\label{phuxe2n-biux1ec7t-3-chux1ec9-sux1ed1-tham-nhux169ng-wbes}}

\textbf{Bảng 2.3.6.2}. \emph{Phân biệt 3 chỉ số tham nhũng WBES --- định
nghĩa chính thức và cách đo.}

\begin{longtable}[]{@{}
  >{\raggedright\arraybackslash}p{(\columnwidth - 8\tabcolsep) * \real{0.2000}}
  >{\raggedright\arraybackslash}p{(\columnwidth - 8\tabcolsep) * \real{0.2000}}
  >{\raggedright\arraybackslash}p{(\columnwidth - 8\tabcolsep) * \real{0.2000}}
  >{\raggedright\arraybackslash}p{(\columnwidth - 8\tabcolsep) * \real{0.2000}}
  >{\raggedright\arraybackslash}p{(\columnwidth - 8\tabcolsep) * \real{0.2000}}@{}}
\toprule\noalign{}
\begin{minipage}[b]{\linewidth}\raggedright
\#
\end{minipage} & \begin{minipage}[b]{\linewidth}\raggedright
Chỉ số
\end{minipage} & \begin{minipage}[b]{\linewidth}\raggedright
Tiếng Việt chuẩn WB
\end{minipage} & \begin{minipage}[b]{\linewidth}\raggedright
Định nghĩa
\end{minipage} & \begin{minipage}[b]{\linewidth}\raggedright
Ngưỡng cảnh báo
\end{minipage} \\
\midrule\noalign{}
\endhead
\bottomrule\noalign{}
\endlastfoot
(a-1) & \textbf{Bribery Incidence} & \textbf{Tỷ lệ hối lộ} & \% doanh
nghiệp gặp ít nhất 1 yêu cầu hối lộ trong 6 giao dịch công &
\textgreater20\% = báo động đỏ \\
(a-2) & \textbf{Bribery Depth} & \textbf{Độ sâu hối lộ} & \% giao dịch
công có hối lộ trên tổng số giao dịch công của doanh nghiệp &
\textgreater15\% = báo động đỏ \\
(b-1) & \textbf{Graft Index} & \textbf{Chỉ số tham ô} & Tỷ lệ hối lộ
trong 6 giao dịch nhưng \textbf{loại trừ thanh tra thuế} &
\textgreater10\% = báo động đỏ \\
\end{longtable}

Lý do loại trừ thanh tra thuế khỏi Graft Index: doanh nghiệp khai báo
``phải hối lộ thanh tra thuế'' có lệch chọn lọc vì những doanh nghiệp
trốn thuế dễ thừa nhận hơn. Loại trừ thanh tra thuế giúp Graft Index
phản ánh \textbf{tham ô hệ thống} thay vì hành vi cá biệt.

\hypertarget{khung-phuxe2n-tuxedch-cux1ea5p-nguxe0nh}{%
\paragraph{2.3.6.3 Khung phân tích cấp
ngành}\label{khung-phuxe2n-tuxedch-cux1ea5p-nguxe0nh}}

\textbf{Bảng 2.3.6.3}. \emph{Khung phân loại 9 ngành ISIC Rev.~4, bao
gồm cột Dynamism profile per Kafouros et al.~(2023).}

\begin{longtable}[]{@{}
  >{\raggedright\arraybackslash}p{(\columnwidth - 8\tabcolsep) * \real{0.2000}}
  >{\raggedright\arraybackslash}p{(\columnwidth - 8\tabcolsep) * \real{0.2000}}
  >{\raggedright\arraybackslash}p{(\columnwidth - 8\tabcolsep) * \real{0.2000}}
  >{\raggedright\arraybackslash}p{(\columnwidth - 8\tabcolsep) * \real{0.2000}}
  >{\raggedright\arraybackslash}p{(\columnwidth - 8\tabcolsep) * \real{0.2000}}@{}}
\toprule\noalign{}
\begin{minipage}[b]{\linewidth}\raggedright
Ngành
\end{minipage} & \begin{minipage}[b]{\linewidth}\raggedright
Mã ISIC
\end{minipage} & \begin{minipage}[b]{\linewidth}\raggedright
Đặc điểm
\end{minipage} & \begin{minipage}[b]{\linewidth}\raggedright
Pattern dự kiến
\end{minipage} & \begin{minipage}[b]{\linewidth}\raggedright
Dynamism profile
\end{minipage} \\
\midrule\noalign{}
\endhead
\bottomrule\noalign{}
\endlastfoot
Sản xuất --- high-tech & C (20--21, 26--30) & Thâm dụng vốn, định hướng
xuất khẩu & FSTS cao; R\&D cao; ISO cao & \textbf{Dynamism công nghệ
cao} \\
Sản xuất --- low-tech & C (10--18, 22--25, 31--33) & Thâm dụng lao động
& FSTS trung bình; R\&D thấp & \textbf{Dynamism công nghệ thấp} \\
Bán buôn/bán lẻ & G (45--47) & Thâm dụng lao động & FSTS thấp; website
cao & \textbf{Dynamism thị trường cao} \\
Du lịch/khách sạn & I (55--56) & Dịch vụ phụ thuộc du lịch & FDI cao ở
SIDS & \textbf{Dynamism thị trường cao} \\
Vận tải & H (49--53) & Kinh tế mạng lưới & FDI cao; ISO cao & Trung
gian \\
ICT & J (58--63) & Thâm dụng tri thức & R\&D cao; website
\textasciitilde100\% & \textbf{Dynamism công nghệ cao} \\
Xây dựng & F (41--43) & Theo dự án & FSTS thấp; FDI ở GCC &
\textbf{Dynamism thấp} \\
Khai khoáng & B (05--09) & Tài nguyên dẫn dắt & FDI cao; ISO ở MNE &
\textbf{Dynamism công nghệ thấp} \\
Tài chính & K (64--66) & Có quy chế, xuyên biên giới & FDI cao; website
\textasciitilde100\% & \textbf{Dynamism thị trường cao} \\
\end{longtable}

Kafouros et al.~(2023) chứng minh chất lượng thể chế tương tác hai chiều
với dynamism ngành: ở ngành dynamism công nghệ cao, thể chế tốt khuếch
đại hiệu quả; ở ngành dynamism thị trường cao, thể chế tốt có thể làm
yếu hiệu quả (doanh nghiệp phải nội bộ hóa chức năng). Phát hiện này gợi
ý hướng kiểm định cho Chuyên đề 2.

\hypertarget{dux1ecb-biux1ec7t-nux1ed9i-bux1ed9-khux1ed1i-chuxe2u-uxe1-mux1edbi-nux1ed5i-7-nux1ec1n-thuux1ed9c-nhuxf3m-ivv}{%
\paragraph{2.3.6.4 Dị biệt nội bộ khối châu Á mới nổi (7 nền thuộc Nhóm
IV--V)}\label{dux1ecb-biux1ec7t-nux1ed9i-bux1ed9-khux1ed1i-chuxe2u-uxe1-mux1edbi-nux1ed5i-7-nux1ec1n-thuux1ed9c-nhuxf3m-ivv}}

Bảy nền kinh tế mới nổi đông dân nhất của khung phân tích (Việt Nam,
Indonesia, Philippines, Ấn Độ, Mông Cổ thuộc Nhóm IV; Sri Lanka, Jordan
thuộc Nhóm V) thường bị văn liệu gộp thành một khối ``Emerging Asia''
đồng nhất. Dữ liệu cho thấy ba phân nhóm con khác biệt rõ:

\textbf{Bảng 2.3.6.4}. \emph{Dị biệt nội bộ khối châu Á mới nổi --- số
liệu tổng hợp pool mô tả 2006--2026 (n=43.502).}

\begin{longtable}[]{@{}
  >{\raggedright\arraybackslash}p{(\columnwidth - 18\tabcolsep) * \real{0.1000}}
  >{\raggedright\arraybackslash}p{(\columnwidth - 18\tabcolsep) * \real{0.1000}}
  >{\raggedleft\arraybackslash}p{(\columnwidth - 18\tabcolsep) * \real{0.1000}}
  >{\raggedleft\arraybackslash}p{(\columnwidth - 18\tabcolsep) * \real{0.1000}}
  >{\raggedleft\arraybackslash}p{(\columnwidth - 18\tabcolsep) * \real{0.1000}}
  >{\raggedleft\arraybackslash}p{(\columnwidth - 18\tabcolsep) * \real{0.1000}}
  >{\raggedleft\arraybackslash}p{(\columnwidth - 18\tabcolsep) * \real{0.1000}}
  >{\raggedleft\arraybackslash}p{(\columnwidth - 18\tabcolsep) * \real{0.1000}}
  >{\raggedleft\arraybackslash}p{(\columnwidth - 18\tabcolsep) * \real{0.1000}}
  >{\raggedleft\arraybackslash}p{(\columnwidth - 18\tabcolsep) * \real{0.1000}}@{}}
\toprule\noalign{}
\begin{minipage}[b]{\linewidth}\raggedright
Phân nhóm con
\end{minipage} & \begin{minipage}[b]{\linewidth}\raggedright
Nền kinh tế
\end{minipage} & \begin{minipage}[b]{\linewidth}\raggedleft
n\_firms
\end{minipage} & \begin{minipage}[b]{\linewidth}\raggedleft
FSTS (\%)
\end{minipage} & \begin{minipage}[b]{\linewidth}\raggedleft
DN xuất khẩu (\%)
\end{minipage} & \begin{minipage}[b]{\linewidth}\raggedleft
FDI ≥10\% (\%)
\end{minipage} & \begin{minipage}[b]{\linewidth}\raggedleft
Website (\%)
\end{minipage} & \begin{minipage}[b]{\linewidth}\raggedleft
R\&D (\%)
\end{minipage} & \begin{minipage}[b]{\linewidth}\raggedleft
ISO (\%)
\end{minipage} & \begin{minipage}[b]{\linewidth}\raggedleft
sd ln(LP)¹
\end{minipage} \\
\midrule\noalign{}
\endhead
\bottomrule\noalign{}
\endlastfoot
FDI dẫn dắt SEA & VNM, IDN, PHL & 9.798 & \textbf{12,4} & 20,5 & 8,6 &
43,0 & 6,8 & 16,0 & 1,47 \\
Dân số lớn & IND, LKA, JOR & 32.622 & \textbf{7,1} & 13,8 & 2,0 & 51,5 &
17,9 & 28,0 & 1,25 \\
Tài nguyên & MNG & 1.082 & \textbf{4,7} & 9,0 & 5,7 & 43,6 & 33,7 & 14,0
& 1,27 \\
\textbf{Tổng 7 nền} & & \textbf{43.502} & 8,2 & 15,1 & 3,5 & 49,4 & 16,5
& 25,1 & 1,31 \\
\end{longtable}

\emph{¹ Trung vị độ lệch chuẩn ln(năng suất lao động) tính trong từng
đợt khảo sát (bất biến tiền tệ).}

Năm phát hiện: (1) FSTS phân tầng 4,7\%--7,1\%--12,4\% bị che giấu khi
gộp; (2) FDI chênh lệch hơn 4 lần giữa ASEAN-3 và khối dân số lớn; (3)
cấu trúc năng lực ngược chiều --- khối FDI dẫn dắt có FSTS cao nhất
nhưng R\&D thấp nhất (6,8\%), khối dân số lớn ngược lại (R\&D 17,9\%,
ISO 28,0\%); (4) Mông Cổ là trường hợp biên với đặc điểm tài nguyên
(R\&D 33,7\% nhưng FSTS chỉ 4,7\%); (5) phân tán trong đợt 1,47 so với
1,25 --- phân tầng hai tầng rõ hơn trong mẫu con FDI dẫn dắt (Hsieh \&
Klenow, 2009).

\hypertarget{phuxe2n-tuxedch-suxe2u-ux111ux1ee3t-khux1ea3o-suxe1t-2025}{%
\paragraph{2.3.6.5 Phân tích sâu đợt khảo sát
2025}\label{phuxe2n-tuxedch-suxe2u-ux111ux1ee3t-khux1ea3o-suxe1t-2025}}

\textbf{Bảng 2.3.6.5}. \emph{12 nền kinh tế trong đợt khảo sát 2025
(n=14.712; xếp theo quy mô mẫu).}

\begin{longtable}[]{@{}
  >{\raggedright\arraybackslash}p{(\columnwidth - 16\tabcolsep) * \real{0.1111}}
  >{\raggedright\arraybackslash}p{(\columnwidth - 16\tabcolsep) * \real{0.1111}}
  >{\raggedright\arraybackslash}p{(\columnwidth - 16\tabcolsep) * \real{0.1111}}
  >{\raggedleft\arraybackslash}p{(\columnwidth - 16\tabcolsep) * \real{0.1111}}
  >{\raggedleft\arraybackslash}p{(\columnwidth - 16\tabcolsep) * \real{0.1111}}
  >{\raggedleft\arraybackslash}p{(\columnwidth - 16\tabcolsep) * \real{0.1111}}
  >{\raggedleft\arraybackslash}p{(\columnwidth - 16\tabcolsep) * \real{0.1111}}
  >{\raggedleft\arraybackslash}p{(\columnwidth - 16\tabcolsep) * \real{0.1111}}
  >{\raggedleft\arraybackslash}p{(\columnwidth - 16\tabcolsep) * \real{0.1111}}@{}}
\toprule\noalign{}
\begin{minipage}[b]{\linewidth}\raggedright
Nền kinh tế
\end{minipage} & \begin{minipage}[b]{\linewidth}\raggedright
ISO3
\end{minipage} & \begin{minipage}[b]{\linewidth}\raggedright
Nhóm ICRV
\end{minipage} & \begin{minipage}[b]{\linewidth}\raggedleft
n\_firms
\end{minipage} & \begin{minipage}[b]{\linewidth}\raggedleft
FSTS (\%)
\end{minipage} & \begin{minipage}[b]{\linewidth}\raggedleft
FDI (\%)
\end{minipage} & \begin{minipage}[b]{\linewidth}\raggedleft
R\&D (\%)
\end{minipage} & \begin{minipage}[b]{\linewidth}\raggedleft
Website (\%)
\end{minipage} & \begin{minipage}[b]{\linewidth}\raggedleft
sd ln(LP)¹
\end{minipage} \\
\midrule\noalign{}
\endhead
\bottomrule\noalign{}
\endlastfoot
Ấn Độ & IND & IV & 10.479 & 2,7 & 1,9 & 2,2 & 41,8 & 0,84 \\
Saudi Arabia & SAU & II & 1.002 & 2,7 & 9,5 & 1,7 & 30,2 &
\textbf{0,49} \\
Thái Lan & THA & III & 813 & 9,3 & 6,3 & 8,7 & 61,9 & 1,40 \\
Sri Lanka & LKA & V & 607 & 16,1 & 1,8 & 4,1 & 48,8 & 1,02 \\
Qatar & QAT & II & 480 & 2,3 & 19,4 & 0,6 & 63,3 & \textbf{0,33} \\
Afghanistan & AFG & V & 426 & 5,9 & 1,2 & 25,1 & 47,2 & 1,48 \\
Maldives & MDV & V & 154 & 5,8 & 4,5 & 22,4 & \textbf{73,4} & 1,38 \\
Fiji & FJI & VI & 151 & 12,5 & 9,9 & 18,8 & \textbf{74,8} & 1,12 \\
Brunei & BRN & II & 150 & 5,1 & 26,2 & 18,9 & \textbf{80,7} & 1,12 \\
Kuwait & KWT & II & 150 & \textbf{0,4} & 0,0 & 20,7 & 69,3 & 1,16 \\
Solomon Islands & SLB & VI & 150 & 4,8 & 20,0 & 11,3 & 53,3 & 1,29 \\
\textbf{Kiribati} & \textbf{KIR} & \textbf{VI} & \textbf{150} &
\textbf{1,0} & \textbf{0,7} & \textbf{8,0} & \textbf{18,7} &
\textbf{1,48} \\
\end{longtable}

\emph{¹ Độ lệch chuẩn ln(năng suất lao động) trong đợt --- bất biến tiền
tệ, so sánh được giữa các nền kinh tế. Ghi chú: đợt khảo sát Nepal 2025
dùng bảng hỏi BREADY Micro (doanh nghiệp siêu nhỏ, n=1.740;
\textasciitilde10\% bán ra thị trường quốc tế theo biến nhị phân
\texttt{me1c}, 22,6\% dùng máy tính/tablet) --- bộ công cụ không tương
đương với cross-section chuẩn nên không thuộc pool mô tả và chỉ nêu ở
đây như tham chiếu.}

Bảy phát hiện từ đợt 2025: (1) Ấn Độ FSTS sụt 5 điểm phần trăm so với
đợt 2022 (hiệu ứng schema); (2) Thái Lan ``số hóa tăng, xuất khẩu
giảm''; (3) Fiji website 74,8\% \textgreater{} Singapore 66,1\% (nhảy
vọt số); (4) Vùng Vịnh + Brunei xác nhận đặc trưng Advanced tài nguyên
dẫn dắt --- Saudi Arabia và Qatar có phân tán năng suất hẹp nhất toàn
pool (0,49 và 0,33); (5) tỷ lệ R\&D tăng đột biến ở một số nền
(Afghanistan 25,1\%, Kuwait 20,7\%) là dấu hiệu ước lượng quá mức do
bảng hỏi (schema-induced overestimation) cần kiểm chứng; (6) đợt 2025 là
tập xác thực ngoài mẫu tự nhiên cho Chuyên đề 2; (7) \textbf{Kiribati
cực đoan --- FSTS 1,0\%, FDI 0,7\%, website 18,7\% --- đối lập hoàn toàn
với Fiji. Dị biệt SIDS ``số hóa cao'' (FJI, MDV) so với ``cô lập nông
thôn'' (KIR) --- Chuyên đề 2 cần tách 2 phân nhóm con SIDS.}

\begin{center}\rule{0.5\linewidth}{0.5pt}\end{center}

\hypertarget{bux1ea3y-tiux1ec3u-cux1ea3nh-ux111iux1ec3n-huxecnh}{%
\subsubsection{2.3.7 Bảy tiểu cảnh điển
hình}\label{bux1ea3y-tiux1ec3u-cux1ea3nh-ux111iux1ec3n-huxecnh}}

\hypertarget{singapore-advanced-ux111ux1ed5i-mux1edbi-suxe1ng-tux1ea1o-dux1eabn-dux1eaft-n623-ux111ux1ee3t-2023}{%
\paragraph{2.3.7.1 Singapore (Advanced đổi mới sáng tạo dẫn dắt ---
n=623, đợt
2023)}\label{singapore-advanced-ux111ux1ed5i-mux1edbi-suxe1ng-tux1ea1o-dux1eabn-dux1eaft-n623-ux111ux1ee3t-2023}}

FSTS 7,1\%; doanh nghiệp xuất khẩu 17,8\%; website 66,1\%; ISO 23,3\%;
R\&D 7,5\%; FDI 31,5\%. Độ lệch chuẩn log năng suất trong đợt 1,08.
\textbf{Biên trên} của nhóm Advanced đổi mới sáng tạo dẫn dắt. Theo
nghiên cứu Singapore (Mar, Đỗ \& Phan, 2026 --- MIR under review): R²
hiệu chỉnh = 0,196; tương tác FSTS² × DAI = 3,119 (p=0,005); điểm uốn
\textasciitilde88,6\%.

\textbf{Bối cảnh hệ sinh thái 2025/2026}: Singapore xếp \textbf{hạng 2
toàn cầu} trong Innovators Business Environment Index 2026
(StartupBlink, 2026); hạng 4 toàn cầu trong Global Startup Ecosystem
Index 2025 với tăng trưởng hệ sinh thái 44,9\%; hạng 3 thế giới trong
Khảo sát Chính phủ điện tử Liên Hiệp Quốc 2024. Sáng kiến \textbf{Smart
Nation 2.0} (2024) với 3 trụ cột: Tin tưởng --- Tăng trưởng --- Cộng
đồng.

Hệ sinh thái Singapore vận hành ở Tier 3--4 --- vượt xa lớp DAI Tier
1--2 đo trong WBES. Đây là lý do \textbf{Nghịch lý bão hòa số (Digital
Saturation Paradox)} xuất hiện: DAI cơ bản đã trở thành yếu tố vệ sinh,
phần bù năng suất chỉ còn đo được ở nhóm xuất khẩu cường độ cao.

\hypertarget{saudi-arabia-qatar-vuxe0-kuwait-advanced-tuxe0i-nguyuxean-dux1eabn-dux1eaft-n1.632}{%
\paragraph{2.3.7.2 Saudi Arabia, Qatar và Kuwait (Advanced tài nguyên
dẫn dắt ---
n=1.632)}\label{saudi-arabia-qatar-vuxe0-kuwait-advanced-tuxe0i-nguyuxean-dux1eabn-dux1eaft-n1.632}}

\begin{longtable}[]{@{}lllll@{}}
\toprule\noalign{}
Chỉ số & Saudi Arabia & Qatar & Kuwait & Singapore \\
\midrule\noalign{}
\endhead
\bottomrule\noalign{}
\endlastfoot
FSTS (\%) & 2,7 & 2,3 & \textbf{0,4} & 7,1 \\
FDI ≥10\% (\%) & 9,5 & 19,4 & 0,0 & \textbf{31,5} \\
R\&D dương (\%) & 1,7 & 0,6 & \textbf{20,7} & 7,5 \\
sd log năng suất¹ & \textbf{0,49} & \textbf{0,33} & 1,16 & 1,08 \\
\end{longtable}

Năm phát hiện: (1) \textbf{nhà nước tô} --- phân phối tô lợi tài nguyên
san bằng phân tán năng suất (Beblawi, 1987; Hertog, 2010); (2)
\textbf{phân bổ sai nguồn lực đảo chiều} so với Hsieh \& Klenow (2009)
--- ở nhà nước tô, cơ chế phân bổ sai là bảo hộ FDI chứ không phải thiếu
hụt thể chế; (3) Kuwait Vision 2035 với R\&D 20,7\% --- cố gắng đa dạng
hóa đáng chú ý; (4) thiên lệch của chỉ số áp dụng số đơn thành phần khi
áp dụng ở bối cảnh tài nguyên; (5) \textbf{phân nhóm con Advanced} ---
bằng chứng cho sự đa dạng của chủ nghĩa tư bản (Hall \& Soskice, 2001).
\textbf{Bối cảnh 2026}: IMF điều chỉnh giảm tăng trưởng Saudi Arabia từ
\textasciitilde4,5\% xuống 3,1\% do xung đột Trung Đông.

\hypertarget{viux1ec7t-nam-3-ux111ux1ee3t-khux1ea3o-suxe1t-200920152023-pool-muxf4-tux1ea3-n3.077-mux1eabu-phuxe2n-tuxedch-cux1ee7a-nghiuxean-cux1ee9u-viux1ec7t-nam-n2.958-cux1eadp-nhux1eadt-adb-vietnam-outlook-2026}{%
\paragraph{2.3.7.3 Việt Nam (3 đợt khảo sát 2009/2015/2023; pool mô tả
n=3.077, mẫu phân tích của nghiên cứu Việt Nam n=2.958) --- cập nhật ADB
Vietnam Outlook
2026}\label{viux1ec7t-nam-3-ux111ux1ee3t-khux1ea3o-suxe1t-200920152023-pool-muxf4-tux1ea3-n3.077-mux1eabu-phuxe2n-tuxedch-cux1ee7a-nghiuxean-cux1ee9u-viux1ec7t-nam-n2.958-cux1eadp-nhux1eadt-adb-vietnam-outlook-2026}}

FSTS 23,2\% → 17,9\% → 16,1\% (suy giảm theo thời gian); doanh nghiệp
xuất khẩu 37,1\% → 23,8\%; ISO 17--23\%; R\&D 6,1\% (2023). Pattern
\textbf{kinh tế hai tầng} --- FDI hướng xuất khẩu hiệu quả cao đan xen
với doanh nghiệp nội địa năng suất thấp (CIEM, 2023).

\textbf{Bối cảnh 2026 --- ADB Vietnam Economic Outlook}: Việt Nam dự
phóng GDP \textbf{7,2\% (2026) / 7,0\% (2027)} --- vượt ASEAN 4,6\%. Bốn
động lực: (a) cam kết FDI 2,4 tỷ USD với 18 dự án chiến lược; (b) năng
suất lao động +5,1\% trong giai đoạn 2026--2027; (c) khoản vay ADB dựa
trên chính sách 2,4 tỷ USD nhắm mục tiêu SME nội địa; (d) tái định vị
trong chuỗi giá trị toàn cầu từ lắp ráp downstream sang thiết kế và R\&D
mid-stream.

\textbf{World Bank FCI Snapshot (2025)}: GDP tăng trưởng 7,09\% (2024),
FDI 20,17 tỷ USD (4,23\% GDP), xuất khẩu máy móc và điện tử 45,8\% tổng
xuất khẩu, tỷ lệ nữ/nam lao động 88,61\% (2023).

\hypertarget{trung-quux1ed1c-mux1eabu-phuxe2n-tuxedch-cux1ee7a-nghiuxean-cux1ee9u-trung-quux1ed1c-n4.559-pool-muxf4-tux1ea3-gux1ed3m-ux111ux1ee3t-2024-n2.189}{%
\paragraph{2.3.7.4 Trung Quốc (mẫu phân tích của nghiên cứu Trung Quốc
n=4.559; pool mô tả gồm đợt 2024,
n=2.189)}\label{trung-quux1ed1c-mux1eabu-phuxe2n-tuxedch-cux1ee7a-nghiuxean-cux1ee9u-trung-quux1ed1c-n4.559-pool-muxf4-tux1ea3-gux1ed3m-ux111ux1ee3t-2024-n2.189}}

FSTS 10,9\% → 8,8\%; FDI ≥10\% chiếm 6,0\%. Quan hệ FSTS -- năng suất có
dạng hàm bậc hai (U ngược) với điểm uốn \textasciitilde47,5\% (Đỗ \&
Phan, 2026 --- JFAR). \textbf{Rủi ro tập trung}: Wang, Huang và Hong
(2024 --- IRFA) phân tích 80\% mô hình rủi ro ngân hàng tại Trung Quốc
phụ thuộc các nhà cung cấp công nghệ chi phối thị trường --- cảnh báo
cho lộ trình sản xuất thông minh.

\hypertarget{tux1ed5ng-hux1ee3p-khux1ed1i-chuxe2u-uxe1-mux1edbi-nux1ed5i-7-nux1ec1n-n43.502-trux1b0ux1eddng-hux1ee3p-ux1ea5n-ux111ux1ed9}{%
\paragraph{2.3.7.5 Tổng hợp khối châu Á mới nổi (7 nền, n=43.502) ---
trường hợp Ấn
Độ}\label{tux1ed5ng-hux1ee3p-khux1ed1i-chuxe2u-uxe1-mux1edbi-nux1ed5i-7-nux1ec1n-n43.502-trux1b0ux1eddng-hux1ee3p-ux1ea5n-ux111ux1ed9}}

FSTS 8,2\%; FDI ≥10\% chiếm 3,5\%; tỷ lệ doanh nghiệp R\&D dương 16,5\%.
Phân tán năng suất trong đợt ở mức cao (trung vị 1,31) và dao động mạnh
giữa ba phân nhóm con (1,25--1,47) --- phản ánh dị biệt rộng xuyên 7 nền
(Bảng 2.3.6.4).

\textbf{Trường hợp Ấn Độ --- lan tỏa công nghệ}: Theo Sikdar \&
Mukhopadhyay (2026 --- \emph{Asian Development Review, 43}(1), 37--75),
panel 4.293 doanh nghiệp Ấn Độ 2006--2019: FDI là kênh lan tỏa ngang
quan trọng nhất; doanh nghiệp lớn thu nhận lan tỏa nhiều hơn; FDI tăng
từ 1,705 triệu USD (1990--2000) → 42,396 triệu USD (2012--2022); R\&D
trì trệ ở 0,7\% GDP.

\hypertarget{muxf4ng-cux1ed5-tux1eeb-lux1eddi-nguyux1ec1n-tuxe0i-nguyuxean-ux111ux1ebfn-khouxe1ng-sux1ea3n-thiux1ebft-yux1ebfu-n1.082-3-ux111ux1ee3t-200920132019}{%
\paragraph{2.3.7.6 Mông Cổ --- từ lời nguyền tài nguyên đến khoáng sản
thiết yếu (n=1.082, 3 đợt
2009/2013/2019)}\label{muxf4ng-cux1ed5-tux1eeb-lux1eddi-nguyux1ec1n-tuxe0i-nguyuxean-ux111ux1ebfn-khouxe1ng-sux1ea3n-thiux1ebft-yux1ebfu-n1.082-3-ux111ux1ee3t-200920132019}}

FSTS đình trệ 4--6\% (5,8\% → 3,9\% → 4,2\%); FDI ≥10\% giảm
\textbf{7,2\% → 4,2\%}; website dao động 39--49\%. Pattern \textbf{lời
nguyền tài nguyên} (Sachs \& Warner, 2001). \textbf{Bước ngoặt 2026}: 11
khoáng sản thiết yếu (Cu, Li, Mo, Mn, Co, Ni, W, REE). Xuất khẩu đồng
1,69 triệu tấn (2024) = 22,1\% tổng xuất khẩu. Oyu Tolgoi mục tiêu
500.000 tấn/năm 2028--2036. Trung Quốc chi phối chế biến: 100\%
graphite, 90\% Mn, 70\% Co (ADB, 2026, tháng 5).

\emph{Hình 2.3.7.1. Mông Cổ: tiến triển 2009--2019 (3 đợt khảo sát WBES)
--- FSTS đình trệ; FDI giảm dần phản ánh lời nguyền tài nguyên giai đoạn
1.}

\hypertarget{sids-thuxe1i-buxecnh-dux1b0ux1a1ng-7-nux1ec1n-n1.781}{%
\paragraph{2.3.7.7 SIDS Thái Bình Dương (7 nền ---
n=1.781)}\label{sids-thuxe1i-buxecnh-dux1b0ux1a1ng-7-nux1ec1n-n1.781}}

\textbf{7 nền}: Fiji (FJI), Papua New Guinea (PNG), Solomon Islands
(SLB), Tonga (TON), Vanuatu (VUT), Samoa (WSM), \textbf{Kiribati (KIR
--- 2025)}. FSTS 6,2\%; FDI ≥10\% \textbf{22,9\%}; đổi mới sản phẩm
\textbf{39,6\%}; website \textbf{47,1\%}. Pattern \textbf{thích nghi
trong điều kiện ràng buộc}.

\textbf{Kiribati 2025 --- trường hợp biên cực đoan nhất}: FSTS 1,0\%,
FDI 0,7\%, website 18,7\%, ISO 1,4\%, độ lệch chuẩn log trong đợt 1,48,
SME 93,3\%.

\textbf{Bối cảnh kinh tế vĩ mô PICs Pacific}: 4 trụ cột MIRAB (Bertram,
2006): (1) khu vực công chủ đạo (30--50\% việc làm chính thức); (2) phụ
thuộc 4 nguồn ngoại sinh --- ODA 30\% GDP và remittances 25--40\% GDP;
(3) thanh niên thất nghiệp 15--30\%; (4) lao động di chuyển --- PLMS Úc
và RSE New Zealand (Vanuatu RSE 15--20\% GDP). Chuyên đề 2 cần bổ sung 3
biến kiểm soát đặc trưng PICs: \texttt{aid\_share\_GDP\_pct},
\texttt{remittance\_share\_GDP\_pct},
\texttt{seasonal\_worker\_outflow\_pct}.

\hypertarget{so-suxe1nh-tux1ed5ng-hux1ee3p}{%
\paragraph{2.3.7.8 So sánh tổng
hợp}\label{so-suxe1nh-tux1ed5ng-hux1ee3p}}

\textbf{Bảng 2.3.7.1}. \emph{Ma trận so sánh đa chiều bảy tiểu cảnh điển
hình.}

\begin{longtable}[]{@{}
  >{\raggedright\arraybackslash}p{(\columnwidth - 14\tabcolsep) * \real{0.1250}}
  >{\raggedright\arraybackslash}p{(\columnwidth - 14\tabcolsep) * \real{0.1250}}
  >{\raggedright\arraybackslash}p{(\columnwidth - 14\tabcolsep) * \real{0.1250}}
  >{\raggedright\arraybackslash}p{(\columnwidth - 14\tabcolsep) * \real{0.1250}}
  >{\raggedright\arraybackslash}p{(\columnwidth - 14\tabcolsep) * \real{0.1250}}
  >{\raggedright\arraybackslash}p{(\columnwidth - 14\tabcolsep) * \real{0.1250}}
  >{\raggedright\arraybackslash}p{(\columnwidth - 14\tabcolsep) * \real{0.1250}}
  >{\raggedright\arraybackslash}p{(\columnwidth - 14\tabcolsep) * \real{0.1250}}@{}}
\toprule\noalign{}
\begin{minipage}[b]{\linewidth}\raggedright
Chỉ số
\end{minipage} & \begin{minipage}[b]{\linewidth}\raggedright
Singapore
\end{minipage} & \begin{minipage}[b]{\linewidth}\raggedright
Việt Nam
\end{minipage} & \begin{minipage}[b]{\linewidth}\raggedright
Trung Quốc
\end{minipage} & \begin{minipage}[b]{\linewidth}\raggedright
Khối mới nổi (7 nền)
\end{minipage} & \begin{minipage}[b]{\linewidth}\raggedright
Mông Cổ
\end{minipage} & \begin{minipage}[b]{\linewidth}\raggedright
SIDS Thái Bình Dương
\end{minipage} & \begin{minipage}[b]{\linewidth}\raggedright
Saudi+Qatar+Kuwait
\end{minipage} \\
\midrule\noalign{}
\endhead
\bottomrule\noalign{}
\endlastfoot
n\_firms & 623 & 3.077 & 2.189 & 43.502 & 1.082 & \textbf{1.781} &
1.632 \\
FSTS (\%) & 7,1 & \textbf{19,1} & 8,8 & 8,2 & 4,7 & 6,2 &
\textbf{2,4} \\
FDI ≥10\% (\%) & \textbf{31,5} & 11,4 & 6,0 & 3,5 & 5,7 & \textbf{22,9}
& 11,5 \\
R\&D dương (\%) & 7,5 & 6,1 & \textbf{39,4} & 16,5 & \textbf{33,7} &
11,5 & \textbf{3,1} \\
Website (\%) & \textbf{66,1} & 47,0 & 42,0 & 49,4 & 43,6 & 47,1 &
43,6 \\
sd ln(LP) trong đợt¹ & \textbf{1,08} & 1,46 & 1,36 & 1,31 & 1,27 & 1,29
& \textbf{0,49} \\
\end{longtable}

\emph{¹ Trung vị độ lệch chuẩn ln(năng suất lao động) trong từng đợt
khảo sát (bất biến tiền tệ). Hàng SIDS gồm 7 nền Pacific (bao gồm
Kiribati 2025); cột Saudi+Qatar+Kuwait lấy trung vị ba đợt 2025.}

\begin{center}\rule{0.5\linewidth}{0.5pt}\end{center}

\hypertarget{yux1ebfu-tux1ed1-giux1ea3i-thuxedch-sux1a1-bux1ed9}{%
\subsubsection{2.3.8 Yếu tố giải thích sơ
bộ}\label{yux1ebfu-tux1ed1-giux1ea3i-thuxedch-sux1a1-bux1ed9}}

\hypertarget{hux1ec7-sux1ed1-tux1b0ux1a1ng-quan-pearson}{%
\paragraph{2.3.8.1 Hệ số tương quan
Pearson}\label{hux1ec7-sux1ed1-tux1b0ux1a1ng-quan-pearson}}

Tương quan giữa năng suất và bốn yếu tố giải thích được tính trên năng
suất chuẩn hóa z-score trong từng cặp nền kinh tế × năm (bất biến tiền
tệ), gộp theo 6 nhóm ICRV của pool mô tả.

\textbf{Bảng 2.3.8.1}. \emph{Hệ số tương quan Pearson giữa năng suất lao
động chuẩn hóa (LP-z) và bốn yếu tố giải thích, theo nhóm ICRV (pool mô
tả 49 nền).}

\begin{longtable}[]{@{}
  >{\raggedright\arraybackslash}p{(\columnwidth - 12\tabcolsep) * \real{0.1429}}
  >{\raggedleft\arraybackslash}p{(\columnwidth - 12\tabcolsep) * \real{0.1429}}
  >{\raggedleft\arraybackslash}p{(\columnwidth - 12\tabcolsep) * \real{0.1429}}
  >{\raggedleft\arraybackslash}p{(\columnwidth - 12\tabcolsep) * \real{0.1429}}
  >{\raggedleft\arraybackslash}p{(\columnwidth - 12\tabcolsep) * \real{0.1429}}
  >{\raggedleft\arraybackslash}p{(\columnwidth - 12\tabcolsep) * \real{0.1429}}
  >{\raggedleft\arraybackslash}p{(\columnwidth - 12\tabcolsep) * \real{0.1429}}@{}}
\toprule\noalign{}
\begin{minipage}[b]{\linewidth}\raggedright
Yếu tố
\end{minipage} & \begin{minipage}[b]{\linewidth}\raggedleft
I --- Adv. đổi mới
\end{minipage} & \begin{minipage}[b]{\linewidth}\raggedleft
II --- Adv. tài nguyên
\end{minipage} & \begin{minipage}[b]{\linewidth}\raggedleft
III --- Trung bình cao
\end{minipage} & \begin{minipage}[b]{\linewidth}\raggedleft
IV --- Chuyển đổi TB-thấp
\end{minipage} & \begin{minipage}[b]{\linewidth}\raggedleft
V --- Đang nổi
\end{minipage} & \begin{minipage}[b]{\linewidth}\raggedleft
VI --- SIDS
\end{minipage} \\
\midrule\noalign{}
\endhead
\bottomrule\noalign{}
\endlastfoot
FDI ≥ 10\% & +0,130** & +0,114** & +0,124** & +0,072** & +0,057** &
+0,122** \\
FSTS & \textbf{+0,163**} & +0,131** & +0,094** & +0,047** & +0,018* &
\textbf{−0,009 ns} \\
TCI (R\&D + ISO) & +0,176** & +0,187** & \textbf{+0,240**} & +0,141** &
+0,125** & +0,117** \\
DAI (website) & +0,098** & \textbf{+0,201**} & +0,124** & +0,180** &
+0,111** & +0,092** \\
\end{longtable}

\emph{Ghi chú: ** p \textless{} 0,01; * p \textless{} 0,05; ns = không
có ý nghĩa thống kê (kiểm định hai phía). LP-z = z-score của ln(doanh
thu/lao động) trong từng cặp nền kinh tế × năm; FDI ≥10\% từ b2b; FSTS =
d3b + d3c; TCI = trung bình hai thành phần R\&D (h8) và ISO (b8); DAI =
website (c22b). n theo nhóm: 3.620 / 2.082 / 10.138 / 40.508 / 15.669 /
1.715 (quan sát có LP hợp lệ). Hệ số tương quan Pearson --- không áp
dụng HC1; HC1 dành cho SE hồi quy ở Chuyên đề 2.}

Sáu phát hiện chính: (1) \textbf{TCI là yếu tố tương quan mạnh và đồng
đều nhất} ở mọi nhóm (+0,117 đến +0,240) --- năng lực công nghệ nền tảng
là dịch chuyển mức phổ quát; (2) \textbf{gradient cường độ FSTS theo
chiều thể chế}: tương quan xuất khẩu--năng suất giảm đơn điệu từ +0,163
(Nhóm I) qua +0,047 (IV) và +0,018 (V) xuống \textbf{mất ý nghĩa ở SIDS
(−0,009 ns)} --- tương quan thô của quốc tế hóa gần như biến mất ở thể
chế yếu, gợi ý quan hệ phi tuyến và bị điều tiết thay vì tuyến tính phổ
quát; (3) \textbf{SIDS là ngoại lệ duy nhất về dấu của FSTS} --- phù hợp
với giả thuyết chi phí buộc phải quốc tế hóa: doanh nghiệp SIDS xuất
khẩu vì bắt buộc, không phải vì chọn lọc theo năng suất; (4) \textbf{FDI
dương ở mọi nhóm}, mạnh hơn ở hai cực thể chế (Nhóm I +0,130 và SIDS
+0,122) --- hai cơ chế khác nhau: MNE hub tại Advanced và FDI ngoại sinh
tại SIDS; (5) \textbf{DAI dương ở mọi nhóm nhưng thấp nhất tại Nhóm I
(+0,098)} trong khi cao nhất ở II và IV (+0,201; +0,180) --- phần bù của
hiện diện số Tier-1 co lại nơi nó đã tiệm cận bão hòa (Nghịch lý bão hòa
số ở dạng suy giảm lợi suất biên); (6) \textbf{biến thiên cường độ
cross-regime có hệ thống} --- thể chế điều tiết độ lớn (và cá biệt là
dấu) của các cơ chế, nhất quán với Xu (2024).

\textbf{Cảnh báo schema BREADY 2025}: Bảng 2.3.8.1 có thể bị nhiễu hiệu
ứng schema --- cần kiểm chứng qua hệ thống phòng thủ ba tầng (§2.2)
trước khi suy diễn nhân quả.

\hypertarget{pattern-giux1edbi-tuxednh-trong-quux1ea3n-luxfd-vuxe0-sux1edf-hux1eefu-eap-dux1eabn-ux111ux1ea7u-touxe0n-cux1ea7u}{%
\paragraph{2.3.8.2 Pattern giới tính trong quản lý và sở hữu --- EAP dẫn
đầu toàn
cầu}\label{pattern-giux1edbi-tuxednh-trong-quux1ea3n-luxfd-vuxe0-sux1edf-hux1eefu-eap-dux1eabn-ux111ux1ea7u-touxe0n-cux1ea7u}}

\textbf{Bảng 2.3.8.2}. \emph{Pattern giới tính trong lãnh đạo cấp cao và
sở hữu doanh nghiệp theo khu vực.}

\begin{longtable}[]{@{}llll@{}}
\toprule\noalign{}
Khu vực & \% nữ TMT & \% nữ sở hữu đa số & Khoảng cách \\
\midrule\noalign{}
\endhead
\bottomrule\noalign{}
\endlastfoot
\textbf{East Asia \& Pacific (EAP)} & \textbf{33,4\%} & \textbf{25,7\%}
& +7,7 đpt \\
Latin America \& Caribbean & 27,8\% & 18,3\% & +9,5 đpt \\
Sub-Saharan Africa & 21,3\% & 14,9\% & +6,4 đpt \\
Europe \& Central Asia & 19,8\% & 12,1\% & +7,7 đpt \\
South Asia & 9,2\% & 8,5\% & +0,7 đpt \\
\textbf{MENA} & \textbf{3,3\%} & \textbf{1,7\%} & +1,6 đpt \\
\end{longtable}

Bốn quan sát: (a) EAP paradox --- quản lý dẫn đầu toàn cầu nhưng
ownership ceiling thấp hơn management ceiling 7,7 điểm; (b) MENA
bottleneck --- rào cản thể chế kép trong bối cảnh nhà nước tô (Beblawi,
1987); (c) 43/50 nền kinh tế có tỷ lệ nữ trong nhóm quản trị cấp cao cao
hơn nữ sở hữu đa số; (d) hàm ý cho Chuyên đề 2: tương tác giới tính ×
FSTS dương tại Emerging qua kênh quan hệ (Hambrick \& Mason, 1984) nhưng
không có ý nghĩa thống kê tại Advanced.

\hypertarget{ux111ux1b0ux1eddng-cong-u-cux1ee7a-ux111ux1ed5i-mux1edbi-sux1ed1-tux1eeb-chi-phuxed-ux111ux1ea7u-tux1b0-ux111ux1ebfn-phux1ea7n-buxf9-nux103ng-suux1ea5t}{%
\paragraph{2.3.8.3 Đường cong U của đổi mới số --- từ chi phí đầu tư đến
phần bù năng
suất}\label{ux111ux1b0ux1eddng-cong-u-cux1ee7a-ux111ux1ed5i-mux1edbi-sux1ed1-tux1eeb-chi-phuxed-ux111ux1ea7u-tux1b0-ux111ux1ebfn-phux1ea7n-buxf9-nux103ng-suux1ea5t}}

Quan hệ giữa cường độ đầu tư số và năng suất có hình dạng phi tuyến ---
\textbf{đường cong U theo thời gian} --- phù hợp với trễ công nghệ của
David (1990) và nghịch lý năng suất của Brynjolfsson \& McAfee (2014).

\textbf{Giai đoạn 1} (năm 1--3): Sụt giảm ngắn hạn do chi phí thiết bị,
đào tạo và gián đoạn quy trình. Bằng chứng gián tiếp: phần bù tương quan
DAI thấp nhất tại Nhóm Advanced đổi mới (+0,098 so với +0,180--0,201 ở
các nhóm chuyển đổi/tài nguyên) --- nơi Tier-1 đã tiệm cận bão hòa, lợi
suất biên co lại; giai đoạn ``sân khấu số'' (Verhoef et al., 2021).

\textbf{Giai đoạn 2} (năm 3--5): Vùng chuyển tiếp. Ngưỡng: ≥ 3\% doanh
thu/năm × ≥ 3 năm liên tục. Dưới ngưỡng, phần bù năng suất không bù được
chi phí cơ hội.

\textbf{Giai đoạn 3} (năm 5+): Phần bù năng suất. Bằng chứng: Trung Quốc
R\&D 39,4\% (Tier-2 mature); FSTS²×DAI = +3,119 (p = 0,005) tại
Singapore --- DAI khuếch đại quan hệ quốc tế hóa -- hiệu quả tại doanh
nghiệp FSTS cao.

\textbf{Liên kết CDCM}: Giai đoạn 1 = proxy cơ chế lỗi thời (DAI Tier-1
bão hòa); Giai đoạn 3 = cơ chế phụ thuộc giai đoạn (DAI Tier-2 có phần
bù có điều kiện). Hai cơ chế bổ sung nhau khi kiểm soát thế hệ số.

\begin{center}\rule{0.5\linewidth}{0.5pt}\end{center}

\hypertarget{thux1ea3o-luux1eadn-tux1ed5ng-hux1ee3p}{%
\subsubsection{2.3.9 Thảo luận tổng
hợp}\label{thux1ea3o-luux1eadn-tux1ed5ng-hux1ee3p}}

\hypertarget{mux1b0ux1eddi-kux1ebft-luux1eadn-tux1ed5ng-hux1ee3p-tux1eeb-dux1eef-liux1ec7u-wbes}{%
\paragraph{2.3.9.1 Mười kết luận tổng hợp từ dữ liệu
WBES}\label{mux1b0ux1eddi-kux1ebft-luux1eadn-tux1ed5ng-hux1ee3p-tux1eeb-dux1eef-liux1ec7u-wbes}}

\textbf{(i) Phân tán năng suất trong đợt tăng theo chiều thể chế yếu
dần}: sd ln(LP) trong đợt từ 1,04 (Nhóm I) lên 1,31--1,39 (Nhóm III--V);
tỷ số P90/P10 leo từ \textasciitilde12 lần (Nhóm I) lên
\textasciitilde30 lần (Nhóm V). Pattern này \textbf{nhất quán với giả
thuyết phân bổ sai nguồn lực} (Hsieh \& Klenow, 2009, 2014) --- lưu ý sd
log năng suất lao động là proxy gián tiếp cho cơ chế phân tán TFPR, nên
đây là bằng chứng mô tả phù hợp chứ không phải phân rã trực tiếp.

\textbf{(ii) Dị biệt nội bộ phân nhóm Advanced}: hai nhóm cùng mức thu
nhập cao nhưng cấu trúc đối lập --- R\&D 20,5\% so với 7,0\%; doanh
nghiệp xuất khẩu 28,4\% so với 11,7\%; nữ quản lý cấp cao 27,5\% so với
4,6\%; phân tán năng suất trong đợt đặc biệt hẹp ở hai nền lớn nhất Vùng
Vịnh (Saudi Arabia 0,49; Qatar 0,33 so với Singapore 1,08 --- chênh hơn
\textbf{2 lần}, đặc trưng nhà nước tô) --- cần phân nhóm con Advanced
trong Chuyên đề 2.

\textbf{(iii) Dị biệt nội bộ khối châu Á mới nổi --- 3 phân nhóm con}:
FDI dẫn dắt SEA (VNM+IDN+PHL, FSTS 12,4\%) ≠ dân số lớn (IND+LKA+JOR,
FSTS 7,1\%) ≠ tài nguyên (MNG, FSTS 4,7\%).

\textbf{(iv) SIDS Pacific --- Kiribati 2025 là trường hợp biên cực đoan
nhất}: FSTS 1,0\%, FDI 0,7\%, website 18,7\%, ISO 1,4\% --- đối lập với
Fiji nhảy vọt số (74,8\%). Cơ sở cho giả thuyết chi phí buộc phải quốc
tế hóa.

\textbf{(v) Quốc tế hóa là hiện tượng phân cực}: trung vị FSTS = 0\%;
chỉ khoảng 12--28\% doanh nghiệp xuất khẩu tùy nhóm. Cần lựa chọn hai
giai đoạn trong Chuyên đề 2.

\textbf{(vi) Nhảy vọt số 2018--2026} ở các nhóm IV/V/VI (+21 đến +31
điểm phần trăm website). Phân biệt \textbf{Tier-1 Sự hiện diện số cơ
bản} (tiệm cận bão hòa ở Advanced) và \textbf{Tier-2 Chuyển đổi số cơ
bản} (đang phát triển). Phần bù tương quan DAI thấp nhất ở Nhóm I
(+0,098) là dấu hiệu bão hòa Tier-1 --- xem kiểm định loại trừ ICT tại
§2.3.4.3.

\textbf{(vii) Pattern phi tuyến FDI ≥10\%} --- dạng chữ U với cực tiểu ở
Nhóm IV (3,1\%): Advanced 10,6--13,6\% → Nhóm IV 3,1\% → SIDS 20,0\%.
Hai mô hình FDI: MNE hub (Advanced) và du lịch/viễn thông (SIDS).

\textbf{(viii) Đợt khảo sát 2025 --- mẫu đơn năm lớn nhất}: 12 nền kinh
tế × 14.712 doanh nghiệp (17,0\% pool). Ấn Độ FSTS sụt 5 điểm phần trăm
(hiệu ứng schema); Fiji website 74,8\% \textgreater{} Singapore 66,1\%;
Kiribati đối lập cực đoan với Fiji.

\textbf{(ix) Khung phân tích cấp ngành}: 9 ngành ISIC Rev.~4; kiểm định
loại trừ ICT; bổ sung Dynamism profile theo Kafouros et al.~(2023).

\textbf{(x) Hệ thống tái tạo được cho 4 thế hệ schema WBES}: Tất cả bảng
và hình đều tái tạo được. Gói tái lập mở cho cộng đồng nghiên cứu.

\hypertarget{nux103m-khoux1ea3ng-trux1ed1ng-nghiuxean-cux1ee9u}{%
\paragraph{2.3.9.2 Năm khoảng trống nghiên
cứu}\label{nux103m-khoux1ea3ng-trux1ed1ng-nghiuxean-cux1ee9u}}

\textbf{KT1 --- Điều tiết ba chiều TCI×DAI×FSTS chưa được kiểm định đồng
thời trên mẫu liên quốc gia}: Các nghiên cứu hiện tại kiểm định từng cặp
điều tiết riêng biệt. Mô hình M7 của Chuyên đề 2 sẽ lấp khoảng trống này
với 102 cặp nền kinh tế--năm.

\textbf{KT2 --- Phân nhóm con Advanced: innovation-driven so với
resource-driven chưa được hệ thống hóa}: Nghiên cứu hiện hành (Bausch \&
Krist, 2007; Kirca et al., 2012) gộp tất cả OECD vào ``phát triển''.
Bằng chứng Singapore so với Saudi Arabia/Qatar/Kuwait cho thấy hai phân
nhóm có pattern hoàn toàn khác nhau.

\textbf{KT3 --- Xác nhận nhân quả Nghịch lý bão hòa số qua ước lượng
biến công cụ}: Phần bù tương quan DAI suy giảm rõ ở Advanced (+0,098 ---
thấp nhất sáu nhóm) có thể phản ánh: (a) bão hòa thực sự của Tier-1; (b)
nhân quả ngược; (c) bias biến bỏ sót. Biến công cụ hợp lệ (độ phủ băng
thông rộng) cần thiết để phân biệt.

\textbf{KT4 --- Dị biệt theo thời gian ba giai đoạn 2006--2026 chưa được
tổng hợp}: Bằng chứng từ kiểm định Paternoster cho thấy bất ổn định tham
số quan trọng giữa ba giai đoạn.

\textbf{KT5 --- SIDS Pacific như trường hợp biên lý thuyết của thuyết
Uppsala}: Chi phí buộc phải quốc tế hóa tại 7 SIDS không khớp với bất kỳ
dạng hàm quốc tế hóa -- hiệu quả nào trong văn liệu. Cần khung kết hợp
MIRAB và Briguglio (1995) và quan điểm dựa trên nguồn lực.

\textbf{Ánh xạ khoảng trống → câu hỏi nghiên cứu cho Chuyên đề 2.} Để
chính thức hóa cầu nối từ thực trạng mô tả (CĐ1) sang thiết kế kiểm định
suy diễn (CĐ2), Bảng 2.3.9.1 chuyển mỗi khoảng trống thực nghiệm thành
một câu hỏi nghiên cứu có mô hình/kiểm định tương ứng.

\textbf{Bảng 2.3.9.1.} \emph{Ánh xạ khoảng trống thực trạng (CĐ1) → câu
hỏi nghiên cứu (CĐ2).}

\begin{longtable}[]{@{}
  >{\raggedright\arraybackslash}p{(\columnwidth - 4\tabcolsep) * \real{0.3333}}
  >{\raggedright\arraybackslash}p{(\columnwidth - 4\tabcolsep) * \real{0.3333}}
  >{\raggedright\arraybackslash}p{(\columnwidth - 4\tabcolsep) * \real{0.3333}}@{}}
\toprule\noalign{}
\begin{minipage}[b]{\linewidth}\raggedright
Khoảng trống (CĐ1)
\end{minipage} & \begin{minipage}[b]{\linewidth}\raggedright
Câu hỏi nghiên cứu cho CĐ2
\end{minipage} & \begin{minipage}[b]{\linewidth}\raggedright
Mô hình/kiểm định dự kiến
\end{minipage} \\
\midrule\noalign{}
\endhead
\bottomrule\noalign{}
\endlastfoot
\textbf{KT1} --- Điều tiết ba chiều TCI×DAI×FSTS chưa kiểm định đồng
thời & Năng lực công nghệ và chấp nhận số đồng thời điều tiết quan hệ
I--P như thế nào xuyên 6 chế độ ICRV? & Mô hình M7 tương tác ba chiều,
102 cặp nền kinh tế--năm \\
\textbf{KT2} --- Phân nhóm con Advanced chưa hệ thống hóa & Cơ chế
innovation-driven và resource-driven khác nhau ra sao trong chuyển hóa
quốc tế hóa thành hiệu quả? & Tách Nhóm I/II + kiểm định khác biệt hệ
số \\
\textbf{KT3} --- Nghịch lý bão hòa số chưa xác nhận nhân quả & Phần bù
DAI suy giảm ở Advanced là bão hòa thực hay nhân quả ngược/bỏ sót biến?
& Ước lượng biến công cụ (độ phủ băng thông rộng) \\
\textbf{KT4} --- Dị biệt thời gian ba giai đoạn chưa tổng hợp & Vị trí
điểm uốn và độ cong I--P có ổn định qua 2006--2026 không? & Kiểm định
Paternoster + tương tác thời gian \\
\textbf{KT5} --- SIDS là trường hợp biên của Uppsala & Có tồn tại bối
cảnh nơi quan hệ I--P chuyển thành âm đơn điệu (FIP) không? & Hồi quy
SIDS + khung MIRAB/Briguglio (1995) \\
\end{longtable}

\emph{Ghi chú: Bảng này bảo đảm mỗi khoảng trống thực nghiệm rút ra từ
bức tranh thực trạng được chuyển thành một câu hỏi nghiên cứu có thiết
kế kiểm định cụ thể trong Chuyên đề 2, duy trì tính liên tục logic giữa
hai chuyên đề.}

\hypertarget{huxe0m-uxfd-luxfd-thuyux1ebft}{%
\paragraph{2.3.9.3 Hàm ý lý
thuyết}\label{huxe0m-uxfd-luxfd-thuyux1ebft}}

\textbf{(1) Uppsala có điều kiện --- không phổ quát}: Tương quan
FSTS--năng suất dương mạnh tại Advanced (+0,163), suy yếu dần theo chiều
thể chế và mất ý nghĩa tại SIDS (−0,009 ns) --- xác nhận điều kiện biên:
lý thuyết Uppsala vận hành chỉ trong môi trường thể chế đủ mạnh
(Williamson, 1985).

\textbf{(2) CDCM như khung lý thuyết tích hợp mới cho kỷ nguyên số trong
kinh doanh quốc tế}: Phần bù năng suất của áp dụng số là hàm của (i) mức
độ phát triển thể chế, (ii) vị trí trên đường cong quốc tế hóa -- hiệu
quả, (iii) tầng của năng lực số được đo lường. Ba điều kiện tương tác
phi tuyến --- vượt qua ``phần bù số'' đồng nhất (Verhoef et al., 2021).

\textbf{(3) Giao diện Đổi mới-Hiệu suất Nâng cao (AIPI)}: FSTS²×DAI =
+3,119 tại Singapore gợi ý cơ chế mới --- tại FSTS cao, doanh nghiệp cần
DAI như ``phối hợp kỹ thuật số xuyên biên giới'' để quản lý độ phức tạp
chuỗi giá trị toàn cầu. Cần kiểm định tại Advanced khác (Hàn Quốc, Đài
Loan).

\textbf{(4) Sân khấu số là thực tế đo lường ở cả hai chiều}: Tại
Advanced, DAI Tier-1 là yếu tố vệ sinh bão hòa. Tại Emerging, DAI Tier-1
có thể là proxy lỗi thời --- đo lường không còn phản ánh được cơ chế
kinh tế thực.

\textbf{(5) Bối cảnh thể chế là tầng điều kiện, không phải yếu tố nền}:
Biến thiên cường độ có hệ thống (và cá biệt là đổi dấu --- FSTS tại
SIDS) của các hệ số cross-regime không phải nhiễu mà là tín hiệu lý
thuyết. Nghiên cứu kinh doanh quốc tế tương lai nên kiểm định hiệu ứng
điều tiết theo nhóm ICRV trước khi báo cáo kích thước hiệu ứng trung
bình.

\hypertarget{huxe0m-uxfd-phux1b0ux1a1ng-phuxe1p-luux1eadn}{%
\paragraph{2.3.9.4 Hàm ý phương pháp
luận}\label{huxe0m-uxfd-phux1b0ux1a1ng-phuxe1p-luux1eadn}}

\textbf{(1) Pool WBES xuyên quốc gia như nền tảng dữ liệu cho kinh doanh
quốc tế thị trường mới nổi}: Khung 49 nền với 91.982 doanh nghiệp ở tầng
ước lượng (pool phân loại 96.415 DN / 52 nền) và 86.701 doanh nghiệp /
102 cặp nền kinh tế × năm ở tầng mô tả là một trong những mẫu lớn nhất
từ trước đến nay cho phân tích quốc tế hóa -- hiệu quả toàn diện tại
châu Á--Thái Bình Dương.

\textbf{(2) Hệ thống phòng thủ ba tầng cho panel đa thế hệ là đóng góp
phương pháp luận}: Ba tầng kiểm chứng --- biến giả Post\_BREADY\_2024,
kiểm định ổn định mô hình neo, panel hậu đại dịch độc lập --- là giao
thức có thể chuẩn hóa cho bất kỳ nghiên cứu dùng WBES với nhiều thế hệ
schema.

\textbf{(3) Winsorize trong cụm quốc gia × năm tốt hơn winsorize tổng
thể}: Winsorize tổng thể bóp méo phân phối tại các nhóm phân tán rộng
(Nhóm V: P90/P10 \textasciitilde30 lần trong đợt). Winsorize trong cụm
bảo tồn dị biệt cross-regime quan trọng về lý thuyết.

\textbf{(4) HC1 robust SE như chuẩn mực cho hồi quy đa quốc gia WBES}:
HC3 cần thiết cho mẫu con nhỏ (SIDS n \textless{} 200). HC1 là cân bằng
tốt giữa hiệu suất ước lượng và độ bền vững cho mẫu đầy đủ.

\textbf{(5) Chiến lược mô hình neo bảo vệ suy luận khi schema chuyển
đổi}: 6 biến bất biến cross-schema (ln\_employees, FSTS, ISO,
innovate\_product, website, ln\_LP) cho phép so sánh hệ số cấu trúc giữa
PICS3 và BREADY.

\begin{center}\rule{0.5\linewidth}{0.5pt}\end{center}

\hypertarget{kux1ebft-luux1eadn}{%
\subsection{2.4 KẾT LUẬN}\label{kux1ebft-luux1eadn}}

\hypertarget{kux1ebft-luux1eadn-chuxednh}{%
\subsubsection{2.4.1 Kết luận chính}\label{kux1ebft-luux1eadn-chuxednh}}

Tám kết luận từ 49 nền kinh tế châu Á -- Thái Bình Dương (pool mô tả
86.701 doanh nghiệp / 102 cặp nền kinh tế × năm, WBES 2006--2026; khung
ước lượng của luận án 91.982 doanh nghiệp):

\textbf{KL1 --- Phi tuyến xuyên chế độ thể chế là quy luật}: U ngược tại
các nhóm chuyển đổi (Việt Nam điểm uốn 39--46\%; Trung Quốc điểm uốn
47--49\%); gần tuyến tính tại Advanced (Singapore điểm uốn 88,6\%); chi
phí buộc phải tại SIDS.

\textbf{KL2 --- TCI là dịch chuyển mức phổ quát, DAI là khuếch đại phụ
thuộc bối cảnh}: TCI tương quan dương mạnh và đồng đều ở mọi nhóm; phần
bù DAI biến thiên cường độ gấp đôi giữa các nhóm ICRV (+0,098 ở Advanced
so với +0,201 ở Advanced tài nguyên) --- đây là nội hàm cốt lõi của
CDCM.

\textbf{KL3 --- Nghịch lý bão hòa số là cơ chế đo lường, không phải thất
bại thị trường}: Phần bù DAI thấp nhất tại Advanced là hệ quả của đo
lường Tier-1 trong bối cảnh Tier-2/3 đã được áp dụng rộng.

\textbf{KL4 --- Thể chế là biến điều tiết, không phải biến trung gian}:
Gap de jure--de facto (Xu, 2024) giải thích biến thiên cường độ
cross-regime của FDI và FSTS. Chính sách nâng cao thể chế hiệu quả nhất
khi song hành với nâng cấp TCI và DAI.

\textbf{KL5 --- Dị biệt theo thời gian ba giai đoạn là thực tế cấu
trúc}: Hiệu ứng DAI mạnh hơn trong BREADY (sau tăng tốc số hóa hậu
COVID-19). Hệ thống phòng thủ ba tầng đảm bảo kết luận không bị nhiễu
schema.

\textbf{KL6 --- EAP dẫn đầu toàn cầu về nữ quản lý cấp cao (33,4\%)
nhưng khoảng cách sở hữu dai dẳng}: Trần kính ở lớp sở hữu chưa được phá
vỡ dù management ceiling đang thu hẹp.

\textbf{KL7 --- SIDS Pacific là trường hợp biên lý thuyết với chi phí
buộc phải quốc tế hóa}: FDI (+0,122) là nguồn năng suất chủ đạo; tương
quan FSTS--năng suất bằng không (−0,009, không có ý nghĩa thống kê) ---
nhóm duy nhất nơi xuất khẩu không gắn với phần bù năng suất. Mô hình
MIRAB phù hợp hơn mô hình Uppsala.

\textbf{KL8 --- Phân nhóm con Advanced là cần thiết cho phân tích kinh
doanh quốc tế châu Á}: Singapore và Saudi Arabia/Qatar/Kuwait cùng nhãn
``Advanced'' nhưng pattern năng suất doanh nghiệp hoàn toàn khác nhau.

\hypertarget{huxe0m-uxfd-vuxe0-ux111ux1ec1-xuux1ea5t}{%
\subsubsection{2.4.2 Hàm ý và đề
xuất}\label{huxe0m-uxfd-vuxe0-ux111ux1ec1-xuux1ea5t}}

\hypertarget{huxe0m-uxfd-chuxednh-suxe1ch-cho-doanh-nghiux1ec7p-viux1ec7t-nam-vuxe0-khu-vux1ef1c-chuxe2u-uxe1}{%
\paragraph{2.4.2.1 Hàm ý chính sách cho doanh nghiệp Việt Nam và khu vực
châu
Á}\label{huxe0m-uxfd-chuxednh-suxe1ch-cho-doanh-nghiux1ec7p-viux1ec7t-nam-vuxe0-khu-vux1ef1c-chuxe2u-uxe1}}

\textbf{(1) Sequencing năng lực trước --- nâng cấp TCI trước khi mở rộng
FSTS}: Hệ số TCI = 0,179 (nhân quả ước lượng biến công cụ, nghiên cứu
Việt Nam) xác nhận: xuất khẩu chỉ tạo năng suất khi doanh nghiệp có năng
lực công nghệ nền tảng. Chương trình hỗ trợ xuất khẩu cần song hành với
nâng cấp TCI.

\textbf{(2) Ưu tiên Tier-2 số cho SME Việt Nam}: Website 47\% là điểm
đầu vào; VietQR, TikTok Shop, Shopee Asia Pacific là Tier-1.5/Tier-2
thực tế hơn. Voucher AI Adoption (\textasciitilde50 triệu VND/doanh
nghiệp) giúp vượt qua giai đoạn đầu tư ban đầu.

\textbf{(3) Phân biệt FDI chiến lược với FDI lắp ráp}: FSTS suy giảm
23,2\% → 16,1\% phản ánh FDI lắp ráp tăng nhưng giá trị gia tăng nội địa
thấp. Các chỉ tiêu cần theo dõi: tỷ lệ R\&D/doanh thu, tỷ lệ nhà cung
cấp nội địa Tier-1, bằng sáng chế đồng nộp.

\textbf{(4) Tái định vị chuỗi giá trị toàn cầu: lắp ráp downstream →
thiết kế và R\&D mid-stream}: Việt Nam ở điểm uốn (39--46\%). Trung Quốc
R\&D 39,4\% và điểm uốn cao hơn (47--49\%) cho thấy: Sự chênh lệch điểm
uốn giữa Trung Quốc (TP \textasciitilde47--49\%, Nhóm III) và Việt Nam
(TP \textasciitilde39--46\%, Nhóm IV) phản ánh gradient ICRV, không phải
gradient TCI. TCI nâng mặt bằng năng suất ở mọi mức FSTS (hiệu ứng
level-shift) nhưng không trực tiếp quyết định vị trí điểm uốn --- phân
biệt hai cơ chế này là quan trọng cho hàm ý chính sách.

\textbf{(5) Nền tảng học hỏi chính sách dựa trên ICRV cho ASEAN}: Trao
đổi thực hành tốt nhất giữa nền kinh tế cùng nhóm IV Emerging (Việt Nam,
Indonesia, Philippines, Ấn Độ) thay vì học từ Singapore/Hàn Quốc (nhóm I
--- bối cảnh khác biệt quá lớn).

\textbf{(6) Chính sách số hóa phân cấp theo vùng và quy mô}: Hiệu ứng
DAI phụ thuộc bối cảnh --- đầu tư website khác nhau tại TP.HCM
(Tier-1.5) so với nông thôn đồng bằng sông Cửu Long (Tier-1 chưa bão
hòa).

\textbf{(7) Hợp tác lao động di chuyển Pacific cho Việt Nam --- bài học
PLMS/RSE}: (a) \textbf{PLMS Úc}: Học mô hình 30.000 lao động PICs/năm →
triển khai Vietnamese Skilled Worker Mobility Scheme (nông nghiệp công
nghệ cao và chăm sóc người cao tuổi); (b) \textbf{RSE New Zealand}: Nuôi
trồng thủy sản, cà phê, đào tạo kỹ năng điều hàng --- tận dụng chuyên
môn đặc thù Việt Nam; (c) \textbf{Kênh remittance kỹ thuật số}: Hợp tác
Wise/MoMo giảm phí từ \textasciitilde7\% xuống \textless3\%; (d)
\textbf{Nâng cấp kỹ năng theo mô hình EPS Hàn Quốc}: Đào tạo kỹ năng
trước khi xuất khẩu lao động và kênh tái di cư.

\textbf{(8) Bốn lộ trình số hóa châu Á tham khảo cho SME Việt Nam}: (a)
Singapore/Hàn Quốc AI full-stack (Tier-3) --- không phù hợp để sao chép
trực tiếp; (b) Nhật Bản 60\% doanh nghiệp AI back-office (Tier-2) ---
\textbf{mục tiêu thực tế 2026--2030}; (c) Trung Quốc sản xuất thông minh
(Tier-2/3) --- học điểm tích cực, tránh rủi ro tập trung; (d) \textbf{Mô
hình ưu tiên Việt Nam Tier-1.5}: VietQR (100\% miễn phí) → Thương mại xã
hội → Xuyên biên giới Shopee/Amazon → Cloud và AI cơ bản.

\hypertarget{ux111ux1ecbnh-hux1b0ux1edbng-cho-chuyuxean-ux111ux1ec1-2}{%
\paragraph{2.4.2.2 Định hướng cho Chuyên đề
2}\label{ux111ux1ecbnh-hux1b0ux1edbng-cho-chuyuxean-ux111ux1ec1-2}}

Tám kiểm định độ bền vững bắt buộc cho Chuyên đề 2:

\textbf{(1) Thay biến phụ thuộc}: ln(LP) → ROS, tăng trưởng doanh thu,
tăng trưởng việc làm. Kết quả chính cần bền vững ≥ 2/3 biến phụ thuộc.

\textbf{(2) Lọc mẫu}: Loại micro firms (\textless{} 5 lao động) và doanh
nghiệp nhà nước. Kết quả cần bền vững trong SME tư nhân (n ≈ 80.000+).

\textbf{(3) Thay phương pháp}: OLS → hồi quy phân vị (p25/p50/p75);
bootstrap SE 500 lần lặp; HC3.

\textbf{(4) Kiểm định dạng hàm}: LOWESS phi tham số; kiểm định
Lind--Mehlum (2010) bắt buộc xác nhận U ngược; AIC/BIC so sánh tuyến
tính so với bậc hai so với bậc ba.

\textbf{(5) Placebo tests}: Thay FSTS + TCI bằng biến nhiễu. Kết quả
chính không được có ý nghĩa thống kê.

\textbf{(6) Nhạy cảm cutoffs ICRV}: ±0,1 đơn vị WGI; PPP GDP/capita thay
GNI; jackknife loại từng nền kinh tế (49 lần).

\textbf{(7) Độ bền vững theo thời gian}: Chow test ổn định tham số;
Paternoster z-test so sánh điểm uốn giữa ba giai đoạn.

\textbf{(8) Kiểm định ranh giới schema}: Hệ thống phòng thủ ba tầng tại
§2.2 --- điều kiện tiên quyết trước khi báo cáo BREADY 2025 như bằng
chứng nhân quả.

\begin{center}\rule{0.5\linewidth}{0.5pt}\end{center}

\hypertarget{tuxe0i-liux1ec7u-tham-khux1ea3o}{%
\subsection{TÀI LIỆU THAM KHẢO}\label{tuxe0i-liux1ec7u-tham-khux1ea3o}}

\emph{Trích dẫn theo APA 7th edition.}

ADB. (2024). \emph{Asian development outlook 2024: Strengthening climate
resilience}. Asian Development Bank.

ADB. (2026, tháng 4). \emph{Asian development outlook April 2026: The
Middle East conflict challenges resilience in Asia and the Pacific}.
Asian Development Bank.

ADB. (2026). \emph{Vietnam economic outlook 2026: Navigating the
crosscurrents}. Asian Development Bank.

Banalieva, E. R., \& Dhanaraj, C. (2019). Internalization theory for the
digital economy. \emph{Journal of International Business Studies,
50}(8), 1372--1387.

Bausch, A., \& Krist, M. (2007). The effect of context-related
moderators on the internationalization--performance relationship:
Evidence from meta-analysis. \emph{Management International Review,
47}(3), 319--347.

Beblawi, H. (1987). The rentier state in the Arab world. In H. Beblawi
\& G. Luciani (Eds.), \emph{The rentier state} (pp.~49--62). Croom Helm.

Bertram, G. (2006). Introduction: The MIRAB model in the twenty-first
century. \emph{Asia Pacific Viewpoint, 47}(1), 1--13.

Briguglio, L. (1995). Small island developing states and their economic
vulnerabilities. \emph{World Development, 23}(9), 1615--1632.

Brynjolfsson, E., \& McAfee, A. (2014). \emph{The second machine age:
Work, progress, and prosperity in a time of brilliant technologies}. W.
W. Norton \& Company.

Cohen, W. M., \& Levinthal, D. A. (1990). Absorptive capacity: A new
perspective on learning and innovation. \emph{Administrative Science
Quarterly, 35}(1), 128--152.

Combs, J. G., Crook, T. R., \& Shook, C. L. (2005). The dimensionality
of organizational performance and its implications for strategic
management research. In D. J. Ketchen \& D. D. Bergh (Eds.),
\emph{Research methodology in strategy and management} (Vol. 2,
pp.~259--286). Emerald. https://doi.org/10.1016/S1479-8387(05)02011-4

Contractor, F. J., Kundu, S. K., \& Hsu, C. C. (2003). A three-stage
theory of international expansion. \emph{Journal of International
Business Studies, 34}(1), 5--18.

Cusolito, A. P., \& Maloney, W. F. (2018). \emph{Productivity revisited:
Shifting paradigms in analysis and policy}. World Bank Group.

David, P. A. (1990). The dynamo and the computer: An historical
perspective on the modern productivity paradox. \emph{American Economic
Review, 80}(2), 355--361.

Đỗ, T. H., \& Phan, A. T. (2026). Internationalization and firm
performance in Vietnam. \emph{Journal of Asia Business Studies} (under
review).

Hall, P. A., \& Soskice, D. (Eds.). (2001). \emph{Varieties of
capitalism: The institutional foundations of comparative advantage}.
Oxford University Press.

Hessels, J., \& Terjesen, S. (2010). Resource dependency and
institutional theory perspectives on direct and indirect export choices.
\emph{Small Business Economics, 34}(2), 203--220.
https://doi.org/10.1007/s11187-008-9156-4

Hertog, S. (2010). \emph{Princes, brokers, and bureaucrats: Oil and the
state in Saudi Arabia}. Cornell University Press.

Hsieh, C.-T., \& Klenow, P. J. (2009). Misallocation and manufacturing
TFP in China and India. \emph{Quarterly Journal of Economics, 124}(4),
1403--1448.

Hsieh, C.-T., \& Klenow, P. J. (2014). The life cycle of plants in India
and Mexico. \emph{Quarterly Journal of Economics, 129}(3), 1035--1084.

Johanson, J., \& Vahlne, J.-E. (1977). The internationalization process
of the firm. \emph{Journal of International Business Studies, 8}(1),
23--32.

Johanson, J., \& Vahlne, J.-E. (2009). The Uppsala internationalization
process model revisited. \emph{Journal of International Business
Studies, 40}(9), 1411--1431.

Kafouros, M., Wang, C., Piperopoulos, P., \& Zhang, M. (2023). Academic
publishing in business and management. \emph{Global Strategy Journal,
13}(1), 3--28.

Kaufmann, D., Kraay, A., \& Mastruzzi, M. (2011). The worldwide
governance indicators: Methodology and analytical issues. \emph{Hague
Journal on the Rule of Law, 3}(2), 220--246.

Khanna, T., \& Palepu, K. G. (2010). \emph{Winning in emerging markets:
A road map for strategy and execution}. Harvard Business Press.

Kirca, A. H., Roth, K., Hult, G. T. M., \& Cavusgil, S. T. (2012). The
role of context in the multinationality-performance relationship: A
meta-analytic review. \emph{Global Strategy Journal, 2}(2), 108--121.

Lu, J. W., \& Beamish, P. W. (2001). The internationalization and
performance of SMEs. \emph{Strategic Management Journal, 22}(6--7),
565--586. https://doi.org/10.1002/smj.184

Lu, J. W., \& Beamish, P. W. (2004). International diversification and
firm performance: The S-curve hypothesis. \emph{Academy of Management
Journal, 47}(4), 598--609.

Mar, J., Đỗ, T. H., \& Phan, A. T. (2026). Digital adoption and
internationalization in Singapore. \emph{Management International
Review} (under review).

Marano, V., Arregle, J.-L., Hitt, M. A., Spadafora, E., \& van Essen, M.
(2016). Home country institutions and the
internationalization-performance relationship. \emph{Journal of
Management, 42}(5), 1075--1110.

North, D. C. (1990). \emph{Institutions, institutional change and
economic performance}. Cambridge University Press.

Paternoster, R., Brame, R., Mazerolle, P., \& Piquero, A. (1998). Using
the correct statistical test for the equality of regression
coefficients. \emph{Criminology, 36}(4), 859--866.

Peng, M. W., \& Ilinitch, A. Y. (1998). Export intermediary firms: A
note on export development research. \emph{Journal of International
Business Studies, 29}(3), 609--620.
https://doi.org/10.1057/palgrave.jibs.8490011

Peng, M. W. (2001). The resource-based view and international business.
\emph{Journal of Management, 27}(6), 803--829.

Peng, M. W. (2003). Institutional transitions and strategic choices.
\emph{Academy of Management Review, 28}(2), 275--296.

Richard, P. J., Devinney, T. M., Yip, G. S., \& Johnson, G. (2009).
Measuring organizational performance: Towards methodological best
practice. \emph{Journal of Management, 35}(3), 718--804.
https://doi.org/10.1177/0149206308330560

Sachs, J. D., \& Warner, A. M. (2001). The curse of natural resources.
\emph{European Economic Review, 45}(4--6), 827--838.

Schwens, C., Zapkau, F. B., Brouthers, K. D., \& Hollender, L. (2018).
How overconfidence can blind internationalization performance
predictions. \emph{Journal of World Business, 53}(3), 290--301.

Sikdar, A., \& Mukhopadhyay, A. (2026). Technology spillovers and firm
performance in India. \emph{Asian Development Review, 43}(1), 37--75.

StartupBlink. (2026). \emph{Global startup ecosystem index 2026}.
StartupBlink Research.

Stallkamp, M., \& Schotter, A. P. J. (2021). Platforms without borders?
\emph{Global Strategy Journal, 11}(1), 58--80.

Sullivan, D. (1994). Measuring the degree of internationalization of a
firm. \emph{Journal of International Business Studies, 25}(2), 325--342.
https://doi.org/10.1057/palgrave.jibs.8490203

Teece, D. J., Pisano, G., \& Shuen, A. (1997). Dynamic capabilities and
strategic management. \emph{Strategic Management Journal, 18}(7),
509--533.

UNCTAD. (2023). \emph{World investment report 2023}. United Nations
Conference on Trade and Development.

Venkatraman, N., \& Ramanujam, V. (1986). Measurement of business
performance in strategy research. \emph{Academy of Management Review,
11}(4), 801--814.

Verhoef, P. C., Broekhuizen, T., Bart, Y., Bhattacharya, A., Dong, J.
Q., Fabian, N., \& Haenlein, M. (2021). Digital transformation: A
multidisciplinary reflection and research agenda. \emph{Journal of
Business Research, 122}, 889--901.

Wang, J., Huang, Y., \& Hong, H. (2024). Technology concentration and
bank risk: Evidence from Chinese commercial banks. \emph{International
Review of Financial Analysis, 91}, 102968.

Wernerfelt, B. (1984). A resource-based view of the firm.
\emph{Strategic Management Journal, 5}(2), 171--180.

Williamson, O. E. (1985). \emph{The economic institutions of
capitalism}. Free Press.

World Bank. (2023). \emph{World Bank Enterprise Surveys 2023}. World
Bank Group.

World Bank. (2024). \emph{World development indicators 2024}. World Bank
Group.

World Bank. (2025). \emph{World Bank Enterprise Surveys 2025}. World
Bank Group.

Wu, J., Xu, L., \& Yang, Z. (2022). Multinationality and performance of
emerging economy multinational enterprises: A twenty-year study.
\emph{International Business Review, 31}(5), 102003.

Xu, D. (2024). De jure versus de facto institutional distance.
\emph{Global Strategy Journal, 14}(2), 187--215.

\begin{center}\rule{0.5\linewidth}{0.5pt}\end{center}

\hypertarget{phux1ee5-lux1ee5c}{%
\subsection{PHỤ LỤC}\label{phux1ee5-lux1ee5c}}

\hypertarget{phux1ee5-lux1ee5c-a-phux1ea1m-vi-thux1ef1c-tux1ebf-cux1ee7a-pool-dux1eef-liux1ec7u-wbes}{%
\subsubsection{Phụ lục A --- Phạm vi thực tế của pool dữ liệu
WBES}\label{phux1ee5-lux1ee5c-a-phux1ea1m-vi-thux1ef1c-tux1ebf-cux1ee7a-pool-dux1eef-liux1ec7u-wbes}}

\textbf{Bảng A.1}. \emph{Phân bổ 49 nền kinh tế theo nhóm ICRV --- pool
mô tả 86.701 doanh nghiệp, 102 cặp nền kinh tế × năm, 2006--2026.}

\begin{longtable}[]{@{}
  >{\raggedright\arraybackslash}p{(\columnwidth - 6\tabcolsep) * \real{0.2500}}
  >{\raggedright\arraybackslash}p{(\columnwidth - 6\tabcolsep) * \real{0.2500}}
  >{\raggedleft\arraybackslash}p{(\columnwidth - 6\tabcolsep) * \real{0.2500}}
  >{\raggedleft\arraybackslash}p{(\columnwidth - 6\tabcolsep) * \real{0.2500}}@{}}
\toprule\noalign{}
\begin{minipage}[b]{\linewidth}\raggedright
Nhóm ICRV
\end{minipage} & \begin{minipage}[b]{\linewidth}\raggedright
Nền kinh tế (năm khảo sát)
\end{minipage} & \begin{minipage}[b]{\linewidth}\raggedleft
Đợt
\end{minipage} & \begin{minipage}[b]{\linewidth}\raggedleft
n
\end{minipage} \\
\midrule\noalign{}
\endhead
\bottomrule\noalign{}
\endlastfoot
\textbf{Nhóm I --- Advanced đổi mới sáng tạo (5 nền)} & Hàn Quốc (2024),
Israel (2013, 2024), Singapore (2023), Đài Loan (2024), Hong Kong SAR
(2023) & 6 & 4.222 \\
\textbf{Nhóm II --- Advanced tài nguyên (6 nền)} & Saudi Arabia (2025),
Qatar (2025), Oman (2026), Bahrain (2024), Brunei (2025), Kuwait (2025)
& 6 & 2.231 \\
\textbf{Nhóm III --- Trung bình cao (6 nền)} & Kazakhstan (2009, 2013,
2019, 2024), Malaysia (2015, 2019, 2024), Trung Quốc (2024), Thái Lan
(2016, 2025), Armenia (2009, 2013, 2020, 2024), Georgia (2013, 2019,
2023) & 17 & 13.993 \\
\textbf{Nhóm IV --- Chuyển đổi TB-thấp (7 nền)} & Ấn Độ (2014, 2022,
2025), Indonesia (2009, 2015, 2023), Việt Nam (2009, 2015, 2023),
Pakistan (2013, 2022), Bangladesh (2013, 2022), Mông Cổ (2009, 2013,
2019), Philippines (2023) & 17 & 45.003 \\
\textbf{Nhóm V --- Đang nổi (17 nền)} & Uzbekistan, Jordan, Lào, Iraq,
Myanmar, Kyrgyzstan, Sri Lanka, Lebanon, Tajikistan, Campuchia,
Afghanistan, Yemen, Bhutan, Nepal, Azerbaijan, Turkmenistan, Maldives
(các đợt 2006--2025) & 40 & 18.957 \\
\textbf{Nhóm VI --- SIDS (8 nền)} & Timor-Leste (2009, 2015, 2021), Fiji
(2009, 2025), Solomon Islands (2015, 2025), Tonga (2009, 2024), Samoa
(2009, 2023), Vanuatu (2009, 2023), Papua New Guinea (2015, 2024),
Kiribati (2025) & 16 & 2.295 \\
\textbf{Tổng} & \textbf{49 nền kinh tế} & \textbf{102} &
\textbf{86.701} \\
\end{longtable}

\emph{Ghi chú: mỗi cặp nền kinh tế × năm lấy đúng một cross-section
chuẩn WBES; loại các điều tra phi chính thức, siêu nhỏ (Micro) và panel
mở rộng; chỉ nhận đợt khảo sát từ 2006 trở đi (bộ câu hỏi so sánh được
toàn cầu). Khung ước lượng của luận án (91.982 doanh nghiệp / 49 nền;
pool phân loại 96.415 / 52 nền) được hợp nhất và hài hòa hóa riêng theo
quy trình tại Phụ lục A của luận án; hai khung thống nhất về danh sách
nền kinh tế và nhãn nhóm ICRV.}

\textbf{Ba thế hệ schema WBES trong pool mô tả}:

\begin{longtable}[]{@{}
  >{\raggedright\arraybackslash}p{(\columnwidth - 8\tabcolsep) * \real{0.2000}}
  >{\raggedright\arraybackslash}p{(\columnwidth - 8\tabcolsep) * \real{0.2000}}
  >{\raggedleft\arraybackslash}p{(\columnwidth - 8\tabcolsep) * \real{0.2000}}
  >{\raggedleft\arraybackslash}p{(\columnwidth - 8\tabcolsep) * \real{0.2000}}
  >{\raggedright\arraybackslash}p{(\columnwidth - 8\tabcolsep) * \real{0.2000}}@{}}
\toprule\noalign{}
\begin{minipage}[b]{\linewidth}\raggedright
Schema
\end{minipage} & \begin{minipage}[b]{\linewidth}\raggedright
Năm
\end{minipage} & \begin{minipage}[b]{\linewidth}\raggedleft
Đợt
\end{minipage} & \begin{minipage}[b]{\linewidth}\raggedleft
n
\end{minipage} & \begin{minipage}[b]{\linewidth}\raggedright
Phạm vi DAI
\end{minipage} \\
\midrule\noalign{}
\endhead
\bottomrule\noalign{}
\endlastfoot
PICS3 / tiền-Standardized & 2006--2012 & 18 & 8.048 & Chỉ Tier-1
(website nhị phân) \\
Standardized & 2013--2017 & 27 & 23.940 & Chỉ Tier-1 \\
BREADY/BEE & 2018--2026 & 57 & 54.713 & Tier-1 + Tier-2 (thanh toán điện
tử k33/k38) \\
\end{longtable}

Hệ thống phòng thủ ba tầng (§2.2) đảm bảo kết quả từ BREADY 2025 không
bị nhiễu hiệu ứng schema.

\begin{center}\rule{0.5\linewidth}{0.5pt}\end{center}

\emph{Chuyên đề tiến sĩ số 1 --- Phiên bản nộp (12/05/2026)}
\emph{Nghiên cứu sinh: Đỗ Thùy Hương --- Trường Đại học Kỹ thuật Công
nghệ Vĩnh Long (VLUTE)} \emph{Người hướng dẫn: TS. Nguyễn Minh Cảnh ---
Trường Kinh tế, Đại học Cần Thơ (CTU)}

\end{document}
